% Modified on 1/13/05.
\newdimen\snellbaselineskip
\newdimen\snellskip
\snellskip=1.5ex
\snellbaselineskip=\baselineskip
\def\srule{\omit\kern.5em\vrule\kern-.5em}
\newbox\bigstrutbox
\setbox\bigstrutbox=\hbox{\vrule height14.5pt depth9.5pt width0pt}
\def\bigstrut{\relax\ifmmode\copy\bigstrutbox\else\unhcopy\bigstrutbox\fi}
\def\middlehrule#1#2{\noalign{\kern-\snellbaselineskip\kern\snellskip}
&\multispan#1\strut\hrulefill
&\omit\hbox to.5em{\hrulefill}\vrule 
height \snellskip\kern-.5em&\multispan#2\hrulefill\cr}

\makeatletter
\def\bordermatrix#1{\begingroup \m@th
  \@tempdima 8.75\p@
  \setbox\z@\vbox{%
    \def\cr{\crcr\noalign{\kern2\p@\global\let\cr\endline}}%
    \ialign{$##$\hfil\kern2\p@\kern\@tempdima&\thinspace\hfil$##$\hfil
      &&\quad\hfil$##$\hfil\crcr
      \omit\strut\hfil\crcr\noalign{\kern-\snellbaselineskip}%
      #1\crcr\omit\strut\cr}}%
  \setbox\tw@\vbox{\unvcopy\z@\global\setbox\@ne\lastbox}%
  \setbox\tw@\hbox{\unhbox\@ne\unskip\global\setbox\@ne\lastbox}%
  \setbox\tw@\hbox{$\kern\wd\@ne\kern-\@tempdima\left(\kern-\wd\@ne
    \global\setbox\@ne\vbox{\box\@ne\kern2\p@}%
    \vcenter{\kern-\ht\@ne\unvbox\z@\kern-\snellbaselineskip}\,\right)$}%
  \null\;\vbox{\kern\ht\@ne\box\tw@}\endgroup}

\makeatother



%\setcounter{chapter}{9}

\chapter{Rantai Markov}\label{chp 11}

\section{Pendahuluan}\label{sec 11.1}

\par
Sebagian besar kajian kita tentang peluang telah membahas proses percobaan
independen. Proses-proses ini merupakan dasar teori peluang klasik dan banyak
bagian statistika. Kita telah membahas dua teorema utama bagi proses-proses
tersebut: Hukum Bilangan Besar dan Teorema Limit Pusat.
\par
Kita telah melihat bahwa ketika suatu barisan percobaan acak membentuk proses
percobaan independen, hasil-hasil yang mungkin dari setiap percobaan adalah sama
dan terjadi dengan peluang yang sama. Selain itu, pengetahuan tentang hasil
percobaan sebelumnya tidak memengaruhi prediksi kita mengenai hasil percobaan
berikutnya. Distribusi hasil satu percobaan cukup untuk membangun pohon dan
ukuran pohon bagi barisan $n$ percobaan, dan kita dapat menjawab sebarang
pertanyaan peluang tentang percobaan-percobaan ini menggunakan ukuran pohon tersebut.
\par
Teori peluang modern mempelajari proses acak ketika pengetahuan tentang hasil
sebelumnya memengaruhi prediksi bagi percobaan mendatang. Pada prinsipnya, ketika
kita mengamati barisan percobaan acak, semua hasil masa lalu dapat memengaruhi
prediksi kita bagi percobaan berikutnya. Sebagai contoh, hal ini seharusnya
berlaku ketika memprediksi nilai seorang siswa dalam serangkaian ujian suatu
mata kuliah. Namun, membolehkan keumuman sebesar ini akan sangat menyulitkan
pembuktian hasil-hasil umum.
\par
Pada 1907, A.~A. Markov mulai mempelajari jenis baru proses acak yang penting.
Dalam proses ini, hasil suatu percobaan dapat memengaruhi hasil percobaan
berikutnya. Jenis proses ini disebut rantai Markov.
\subsection*{Menentukan Rantai Markov}\index{rantai Markov}
Rantai Markov dideskripsikan sebagai berikut. Kita memiliki himpunan \emx
{keadaan,} $S = \{s_1,s_2,\ldots,s_r\}$.\index{keadaan!dari rantai Markov}
Proses dimulai pada salah satu keadaan ini dan berpindah secara berturut-turut
dari satu keadaan ke keadaan lain. Setiap perpindahan disebut \emx {langkah.}
Jika rantai saat ini berada dalam keadaan $s_i$, pada langkah berikutnya rantai
berpindah ke keadaan $s_j$ dengan peluang yang dinyatakan oleh $p_{ij}$, dan
peluang ini tidak bergantung pada keadaan mana saja yang ditempati rantai
sebelum keadaan saat ini.
\par
Peluang~$p_{ij}$ disebut \emx {peluang transisi.}\index{peluang transisi}
\index{peluang!transisi} Proses dapat tetap berada dalam keadaannya saat ini,
dan hal ini terjadi dengan peluang~$p_{ii}$. Distribusi peluang awal yang
didefinisikan pada~$S$ menentukan keadaan awal. Biasanya hal ini dilakukan
dengan menetapkan suatu keadaan tertentu sebagai keadaan awal.

R.~A. Howard\footnote{R.~A. Howard, \emx {Dynamic Probabilistic Systems,}
vol.~1
(New York: John Wiley and Sons, 1971).}\index{HOWARD, R. A.} memberikan
gambaran menarik tentang rantai Markov sebagai seekor katak yang melompat di
atas sekumpulan daun teratai. Katak tersebut mulai pada salah satu daun, lalu
melompat dari satu daun teratai ke daun lain dengan peluang transisi yang sesuai.



\begin{example}\label{exam 11.1.1}
Menurut Kemeny, Snell, dan Thompson,\footnote{J.~G. Kemeny, J.~L. Snell,
G.~L. Thompson, \emx {Introduction to Finite Mathematics,} 3rd ed.\ (Englewood
Cliffs, NJ: Prentice-Hall, 1974).}\index{KEMENY, J. G.}\index{SNELL, J.
L.}\index{THOMPSON, G. L.} Negeri Oz\index{Oz, Negeri} dianugerahi banyak hal,
tetapi bukan cuaca yang baik. Di sana tidak pernah terjadi dua hari cerah
berturut-turut. Jika suatu hari cerah, keesokan harinya salju dan hujan sama
mungkin terjadi. Jika turun salju atau hujan, peluang cuaca yang sama terjadi
keesokan harinya adalah setengah. Jika cuaca berubah dari salju atau hujan,
hanya separuh perubahan itu yang menghasilkan hari cerah. Dengan informasi ini
kita membentuk rantai Markov sebagai berikut. Jenis cuaca R, N, dan S kita ambil
sebagai keadaan. Dari informasi di atas kita menentukan peluang transisinya.
Peluang-peluang tersebut paling mudah dinyatakan dalam larik persegi sebagai
$$
\mat {P} = \bordermatrix{
        & \mbox {R} & \mbox {N} & \mbox {S} \cr
\mbox {R} &     1/2 &     1/4 &     1/4 \cr
\mbox {N} &     1/2 &       0 &     1/2 \cr
\mbox {S} &     1/4 &     1/4 &     1/2}\ .
$$
\end{example}

\subsection*{Matriks Transisi}
Entri-entri pada baris pertama matriks $\mat {P}$ dalam Contoh~\ref{exam
11.1.1} menyatakan peluang berbagai jenis cuaca setelah hari hujan. Demikian
pula, entri-entri pada baris kedua dan ketiga masing-masing menyatakan peluang
berbagai jenis cuaca setelah hari cerah dan hari bersalju. Larik persegi semacam
ini disebut \emx {matriks peluang transisi}, atau \emx {matriks
transisi}.\index{matriks transisi}
\par
Kita meninjau masalah menentukan peluang bahwa, jika rantai berada dalam
keadaan $i$ hari ini, rantai akan berada dalam keadaan $j$ dua hari lagi.
Peluang ini kita nyatakan dengan $p_{ij}^{(2)}$. Dalam
Contoh~\ref{exam 11.1.1}, kita melihat bahwa jika hari ini hujan, kejadian bahwa
dua hari lagi turun salju merupakan gabungan saling lepas dari tiga kejadian
berikut: 1) besok hujan dan dua hari lagi turun salju, 2) besok cerah dan dua
hari lagi turun salju, dan 3) besok turun salju dan dua hari lagi juga turun
salju. Peluang kejadian pertama adalah hasil kali peluang bersyarat bahwa besok
hujan jika hari ini hujan, dan peluang bersyarat bahwa dua hari lagi turun
salju jika besok hujan. Dengan menggunakan matriks transisi $\mat{P}$, hasil
kali ini dapat ditulis sebagai $p_{11}p_{13}$. Kedua kejadian lain juga
memiliki peluang yang dapat ditulis sebagai hasil kali entri-entri $\mat{P}$.
Jadi, kita mempunyai
$$p_{13}^{(2)} = p_{11}p_{13} + p_{12}p_{23} + p_{13}p_{33}\ .$$
Persamaan ini mengingatkan kita pada hasil kali titik dua vektor; kita mengambil
hasil kali titik baris pertama $ \mat {P}$ dengan kolom ketiga $ \mat {P}$.
Tepat inilah yang dilakukan untuk memperoleh entri-$1,3$ dari hasil kali $ \mat
{P}$ dengan dirinya sendiri. Secara umum, jika rantai Markov memiliki $r$
keadaan, maka
$$p_{ij}^{(2)} = \sum_{k = 1}^r p_{ik}p_{kj}\ .$$
Teorema umum berikut mudah dibuktikan menggunakan pengamatan di atas dan induksi.
\begin{theorem}\label{thm 11.1.1}
Misalkan $ \mat {P}$ adalah matriks transisi suatu rantai Markov. Entri
ke-$ij$, yaitu~$p_{ij}^{(n)}$, dari matriks~$ \mat {P}^n$ memberikan peluang
bahwa rantai Markov yang bermula dalam keadaan~$s_i$ akan berada dalam
keadaan~$s_j$ setelah $n$ langkah.
\proof
Bukti teorema ini diserahkan sebagai latihan (Latihan~\ref{exer 11.1.18}).
\end{theorem}

\begin{example}\label{exam 11.1.1.5} (lanjutan Contoh~\ref{exam 11.1.1})
Tinjau kembali cuaca di Negeri Oz. Kita mengetahui bahwa pangkat-pangkat matriks
transisi memberi kita informasi menarik tentang perkembangan proses tersebut.
Kita terutama akan tertarik pada keadaan rantai setelah banyak langkah. Program
{\bf MatrixPowers}\index{MatrixPowers (program)} menghitung pangkat-pangkat~$\mat{P}$.
\par
Kita telah menjalankan program {\bf MatrixPowers} bagi contoh Negeri Oz untuk
menghitung pangkat-pangkat berturut-turut~$\mat{P}$ dari 1~sampai~6. Hasilnya
ditampilkan dalam Tabel~\ref{table 11.1}. Perhatikan bahwa setelah enam hari,
hingga ketelitian tiga tempat desimal, prediksi cuaca kita tidak bergantung pada
cuaca hari ini. Peluang ketiga jenis cuaca R, N, dan S adalah .4, .2, dan .4
tanpa memandang keadaan awal rantai. Ini merupakan contoh jenis rantai Markov
yang disebut rantai Markov \emx {reguler}. Untuk jenis rantai ini, prediksi
jangka panjang tidak bergantung pada keadaan awal. Tidak semua rantai bersifat
reguler, tetapi ini merupakan kelas rantai penting yang akan kita pelajari
secara terperinci nanti.
\begin{table}
\centering
$$
\mat{P}^1  = \bordermatrix{
            &\mbox{Hujan}&\mbox{Cerah}&\mbox{Salju} \cr
\mbox{Hujan} & .500      & .250      & .250 \cr
\mbox{Cerah} & .500      & .000      & .500 \cr
\mbox{Salju} & .250      & .250      & .500 \cr}
$$
$$  
\mat {P}^2  = \bordermatrix{
            &\mbox{Hujan}&\mbox{Cerah}&\mbox{Salju} \cr
\mbox{Hujan} & .438     & .188       & .375 \cr
\mbox{Cerah} & .375     & .250       & .375 \cr
\mbox{Salju} & .375     & .188       & .438 \cr}
$$
$$
\mat {P}^3  = \bordermatrix{
            &\mbox{ Hujan}&\mbox{Cerah} &\mbox{Salju} \cr
\mbox{Hujan} & .406       & .203       & .391 \cr
\mbox{Cerah} & .406       & .188       & .406 \cr
\mbox{Salju} & .391       & .203       & .406 \cr}
$$
$$  
\mat {P}^4  = \bordermatrix{
            &\mbox{Hujan}&\mbox{Cerah}&\mbox{Salju} \cr
\mbox{Hujan} & .402      & .199      & .398 \cr
\mbox{Cerah} & .398      & .203      & .398 \cr
\mbox{Salju} & .398      & .199      & .402 \cr}
$$
$$
\mat {P}^5  = \bordermatrix{
           &\mbox{Hujan}&\mbox{Cerah}&\mbox{Salju} \cr
\mbox{Hujan} & .400     & .200     & .399 \cr
\mbox{Cerah} & .400     & .199     & .400 \cr
\mbox{Salju} & .399     & .200     & .400 \cr}
$$
$$  
\mat {P}^6  = \bordermatrix{
           &\mbox{Hujan}&\mbox{Cerah}&\mbox{Salju} \cr
\mbox{Hujan} & .400     & .200     & .400 \cr
\mbox{Cerah} & .400     & .200     & .400 \cr
\mbox{Salju} & .400     & .200     & .400 \cr}
$$
\caption{Pangkat-pangkat matriks transisi Negeri Oz.}
\label{table 11.1}
\end{table}
\end{example}

Sekarang kita meninjau perilaku jangka panjang suatu rantai Markov ketika
rantai bermula dalam keadaan yang dipilih menurut distribusi peluang pada
himpunan keadaan; distribusi ini akan kita sebut \emx {vektor
peluang}.\index{peluang!vektor} Vektor peluang dengan $r$ komponen adalah vektor
baris yang entri-entrinya non-negatif dan berjumlah 1. Jika $\mat {u}$ adalah
vektor peluang yang menyatakan keadaan awal suatu rantai Markov, komponen
ke-$i$ dari $\mat {u}$ kita tafsirkan sebagai peluang bahwa rantai bermula dalam
keadaan $s_i$.
\par
Dengan penafsiran keadaan awal acak ini, teorema berikut mudah dibuktikan.
\begin{theorem}\label{thm 11.1.2}
Misalkan $\mat{P}$ adalah matriks transisi suatu rantai Markov dan $\mat {u}$
adalah vektor peluang yang menyatakan distribusi awal. Maka peluang bahwa
rantai berada dalam keadaan $s_i$ setelah $n$ langkah adalah entri ke-$i$ pada
vektor
$$ \mat{u}^{(n)} = \mat{u}{\mat{P}^n}\ .$$
\proof
Bukti teorema ini diserahkan sebagai latihan (Latihan~\ref{exer 11.1.19}).
\end{theorem}

Perhatikan bahwa jika kita ingin memeriksa perilaku rantai dengan asumsi bahwa
rantai bermula dalam keadaan tertentu $s_i$, kita cukup memilih $\mat {u}$
sebagai vektor peluang yang entri ke-$i$-nya sama dengan 1 dan semua entri lain
sama dengan 0.

\begin{example}\label{exam 11.1.1.6} Dalam contoh Negeri Oz
(Contoh~\ref{exam 11.1.1}), misalkan vektor peluang awal $\mat {u}$ sama dengan
$(1/3, 1/3, 1/3)$. Kita dapat menghitung distribusi keadaan setelah tiga hari
dengan menggunakan Teorema~\ref{thm 11.1.2} dan perhitungan kita sebelumnya
atas ${\mat {P}^3}$. Kita memperoleh
\begin{eqnarray*}
{\mat {u}}^{(3)} = {\mat {u}}{\mat {P}^3} &=& \pmatrix{ 1/3,& 1/3,& 1/3}
\pmatrix{ .406 & .203 &
.391 \cr   .406 & .188 & .406 \cr .391 & .203 & .406 } \cr
&& \cr
&=& \pmatrix{ .401,& .198,& .401} \ .
\end{eqnarray*}
\end{example}

\subsection*{Contoh}
Contoh-contoh rantai Markov berikut akan digunakan dalam latihan di seluruh bab.

\begin{example}\label{exam 11.1.2}
Presiden Amerika Serikat memberi tahu orang~A tentang niatnya untuk maju atau
tidak maju dalam pemilihan berikutnya. Kemudian A meneruskan kabar itu kepada~B,
yang kemudian meneruskan pesan tersebut kepada~C, dan seterusnya, selalu kepada
orang baru. Kita menganggap terdapat peluang~$a$ bahwa seseorang akan mengubah
jawaban dari ya menjadi tidak ketika menyampaikannya kepada orang berikutnya,
dan peluang~$b$ bahwa ia akan mengubahnya dari tidak menjadi ya. Pesan, yaitu
ya atau tidak, kita pilih sebagai keadaan. Matriks transisinya adalah
$$
\mat{P} = \bordermatrix{
           & \mbox{ya} & \mbox{tidak} \cr
\mbox{ya} &      1 - a &         a \cr
\mbox{tidak}  &          b &     1 - b}\ .
$$
Keadaan awal menyatakan pilihan Presiden.
\end{example}

\begin{example}\label{exam 11.1.3}
Setiap kali seekor kuda tertentu berlomba dalam pacuan tiga kuda, peluangnya
untuk menang adalah~1/2, untuk menempati urutan kedua 1/4, dan untuk menempati
urutan ketiga 1/4, terlepas dari hasil pacuan sebelumnya. Kita memiliki proses
percobaan independen, tetapi proses ini juga dapat ditinjau dari sudut pandang
teori rantai Markov. Matriks transisinya adalah
$$
\mat{P} = \bordermatrix{
        & \mbox{W} & \mbox{P} & \mbox{S} \cr
\mbox{W} &      .5 &      .25 &      .25 \cr
\mbox{P} &      .5 &      .25 &      .25 \cr
\mbox{S} &      .5 &      .25 &      .25}\ .
$$
\end{example}
\vskip -3pt
\begin{example}\label{exam 11.1.4}
Pada Zaman Kegelapan, Harvard, Dartmouth, dan Yale hanya menerima mahasiswa
laki-laki. Anggaplah bahwa pada masa itu 80~persen anak laki-laki alumni
Harvard kuliah di Harvard dan sisanya di Yale; 40~persen anak laki-laki alumni
Yale kuliah di Yale dan sisanya terbagi rata antara Harvard dan Dartmouth;
sedangkan dari anak laki-laki alumni Dartmouth, 70~persen kuliah di Dartmouth,
20~persen di Harvard, dan 10~persen di Yale. Kita membentuk rantai Markov dengan
matriks transisi
$$
\mat{P} = \bordermatrix{
        & \mbox{H} & \mbox{Y} & \mbox{D} \cr
\mbox{H} &      .8 &      .2 &       0 \cr
\mbox{Y} &      .3 &      .4 &      .3 \cr 
\mbox{D} &      .2 &      .1 &      .7}\ .
$$
\end{example}
\vskip -3pt
\begin{example}\label{exam 11.1.5}
Ubahlah Contoh~\ref{exam 11.1.4} dengan menganggap bahwa anak laki-laki seorang
alumni Harvard selalu kuliah di Harvard. Matriks transisinya sekarang adalah
$$
\mat{P} = \bordermatrix{
        & \mbox{H} & \mbox{Y} & \mbox{D} \cr
\mbox{H} &       1 &       0 &       0 \cr
\mbox{Y} &      .3 &      .4 &      .3 \cr
\mbox{D} &      .2 &      .1 &      .7}\ .
$$
\end{example}
\vskip -3pt
\begin{example}\label{exam 11.1.6}\index{model Ehrenfest}
\index{difusi gas!model Ehrenfest}
(Model Ehrenfest) Berikut ini merupakan kasus khusus suatu model yang disebut
model Ehrenfest,\footnote{P. and T. Ehrenfest, ``\"{U}ber zwei bekannte Einw\"{a}nde
gegen das
Boltzmannsche  H-Theorem," \emx {Physikalishce Zeitschrift,} vol.~8 (1907),
pp.~311-314.}\index{EHRENFEST, P.}
\index{EHRENFEST, T}
yang telah digunakan untuk menjelaskan difusi gas. Model umumnya akan dibahas
secara terperinci dalam Bagian~\ref{sec 11.5}. Kita memiliki dua guci yang
bersama-sama memuat empat bola. Pada setiap langkah, salah satu dari keempat
bola dipilih secara acak dan dipindahkan dari guci tempatnya berada ke guci
lain. Banyaknya bola dalam guci pertama kita pilih sebagai keadaan. Matriks
transisinya adalah
$$
 \mat {P} = \bordermatrix{
  &  0  &  1  &  2  &  3  &  4  \cr
0 &  0  &  1  &  0  &  0  &  0  \cr
1 & 1/4 &  0  & 3/4 &  0  &  0  \cr
2 &  0  & 1/2 &  0  & 1/2 &  0  \cr
3 &  0  &  0  & 3/4 &  0  & 1/4 \cr
4 &  0  &  0  &  0  &  1  &  0 \cr}\ . 
$$
\end{example}

\begin{example}\label{exam 11.1.7}\index{gen}
(Model Gen) Jenis pewarisan sifat paling sederhana pada hewan terjadi ketika
suatu sifat dikendalikan oleh sepasang gen yang masing-masing dapat terdiri
atas dua jenis, misalnya G~dan~g. Seorang individu dapat memiliki kombinasi GG,
Gg (yang secara genetik sama dengan gG), atau gg. Sangat sering tipe GG dan Gg
tidak dapat dibedakan dari penampilannya; dalam hal ini kita mengatakan bahwa
gen G~mendominasi gen g. Individu disebut \emx {dominan} jika memiliki gen~GG,
\emx {resesif} jika memiliki gg, dan \emx {hibrida} jika memiliki campuran Gg.
\par
Dalam perkawinan dua hewan, keturunannya mewarisi satu gen dari pasangan gen
tersebut dari setiap induk, dan asumsi dasar genetika adalah bahwa gen-gen ini
dipilih secara acak dan saling independen. Asumsi ini menentukan peluang
munculnya setiap jenis keturunan. Keturunan dari dua induk dominan murni harus
dominan, keturunan dari dua induk resesif harus resesif, dan keturunan dari satu
induk dominan dan satu induk resesif harus hibrida.
\par
Dalam perkawinan hewan dominan dan hibrida, setiap keturunan harus memperoleh
gen G~dari hewan pertama dan memiliki peluang yang sama untuk memperoleh G~atau
g dari hewan kedua. Jadi, peluang memperoleh keturunan dominan sama dengan
peluang memperoleh keturunan hibrida. Demikian pula, dalam perkawinan hewan
resesif dan hibrida, peluang memperoleh keturunan resesif sama dengan peluang
memperoleh keturunan hibrida. Dalam perkawinan dua hibrida, keturunannya
memiliki peluang yang sama untuk memperoleh G~atau g dari setiap induk. Jadi,
peluangnya adalah 1/4 untuk GG, 1/2 untuk Gg, dan 1/4 untuk gg.
\par
Tinjaulah proses perkawinan berkelanjutan. Kita mulai dengan individu yang
karakter genetiknya diketahui dan mengawinkannya dengan hibrida. Kita
menganggap terdapat sedikitnya satu keturunan. Satu keturunan dipilih secara
acak dan dikawinkan dengan hibrida, lalu proses ini diulangi selama beberapa
generasi. Jenis genetik keturunan terpilih dalam generasi-generasi berturut-turut
dapat dinyatakan oleh rantai Markov. Keadaannya adalah dominan, hibrida, dan
resesif, yang masing-masing ditunjukkan oleh GG, Gg, dan gg.
\par
Peluang transisinya adalah
$$
\mat{P} = \bordermatrix{
          &\mbox{GG} & \mbox{Gg} & \mbox{gg} \cr
\mbox{GG} &  .5      & .5        &   0       \cr
\mbox{Gg} & .25      & .5        & .25       \cr
\mbox{gg} &   0      & .5        &  .5       }\ . 
$$
\end{example}

\begin{example}\label{exam 11.1.8}
Ubahlah Contoh~\ref{exam 11.1.7} sebagai berikut: Alih-alih mengawinkan
keturunan tertua dengan hibrida, kita mengawinkannya dengan individu dominan.
Matriks transisinya adalah
$$
\mat{P} = \bordermatrix{
   & \mbox{GG} & \mbox{Gg} &\mbox{gg} \cr
\mbox{GG} &  1 &  0 &  0 \cr
\mbox{Gg} & .5 & .5 &  0 \cr
\mbox{gg} &  0 &  1 &  0}\ .
$$
\end{example}

\begin{example}\label{exam 11.1.9}
Kita mulai dengan dua hewan berjenis kelamin berbeda, mengawinkan keduanya,
memilih dua keturunan mereka yang berjenis kelamin berbeda, lalu mengawinkan
keduanya, dan seterusnya. Untuk menyederhanakan contoh, kita akan menganggap
bahwa sifat yang ditinjau tidak bergantung pada jenis kelamin.
\par
Di sini suatu keadaan ditentukan oleh sepasang hewan. Jadi, keadaan-keadaan
proses kita adalah: $s_1 = (\mbox{GG},\mbox{GG})$, $s_2 =
(\mbox{GG},\mbox{Gg})$, $s_3 = (\mbox{GG},\mbox{gg})$, $s_4 =
(\mbox{Gg},\mbox{Gg})$, $s_5 = (\mbox{Gg},\mbox{gg})$, dan $s_6 =
(\mbox{gg},\mbox{gg})$.
\par
Kita mengilustrasikan perhitungan peluang transisi melalui keadaan~$s_2$.
Ketika proses berada dalam keadaan ini, satu induk memiliki gen~GG dan induk
lainnya Gg. Jadi, peluang memperoleh keturunan dominan adalah~1/2. Kemudian
peluang transisi ke~$s_1$ (memilih dua individu dominan) adalah~1/4, peluang
transisi ke~$s_2$ adalah~1/2, dan peluang transisi ke~$s_4$ adalah~1/4.
Keadaan-keadaan lain diperlakukan dengan cara yang sama. Matriks transisi
rantai ini adalah:
\par

$$
{\mat{P}^1} = \bordermatrix{            
&\mbox{GG,GG}&\mbox{GG,Gg}&\mbox{GG,gg}&\mbox{Gg,Gg}&\mbox{Gg,gg}&\mbox{gg,gg}\cr
\mbox{GG,GG} & 1.000 & .000 &  .000 &  .000 &  .000 &  .000\cr
\mbox{GG,Gg} &  .250 & .500 &  .000 &  .250 &  .000 &  .000\cr
\mbox{GG,gg} &  .000 & .000 &  .000 & 1.000 &  .000 &  .000\cr
\mbox{Gg,Gg} &  .062 & .250 &  .125 &  .250 &  .250 &  .062\cr
\mbox{Gg,gg} &  .000 & .000 &  .000 &  .250 &  .500 &  .250\cr
\mbox{gg,gg} &  .000 & .000 &  .000 &  .000 &  .000 &  1.000}\ .
$$
 
\end{example}

\begin{example} \label{exam 11.1.10}\index{batu loncatan}
(Model Batu Loncatan) Contoh terakhir kita merupakan contoh lain yang telah
digunakan dalam kajian genetika. Contoh ini disebut model \emx {batu
loncatan}.\footnote{S.
Sawyer, ``Results for The Stepping Stone Model for Migration in Population
Genetics," \emx {Annals of Probability,} vol.~4 (1979),
pp.~699--728.}\index{SAWYER, S.} Dalam model ini kita memiliki larik persegi
$n$-kali-$n$, dan setiap persegi pada awalnya memiliki salah satu dari $k$
warna berbeda. Pada setiap langkah, sebuah persegi dipilih secara acak. Persegi
ini kemudian memilih salah satu dari delapan tetangganya secara acak dan
mengambil warna tetangga tersebut. Untuk menghindari masalah batas, kita
menganggap bahwa jika suatu persegi $S$ berada pada batas kiri, misalnya, tetapi
tidak di sudut, persegi itu bertetangga dengan persegi $T$ pada batas kanan di
baris yang sama dengan $S$, dan $S$ juga bertetangga dengan persegi tepat di
atas dan di bawah $T$. Asumsi serupa dibuat untuk persegi pada batas atas dan
bawah. Persegi sudut kiri atas bertetangga dengan tiga tetangga yang jelas,
yaitu persegi di bawahnya, di kanannya, dan secara diagonal di kanan bawahnya.
Persegi ini memiliki lima tetangga lain, yaitu: tiga persegi sudut lainnya,
persegi di bawah sudut kanan atas, dan persegi di kanan sudut kiri bawah.
Ketiga sudut lainnya, dengan cara serupa, juga memiliki delapan tetangga.
(Ketetanggaan ini jauh lebih mudah dipahami jika kita membayangkan larik
dibentuk menjadi silinder dengan merekatkan sisi atas dan bawah, lalu membentuk
silinder itu menjadi donat dengan merekatkan kedua batas lingkarannya.) Dengan
ketetanggaan ini, setiap persegi dalam larik bertetangga dengan tepat delapan
persegi lain.
\par
Suatu keadaan dalam rantai Markov ini merupakan deskripsi warna setiap persegi.
Untuk rantai Markov ini, banyaknya keadaan adalah~$k^{n^2}$, yang bahkan untuk
larik persegi kecil pun sangat besar. Ini merupakan contoh rantai Markov yang
mudah disimulasikan, tetapi sulit dianalisis melalui matriks transisinya. Program
{\bf SteppingStone}\index{SteppingStone (program)} menyimulasikan rantai ini.
Kita memulai dengan konfigurasi awal acak dari dua warna dengan $n = 20$ dan
menampilkan hasil setelah proses berjalan selama beberapa waktu dalam
Gambar~\ref{fig 11.2}.
\noindent
\par
Ini merupakan contoh rantai Markov \emx {menyerap}. Jenis rantai ini akan
dipelajari dalam Bagian~\ref{sec 11.2}. Salah satu teorema yang dibuktikan dalam
bagian tersebut, ketika diterapkan pada contoh ini, mengakibatkan bahwa dengan
peluang 1, semua batu pada akhirnya akan memiliki warna yang sama. Dengan
mengamati program berjalan, Anda dapat melihat wilayah-wilayah terbentuk dan
pertempuran berkembang untuk menentukan warna mana yang bertahan. Pada sebarang
saat, peluang bahwa suatu warna tertentu akan menang sama dengan proporsi larik
yang memiliki warna tersebut. Anda diminta membuktikan hal ini dalam
Latihan~\ref{sec 11.2}.\ref{exer 11.2.31}.
\end{example}

\putfig{2truein}{PSfig11-1}{Keadaan awal model batu loncatan.}{fig 11.1}

%\begin{figure}
%\centerline{\epsfxsize=2truein 
%\epsffile{PSfig11-1}}
%\caption{Initial state of the stepping stone model.}\label{fig 11.1}
%\end{figure}
\nopagebreak[4]
\putfig{2truein}{PSfig11-2}{Keadaan model batu loncatan setelah 10,000 langkah.}{fig 11.2}

%\begin{figure}
%\centerline{\epsfxsize=2truein 
%\epsffile{PSfig11-2}}
%\caption{State of the stepping stone model after 10,000 steps.}\label{fig 11.2}
%\end{figure}

\exercises
\begin{LJSItem}

\i\label{exer 11.1.1} Di Negeri Oz sedang hujan. Tentukan pohon dan ukuran
pohon untuk cuaca tiga hari berikutnya. Carilah $\mat {w}^{(1)},  \mat
{w}^{(2)},$ dan $ \mat {w}^{(3)}$, lalu bandingkan dengan hasil yang diperoleh
dari $ \mat {P},~\mat {P}^2,$ dan~$ \mat {P}^3$.

\i\label{exer 11.1.2} Dalam Contoh~\ref{exam 11.1.2}, misalkan $a = 0$ dan $b =
1/2$. Carilah $ \mat {P},~ \mat {P}^2,$ dan $ \mat {P}^3.$ Bagaimanakah bentuk $
\mat {P}^n$? Apa yang terjadi pada~$ \mat {P}^n$ ketika $n$ menuju tak
hingga? Tafsirkan hasil ini.

\i\label{exer 11.1.3} Dalam Contoh~\ref{exam 11.1.3}, carilah $\mat{P}$,~$
\mat {P}^2,$ dan~$ \mat {P}^3.$ Tentukan $ \mat {P}^n$.

\i\label{exer 11.1.4} Untuk Contoh~\ref{exam 11.1.4}, carilah peluang bahwa
cucu laki-laki seorang alumni Harvard kuliah di Harvard.

\i\label{exer 11.1.5} Dalam Contoh~\ref{exam 11.1.5}, carilah peluang bahwa
cucu laki-laki seorang alumni Harvard kuliah di Harvard.

\i\label{exer 11.1.6} Dalam Contoh~\ref{exam 11.1.7}, anggaplah kita mulai
dengan hibrida yang dikawinkan dengan hibrida. Carilah $ \mat {u}^{(1)},$~$
\mat {u}^{(2)},$ dan~$ \mat {u}^{(3)}.$ Apakah bentuk $ \mat {u}^{(n)}$?

\i\label{exer 11.1.7} Carilah matriks $\mat{ P}^2,~\mat {P}^3,~\mat {P}^4,$
dan~$ \mat {P}^n$ bagi rantai Markov yang ditentukan oleh matriks transisi $
\mat {P} = \pmatrix{ 1 & 0 \cr 0 & 1 \cr}$. Lakukan hal yang sama bagi matriks
transisi $ \mat {P} = \pmatrix{ 0 & 1 \cr 1 & 0 \cr}$. Tafsirkan apa yang
terjadi dalam setiap proses ini.

\i\label{exer 11.1.8} Sebuah mesin hitung tertentu hanya menggunakan angka
0~dan~1. Mesin tersebut dimaksudkan untuk meneruskan salah satu angka ini
melalui beberapa tahap. Namun, pada setiap tahap terdapat peluang~$p$ bahwa
angka yang memasuki tahap tersebut akan berubah ketika keluar, dan peluang $q
= 1 - p$ bahwa angka itu tidak berubah. Bentuklah rantai Markov untuk
menyatakan proses transmisi dengan mengambil angka 0~dan~1 sebagai keadaan.
Apakah matriks peluang transisinya?

\i\label{exer 11.1.9} Untuk rantai Markov dalam Latihan~\ref{exer 11.1.8},
gambarlah pohon dan tetapkan ukuran pohon dengan menganggap bahwa proses bermula
dalam keadaan~0 dan melewati dua tahap transmisi. Berapakah peluang bahwa
setelah dua tahap mesin menghasilkan angka~0 (yakni, angka yang benar)?
Berapakah peluang bahwa mesin tidak pernah mengubah angka~0? Sekarang misalkan
$p = .1$. Dengan menggunakan program {\bf MatrixPowers}, hitunglah pangkat
ke-100 matriks transisi. Tafsirkan entri-entri matriks ini. Ulangi dengan $p =
.2$. Mengapa pangkat ke-100 tersebut tampak sama?

\i\label{exer 11.1.10} Ubahlah program {\bf MatrixPowers} agar mencetak rataan
$ \mat {A}_n$ dari pangkat-pangkat $\mat {P}^n$, untuk $n = 1$ sampai $N$.
Cobalah program Anda pada contoh Negeri Oz dan bandingkan $\mat {A}_n$~dengan
$\mat {P}^n.$

\i\label{exer 11.1.11} Anggaplah profesi seorang pria dapat diklasifikasikan
sebagai tenaga profesional, pekerja terampil, atau pekerja tidak terampil.
Anggaplah bahwa dari anak laki-laki tenaga profesional, 80~persen menjadi
tenaga profesional, 10~persen menjadi pekerja terampil, dan 10~persen menjadi
pekerja tidak terampil. Dari anak laki-laki pekerja terampil, 60~persen menjadi
pekerja terampil, 20~persen menjadi tenaga profesional, dan 20~persen menjadi
pekerja tidak terampil. Terakhir, dari anak laki-laki pekerja tidak terampil,
50~persen menjadi pekerja tidak terampil dan masing-masing 25~persen berada
dalam dua kategori lainnya. Anggaplah setiap pria memiliki sedikitnya satu
anak laki-laki, lalu bentuklah rantai Markov dengan mengikuti profesi seorang
anak laki-laki yang dipilih secara acak dari suatu keluarga selama beberapa
generasi. Susunlah matriks peluang transisinya. Carilah peluang bahwa cucu
laki-laki yang dipilih secara acak dari seorang pekerja tidak terampil menjadi
tenaga profesional.

\i\label{exer 11.1.12} Dalam Latihan~\ref{exer 11.1.11}, kita menganggap setiap
pria memiliki seorang anak laki-laki. Sebagai gantinya, anggaplah peluang bahwa
seorang pria memiliki sedikitnya satu anak laki-laki adalah~.8. Bentuklah
rantai Markov dengan empat keadaan. Jika seorang pria memiliki anak laki-laki,
peluang bahwa anak tersebut berada dalam profesi tertentu sama seperti dalam
Latihan~\ref{exer 11.1.11}. Jika tidak ada anak laki-laki, proses berpindah ke
keadaan keempat yang menyatakan keluarga yang garis keturunan laki-lakinya
telah berakhir. Carilah matriks peluang transisi dan peluang bahwa cucu
laki-laki yang dipilih secara acak dari seorang pekerja tidak terampil menjadi
tenaga profesional.

\i\label{exer 11.1.14} Tulislah program untuk menghitung $\mat {u}^{(n)}$ jika
$\mat {u}$ dan $\mat{P}$ diberikan. Gunakan program ini untuk menghitung $\mat
{u}^{(10)}$ bagi contoh Negeri Oz, dengan $\mat {u} = (0, 1, 0)$, dan dengan
$\mat {u} = (1/3, 1/3, 1/3)$.

\i\label{exer 11.1.15} Dengan menggunakan program {\bf  MatrixPowers}, carilah
$\mat {P}^1$ sampai $\mat {P}^6$ untuk Contoh~\ref{exam 11.1.7} dan
\ref{exam 11.1.8}. Periksalah apakah Anda dapat memprediksi peluang jangka
panjang menemukan proses dalam setiap keadaan bagi contoh-contoh ini.

\i\label{exer 11.1.16} Tulislah program untuk menyimulasikan hasil rantai Markov
setelah $n$~langkah jika keadaan awal dan matriks transisi $\mat{P}$ diberikan
sebagai data (lihat Contoh~\ref{exam 11.1.10}). Simpanlah program ini untuk
digunakan dalam soal-soal berikutnya.

\i\label{exer 11.1.17} Ubahlah program dalam Latihan~\ref{exer 11.1.16} agar
mencatat proporsi waktu yang dihabiskan dalam setiap keadaan selama $n$~langkah.
Jalankan program yang telah diubah untuk keadaan awal yang berbeda pada
Contoh~\ref{exam 11.1.1} dan Contoh~\ref{exam 11.1.6}. Jika $n$ besar, apakah
keadaan awal memengaruhi proporsi waktu yang dihabiskan dalam setiap keadaan?

\i\label{exer 11.1.18} Buktikan Teorema~\ref{thm 11.1.1}.

\i\label{exer 11.1.19} Buktikan Teorema~\ref{thm 11.1.2}.

\i\label{exer 11.1.20} Tinjaulah proses berikut. Kita memiliki dua koin; satu
koin adil, sedangkan koin lainnya memiliki sisi kepala pada kedua sisinya.
Kedua koin ini kita berikan kepada seorang teman, yang memilih salah satunya
secara acak (masing-masing dengan peluang 1/2). Selama sisa proses, ia hanya
menggunakan koin pilihannya. Kemudian ia melempar koin itu berkali-kali dan
melaporkan hasilnya. Kita memandang proses ini semata-mata terdiri atas apa
yang ia laporkan kepada kita.
\begin{enumerate}
\item Jika ia melaporkan kepala pada lemparan ke-$n$, berapakah peluang kepala
muncul pada lemparan ke-$(n+1)$?

\item Pandanglah proses ini memiliki dua keadaan, kepala dan ekor. Dengan
menghitung tiga peluang transisi lain yang serupa dengan peluang pada bagian
(a), tuliskan sebuah ``matriks transisi" untuk proses ini.

\item Sekarang anggaplah proses berada dalam keadaan ``kepala" pada lemparan
ke-$(n-1)$ maupun ke-$n$. Carilah peluang bahwa kepala muncul pada lemparan
ke-$(n+1)$.

\item Apakah proses ini merupakan rantai Markov?
\end{enumerate}

\end{LJSItem}

\section{Rantai Markov Penyerap}\label{sec 11.2}
Rantai Markov paling baik dipelajari dengan meninjau jenis-jenis khusus rantai
Markov. Jenis pertama yang akan kita pelajari disebut \emx {rantai Markov
penyerap.}\index{rantai Markov!penyerap}\index{rantai Markov penyerap}
\begin{definition}
Suatu keadaan~$s_i$ dalam rantai Markov disebut \emx
{penyerap}\index{keadaan!penyerap}\index{keadaan penyerap} jika keadaan tersebut
tidak mungkin ditinggalkan (yakni,
$p_{ii} = 1$). Suatu rantai Markov disebut \emx {penyerap} jika memiliki
sedikitnya satu keadaan penyerap dan, dari setiap keadaan, terdapat kemungkinan
untuk mencapai suatu keadaan penyerap (tidak harus dalam satu langkah).
\end{definition}
\begin{definition}
Dalam rantai Markov penyerap, keadaan yang bukan penyerap disebut
\emx{transien.}\index{keadaan!transien}\index{keadaan transien}
\end{definition}

\subsection*{Perjalanan Acak Seorang Pemabuk}\index{contoh perjalanan acak pemabuk}
\begin{example}\label{exam 11.2.1}
Seorang pria berjalan sepanjang ruas empat blok di Park Avenue (lihat
Gambar~\ref{fig 11.3}). Jika ia berada di persimpangan 1, 2, atau 3, ia berjalan
ke kiri atau ke kanan dengan peluang yang sama. Ia melanjutkan perjalanan
hingga mencapai persimpangan~4, tempat sebuah bar berada, atau persimpangan~0,
tempat rumahnya berada. Setelah mencapai rumah ataupun bar, ia tetap di sana.

\par
Kita membentuk rantai Markov dengan keadaan 0,~1, 2, 3, dan~4. Keadaan 0~dan~4
merupakan keadaan penyerap. Matriks transisinya adalah

$$
\mat{P} =\bordermatrix{
  &  0  &  1  &  2  &  3  &  4  \cr
0 &  1  &  0  &  0  &  0  &  0  \cr
1 & 1/2 &  0  & 1/2 &  0  &  0  \cr
2 &  0  & 1/2 &  0  & 1/2 &  0  \cr
3 &  0  &  0  & 1/2 &  0  & 1/2 \cr
4 &  0  &  0  &  0  &  0  &  1 \cr}\ .
$$
Keadaan 1,~2, dan~3 merupakan keadaan transien, dan dari masing-masing keadaan
ini terdapat kemungkinan untuk mencapai keadaan penyerap 0~dan~4. Dengan
demikian, rantai tersebut merupakan rantai penyerap. Ketika suatu proses
mencapai keadaan penyerap, kita akan mengatakan bahwa proses itu
\emx{terserap}.
\end{example}

Pertanyaan paling jelas mengenai rantai seperti ini adalah: berapakah peluang
bahwa proses tersebut pada akhirnya mencapai suatu keadaan penyerap?
Pertanyaan menarik lainnya meliputi: (a)~Berapakah peluang bahwa proses berakhir
pada suatu keadaan penyerap tertentu? (b)~Secara rata-rata, berapa lama waktu
yang diperlukan hingga proses terserap? (c)~Secara rata-rata, berapa kali proses
berada dalam setiap keadaan transien? Secara umum, jawaban semua pertanyaan ini
bergantung pada keadaan awal proses maupun peluang transisinya.

\putfig{4.5truein}{PSfig11-3}{Perjalanan acak seorang pemabuk.}{fig 11.3}

%\begin{figure}
%\centerline{\epsfxsize=4.5truein 
%\epsffile{PSfig11-3}}
%\caption{Drunkard's walk.}\label{fig 11.3}
%\end{figure}

\subsection*{Bentuk Kanonik}\index{bentuk kanonik rantai Markov\\ penyerap}
Tinjaulah sebarang rantai Markov penyerap. Nomori ulang keadaan-keadaannya agar
keadaan transien ditempatkan terlebih dahulu. Jika terdapat $r$ keadaan
penyerap dan $t$ keadaan transien, matriks transisinya mempunyai \emx{bentuk
kanonik} berikut

\[
\offinterlineskip
\mat{P}\;= \bordermatrix{      
                               &\hbox{TR.}&\omit\hfil&\hbox{ABS.}\cr
           \hbox{TR.}\bigstrut &\mat{Q}   &\srule    &\mat{R}    \cr
\middlehrule{1}{1}
           \hbox{ABS.}\bigstrut&\mat{0}   &\srule    &\mat{I}}
\] 

Di sini $\mat{I}$ adalah matriks identitas berukuran $r$-kali-$r$, $\mat{0}$
adalah matriks nol berukuran $r$-kali-$t$, $\mat{R}$ adalah matriks tak nol
berukuran $t$-kali-$r$, dan $\mat{Q}$ adalah matriks berukuran $t$-kali-$t$.
Sebanyak $t$ keadaan pertama bersifat transien dan $r$ keadaan terakhir
bersifat menyerap.
\par
Dalam Bagian~\ref{sec 11.1}, kita melihat bahwa entri~$p_{ij}^{(n)}$ dari matriks
$\mat{P}^n$ adalah peluang berada dalam keadaan~$s_j$ setelah $n$ langkah jika
rantai dimulai dalam keadaan~$s_i$. Suatu argumen aljabar matriks standar
menunjukkan bahwa $\mat{P}^n$ berbentuk
 \[
\offinterlineskip
\mat{P}^n\;= \bordermatrix{      
                      &\hbox{TR.}&\omit\hfil&\hbox{ABS.}\cr
  \hbox{TR.}\bigstrut &\mat{Q}^n &\srule    &\ast       \cr
\middlehrule{1}{1}
  \hbox{ABS.}\bigstrut&\mat{0}   &\srule    &\mat{I}}
\] 
di mana tanda bintang $*$ menyatakan matriks berukuran $t$-kali-$r$ di sudut
kanan atas~$\mat{P}^n.$ (Submatriks ini dapat dituliskan dalam $\mat{Q}$ dan
$\mat{R}$, tetapi ungkapannya rumit dan belum diperlukan.) Bentuk~$\mat{P}^n$
menunjukkan bahwa entri-entri $\mat{Q}^n$ memberikan peluang berada dalam setiap
keadaan transien setelah $n$ langkah untuk setiap kemungkinan keadaan awal
transien. Dalam teorema pertama, kita membuktikan bahwa peluang berada dalam
keadaan transien setelah $n$ langkah mendekati nol. Dengan demikian, setiap
entri~$\mat{ Q}^n$ harus mendekati nol ketika $n$ menuju tak hingga (yakni,
$\mat{Q}^n \to \mat{ 0}$).

\subsection*{Peluang Penyerapan}
\begin{theorem}\label{thm 11.2.1}
Dalam rantai Markov penyerap, peluang bahwa proses akan terserap adalah~1
(yakni, $\mat{Q}^n \to \mat{0}$ ketika $n \to \infty$).

\proof
Dari setiap keadaan tak menyerap $s_j$, terdapat kemungkinan untuk mencapai
suatu keadaan penyerap. Misalkan $m_j$ adalah jumlah langkah minimum yang
diperlukan untuk mencapai keadaan penyerap jika dimulai dari $s_j$. Misalkan
$p_j$ adalah peluang bahwa, jika dimulai dari $s_j$, proses tidak mencapai
keadaan penyerap dalam $m_j$ langkah. Maka $p_j <1$. Misalkan $m$ adalah yang
terbesar di antara semua $m_j$, dan $p$ yang terbesar di antara semua $p_j$.
Peluang tidak terserap dalam $m$ langkah paling besar $p$, dalam $2m$ langkah
paling besar $p^2$, dan seterusnya. Karena $p<1$, peluang-peluang ini menuju 0.
Karena peluang tidak terserap dalam $n$ langkah menurun secara monoton,
peluang-peluang tersebut juga menuju 0; oleh karena itu
$\lim_{n \rightarrow \infty } \mat{Q}^n = 0.$
\end{theorem}

\subsection*{Matriks Fundamental}
\begin{theorem}\label{thm 11.2.2}
Untuk suatu rantai Markov penyerap, matriks $\mat{I} - \mat{Q}$ memiliki invers
$\mat{N}$ dan
$\mat{N} =\mat{I} + \mat{Q} + \mat{Q}^{2} + \cdots\ $. Entri ke-$ij$, yaitu
$n_{ij}$, dari matriks $\mat{N}$ adalah nilai harapan banyaknya rantai berada
dalam keadaan $s_j$ jika rantai dimulai dalam keadaan $s_i$. Keadaan awal turut
dihitung jika $i = j$.

\proof  
Misalkan $(\mat{I} - \mat{Q})\mat{x}~=~0;$ yakni
$\mat{x}~=~\mat{Q}\mat{x}.$ Dengan mengiterasikan persamaan ini, kita memperoleh
$\mat{x}~=~\mat{Q}^{n}\mat x.$ Karena $\mat{Q}^{n} \rightarrow \mat{0}$, kita
memiliki $\mat{Q}^n\mat{x} \rightarrow \mat{0}$, sehingga
$\mat{x}~=~\mat{0}$. Jadi, $(\mat{I} - \mat{Q})^{-1}~=~\mat{N}$ ada. Selanjutnya,
perhatikan bahwa
$$
(\mat{I} - \mat{Q}) (\mat{I} + \mat{Q} + \mat{Q}^2 + \cdots + \mat{Q}^n) =
\mat{I} -
\mat{Q}^{n + 1}\ .
$$
Dengan demikian, mengalikan kedua ruas dengan $\mat{N}$ memberikan
$$
\mat{I} + \mat{Q} + \mat{Q}^2 + \cdots + \mat{Q}^n = \mat{N} (\mat{I} -
\mat{Q}^{n + 1})\ .
$$
Dengan membiarkan $n$ menuju tak hingga, kita memperoleh
$$
\mat{N} = \mat{I} + \mat{Q} + \mat{Q}^2 + \cdots\ .
$$

Misalkan $s_i$ dan $s_j$ adalah dua keadaan transien, dan untuk seluruh sisa
bukti anggaplah $i$ dan $j$ tetap. Misalkan $X^{(k)}$ adalah peubah acak yang
bernilai 1 jika rantai berada dalam keadaan $s_j$ setelah $k$ langkah, dan
bernilai 0 selainnya. Untuk setiap $k$, peubah acak ini bergantung pada $i$
maupun $j$; demi kejelasan, kita memilih untuk tidak menampilkan ketergantungan
tersebut secara eksplisit. Kita memiliki
$$
P(X^{(k)} = 1) = q_{ij}^{(k)}\ ,
$$
dan
$$
P(X^{(k)} = 0) = 1 - q_{ij}^{(k)}\ ,
$$
di mana $q_{ij}^{(k)}$ adalah entri ke-$ij$ dari $\mat{Q}^k$. Persamaan-persamaan
ini juga berlaku untuk $k = 0$ karena $\mat{Q}^0 = \mat{I}$. Oleh karena itu,
karena $X^{(k)}$ adalah peubah acak 0-1, $E(X^{(k)}) = q_{ij}^{(k)}$.
\par
Nilai harapan banyaknya rantai berada dalam keadaan $s_j$ selama $n$ langkah
pertama, jika rantai dimulai dalam keadaan $s_i$, jelas sama dengan
$$E\Bigl(X^{(0)} + X^{(1)} + \cdots + X^{(n)} \Bigr) = q_{ij}^{(0)} +
q_{ij}^{(1)} +
\cdots + q_{ij}^{(n)}\ .$$
Dengan membiarkan $n$ menuju tak hingga, kita memperoleh
$$
E\Bigl(X^{(0)} + X^{(1)} + \cdots \Bigr) = q_{ij}^{(0)} +
q_{ij}^{(1)} + \cdots  = n_{ij} \ .
$$
\end{theorem}

\begin{definition}
Untuk rantai Markov penyerap~$\mat{P}$, matriks $\mat{N} = (\mat{I} -
\mat{Q})^{-1}$ disebut \emx {matriks fundamental} bagi~$\mat{P}$.
\index{matriks fundamental}\index{matriks!fundamental} Entri~$n_{ij}$ dari
$\mat{N}$ memberikan nilai harapan banyaknya proses berada dalam keadaan
transien~$s_j$ jika proses dimulai dalam keadaan transien~$s_i$.
\end{definition}

\begin{example}\label{exam 11.2.2}
\index{contoh perjalanan acak pemabuk}(Lanjutan Contoh~\ref{exam 11.2.1})
Dalam contoh perjalanan acak seorang pemabuk, matriks transisi dalam bentuk
kanonik adalah
\[
\offinterlineskip
\mat{P}\;= \bordermatrix{&
               \hbox{1} &\hbox{2}&\hbox{3}&\omit\hfil&\hbox{0}&\hbox{4}\cr
\hbox{1}\strut  &  0    &1/2 &  0     &  \srule  & 1/2    &  0     \cr
\hbox{2}\strut  &1/2    &  0 &1/2     &  \srule  & 0      &  0     \cr
\hbox{3}\strut  &  0    &1/2 &  0     &  \srule  & 0      & 1/2\cr
\middlehrule{3}{2}
\hbox{0}\strut  &  0    &  0     &  0     &  \srule  & 1      &  0     \cr
\hbox{4}\strut  &  0    &  0     &  0     &  \srule  & 0      &  1}\ .
\]
Dari sini kita melihat bahwa matriks $\mat{Q}$ adalah
$$
\mat{Q} = \pmatrix{
0 & 1/2 & 0 \cr
1/2 & 0 & 1/2 \cr
0 & 1/2 & 0 \cr}\ ,
$$
dan
$$
\mat{I} - \mat{Q} = \pmatrix{
1 & -1/2 & 0 \cr
-1/2 & 1 & -1/2 \cr
0 & -1/2 & 1 \cr}\ .
$$
Dengan menghitung $(\mat{I} - \mat{Q})^{-1}$, kita memperoleh
$$
\vspace{7pt}\hbox{$\mat{N} = (\mat{I} - \mat{Q})^{-1} = {}$} \bordermatrix{
  & 1 & 2 & 3 \cr
1 & 3/2 & 1 & 1/2 \cr
2 & 1 & 2 & 1 \cr
3 & 1/2 & 1 & 3/2 \cr}\ .
$$
Dari baris tengah~$\mat{N}$, kita melihat bahwa jika proses dimulai dalam
keadaan 2, nilai harapan banyaknya proses berada dalam keadaan 1, 2, dan 3
sebelum terserap berturut-turut adalah 1, 2, dan 1.
\end{example}
 
\subsection*{Waktu hingga Penyerapan}\index{waktu hingga penyerapan}
Sekarang kita meninjau pertanyaan berikut: jika rantai dimulai dalam keadaan
$s_i$, berapakah nilai harapan jumlah langkah sebelum rantai terserap? Jawabannya
diberikan oleh teorema berikut.

\begin{theorem}\label{thm 11.2.2.5} Misalkan $t_i$ adalah nilai harapan jumlah
langkah sebelum rantai terserap jika rantai dimulai dalam keadaan $s_i$, dan
misalkan $\mat{t}$ adalah vektor kolom yang entri ke-$i$-nya adalah $t_i$. Maka
$$\mat{t} = \mat{N}\mat{c}\ ,$$
di mana $\mat{c}$ adalah vektor kolom yang semua entrinya bernilai 1.

\proof
Jika semua entri pada baris ke-$i$ dari~$\mat{N}$ dijumlahkan, kita memperoleh
nilai harapan banyaknya proses berada dalam keadaan transien mana pun untuk
keadaan awal~$s_i$, yakni nilai harapan waktu yang diperlukan sebelum terserap.
Jadi, $t_i$ adalah jumlah entri pada baris ke-$i$ dari $\mat{N}$. Jika pernyataan
ini dituliskan dalam bentuk matriks, kita memperoleh teorema tersebut.
\end{theorem}

\subsection*{Peluang Penyerapan}\index{peluang penyerapan}
\begin{theorem}\label{thm 11.2.3}
Misalkan $b_{ij}$ adalah peluang bahwa suatu rantai penyerap akan terserap dalam
keadaan penyerap~$s_j$ jika dimulai dalam keadaan transien~$s_i$. Misalkan
$\mat{B}$ adalah matriks dengan entri $b_{ij}$. Maka $\mat{B}$ adalah matriks
berukuran $t$-kali-$r$, dan
$$
\mat{B} = \mat{N} \mat{R}\ ,
$$
di mana $\mat{N}$ adalah matriks fundamental dan $\mat{R}$ adalah matriks yang
muncul dalam bentuk kanonik.
\proof
Kita memiliki
\begin{eqnarray*}
\mat{B}_{ij} &=& \sum_n\sum_k q_{ik}^{(n)} r_{kj} \\
&=& \sum_k \sum_n q_{ik}^{(n)} r_{kj} \\
&=& \sum_k n_{ik}r_{kj} \\
&=& (\mat{N}\mat{R})_{ij}\ .
\end{eqnarray*}
Ini menyelesaikan bukti.
\end{theorem}
\par
Bukti lain diberikan dalam Latihan~\ref{exer 11.2.33}.

\begin{example}\label{exam 11.2.3}\index{contoh perjalanan acak
pemabuk}(Lanjutan Contoh~\ref{exam 11.2.2})
Dalam contoh perjalanan acak seorang pemabuk, kita memperoleh
$$
\vspace{7pt}\hbox{$\mat{N} = {}$} \bordermatrix{
  & 1   & 2 & 3   \cr
1 & 3/2 & 1 & 1/2 \cr
2 & 1   & 2 & 1   \cr
3 & 1/2 & 1 & 3/2 \cr}\ .
$$
Dengan demikian,
\vspace{7pt}
\begin{eqnarray*}
\mat{t} = \mat{N}\mat{c} &=& \pmatrix{
3/2 & 1 & 1/2 \cr
  1 & 2 & 1   \cr
1/2 & 1 & 3/2 \cr
} \pmatrix{
1 \cr
1 \cr
1 \cr
}
\\
&=& \pmatrix{
3 \cr
4 \cr
3 \cr
}\ .
\end{eqnarray*}
Jadi, jika proses dimulai dalam keadaan 1, 2, dan 3, nilai harapan waktu hingga
penyerapan berturut-turut adalah 3, 4, dan 3.
\par
Dari bentuk kanonik,
$$
\vspace{7pt}\hbox{$\mat{ R} = {}$} \bordermatrix{
  & 0   & 4   \cr
1 & 1/2 & 0   \cr
2 & 0   & 0   \cr
3 & 0   & 1/2 \cr}\ .
$$
Dengan demikian,
\begin{eqnarray*}
\mat{B} = \mat{N} \mat{R} &=& \pmatrix{
3/2 & 1 & 1/2 \cr
  1 & 2 & 1   \cr
1/2 & 1 & 3/2 \cr} \cdot \pmatrix{
1/2 & 0   \cr
  0 & 0   \cr
  0 & 1/2 \cr} \cr
\cr
\cr
&=& \bordermatrix{
  & 0 & 4 \cr
1 & 3/4 & 1/4 \cr
2 & 1/2 & 1/2 \cr
3 & 1/4 & 3/4 \cr}\ .
\end{eqnarray*}


\noindent Di sini baris pertama memberi tahu kita bahwa, jika dimulai dalam
keadaan $1$, peluang terserap dalam keadaan $0$ adalah~3/4 dan peluang terserap
dalam keadaan $4$ adalah 1/4.
\end{example}

\subsection*{Komputasi}
Karena kita dapat memperoleh ketiga besaran deskriptif ini dalam bentuk
matriks, sangat mudah untuk menulis program komputer yang menentukan
besaran-besaran tersebut bagi suatu matriks rantai penyerap tertentu.

Program {\bf AbsorbingChain}\index{AbsorbingChain (program)} menghitung
besaran-besaran deskriptif dasar suatu rantai Markov penyerap.

Kita telah menjalankan program {\bf AbsorbingChain} untuk contoh perjalanan acak
seorang pemabuk\index{contoh perjalanan acak pemabuk}
(Contoh~\ref{exam 11.2.1}) dengan 5~blok. Hasilnya adalah sebagai berikut:

$$
\mat{Q}  = \bordermatrix{
  & 1     & 2     & 3     & 4  \cr
1 & .00   & .50   & .00   & .00\cr
2 & .50   & .00   & .50   & .00\cr
3 & .00   & .50   & .00   & .50\cr 
4 & .00   & .00   & .50   & .00}\ ;
$$

$$
\mat{R}  = \bordermatrix{
  & 0     & 5     \cr
1 & .50   & .00   \cr
2 & .00   & .00   \cr
3 & .00   & .00   \cr 
4 & .00   & .50   }\ ;
$$

$$
\mat{N}  = \bordermatrix{
  & 1      & 2      & 3      &  4  \cr
1 & 1.60   & 1.20   & .80    & .40 \cr
2 & 1.20   & 2.40   & 1.60   & .80 \cr
3 & .80    & 1.60   & 2.40   & 1.20\cr 
4 & .40    & .80    & 1.20   & 1.60}\ ;
$$

$$
\mat{t}  = \bordermatrix{
 &  \cr
1  & 4.00 \cr
2  & 6.00 \cr
3  & 6.00 \cr
4  & 4.00}\ ;
$$

$$
\mat{B}  = \bordermatrix{
  & 0     & 5     \cr
1 & .80   & .20   \cr
2 & .60   & .40   \cr
3 & .40   & .60   \cr 
4 & .20   & .80   }\ .
$$
 
Perhatikan bahwa peluang mencapai bar sebelum mencapai rumah jika dimulai
di~$x$ adalah~$x/5$ (yakni, sebanding dengan jarak rumah dari titik awal).
(Lihat Latihan~\ref{exer 11.2.23}.)

\exercises
\begin{LJSItem}

\i\label{exer 11.2.1} Dalam Contoh~\ref{exam 11.1.2}, untuk nilai $a$~dan~$b$
manakah kita memperoleh rantai Markov penyerap?

\i\label{exer 11.2.2} Tunjukkan bahwa Contoh~\ref{exam 11.1.5} merupakan rantai
Markov penyerap.

\i\label{exer 11.2.3} Manakah di antara contoh-contoh genetika
(Contoh~\ref{exam
11.1.7},
~\ref{exam 11.1.8}, dan~\ref{exam 11.1.9}) yang bersifat menyerap?

\i\label{exer 11.2.4} Carilah matriks fundamental $\mat{N}$ untuk Contoh~\ref
{exam 11.1.8}.

\i\label{exer 11.2.5} Untuk Contoh~\ref{exam 11.1.9}, periksalah bahwa matriks
berikut adalah invers dari $\mat{I} - \mat{Q}$ sehingga merupakan matriks
fundamental~$\mat{N}$.
$$
\mat{N} = \pmatrix{
8/3 & 1/6 & 4/3 & 2/3 \cr
4/3 & 4/3 & 8/3 & 4/3 \cr
4/3 & 1/3 & 8/3 & 4/3 \cr
2/3 & 1/6 & 4/3 & 8/3 \cr}\ .
$$
Carilah $\mat{N} \mat{c}$ dan $\mat{N} \mat{R}$. Tafsirkan hasilnya.

\i\label{exer 11.2.6} Dalam contoh Negeri Oz (Contoh~\ref{exam 11.1.1}), ubahlah
matriks transisi dengan menjadikan R sebagai keadaan penyerap. Hasilnya adalah
$$\mat{P} = 
\bordermatrix{
  & \mbox{R} & \mbox{N} & \mbox{S} \cr
\mbox{R} & 1 & 0 & 0 \cr
\mbox{N} & 1/2 & 0 & 1/2 \cr
\mbox{S} & 1/4 & 1/4 & 1/2}\ .
$$
Carilah matriks fundamental~$\mat{N}$, serta $\mat{Nc}$ dan $\mat{NR}$.
Tafsirkan hasilnya.

\i\label{exer 11.2.7} Dalam Contoh~\ref{exam 11.1.6}, jadikan keadaan 0~dan~4
sebagai keadaan penyerap. Carilah matriks fundamental~$\mat{N}$, serta
$\mat{Nc}$ dan $\mat{NR}$, untuk rantai penyerap yang dihasilkan. Tafsirkan
hasilnya.

\i\label{exer 11.2.8} Dalam Contoh~\ref{exam 11.2.1} (perjalanan acak seorang
pemabuk)\index{contoh perjalanan acak pemabuk} pada bagian ini, anggaplah
peluang melangkah ke kanan adalah~2/3 dan peluang melangkah ke kiri adalah~1/3.
Carilah $\mat{N},~\mat{N}\mat{c}$, dan~$\mat{N}\mat{R}$. Bandingkan dengan hasil
Contoh~\ref{exam 11.2.3}.

\i\label{exer 11.2.9} Suatu proses bergerak pada bilangan bulat 1,~2, 3, 4,
dan~5. Proses dimulai di~1 dan, pada setiap langkah berikutnya, bergerak ke suatu
bilangan bulat yang lebih besar daripada posisinya saat itu, dengan peluang yang
sama untuk setiap bilangan bulat lebih besar yang tersisa. Keadaan lima adalah
keadaan penyerap. Carilah nilai harapan jumlah langkah untuk mencapai keadaan
lima.

\i\label{exer 11.2.10} Dengan menggunakan hasil Latihan~\ref{exer 11.2.9},
buatlah dugaan tentang bentuk matriks fundamental jika proses bergerak seperti
dalam latihan tersebut, kecuali bahwa sekarang proses bergerak pada bilangan
bulat dari 1~hingga~$n$. Ujilah dugaan Anda untuk beberapa nilai $n$ yang
berbeda. Dapatkah Anda menduga suatu taksiran bagi nilai harapan jumlah langkah
untuk mencapai keadaan $n$ ketika $n$ besar? (Lihat
Latihan~\ref{exer 11.2.10.5} untuk metode menentukan nilai harapan jumlah langkah
ini.)

\istar\label{exer 11.2.10.5} Misalkan $b_k$ menyatakan nilai harapan jumlah
langkah untuk mencapai $n$ dari $n-k$ dalam proses yang dijelaskan dalam
Latihan~\ref{exer 11.2.9}.
\begin{enumerate}
\item
Definisikan $b_0 = 0$. Tunjukkan bahwa untuk $k > 0$, kita memiliki
$$b_k = 1 + \frac 1k \bigl(b_{k-1} + b_{k-2} + \cdots + b_0\bigr)\ .$$
\item
Misalkan
$$f(x) = b_0 + b_1 x + b_2 x^2 + \cdots\ .$$
Dengan menggunakan rekurensi pada bagian (a), tunjukkan bahwa $f(x)$ memenuhi
persamaan diferensial
$$(1-x)^2 y' - (1-x) y - 1 = 0\ .$$
\item
Tunjukkan bahwa solusi umum persamaan diferensial pada bagian (b) adalah
$$y = \frac{-\log(1-x)}{1-x} + \frac c{1-x}\ ,$$
di mana $c$ adalah suatu konstanta.
\item
Gunakan bagian (c) untuk menunjukkan bahwa
$$b_k = 1 + \frac 12 + \frac 13 + \cdots + \frac 1k\ .$$
\end{enumerate}

\i\label{exer 11.2.11} Tiga tank berduel satu sama lain. Tank A memiliki
peluang~1/2 untuk menghancurkan tank yang ditembaknya, tank B memiliki
peluang~1/3, dan tank C memiliki peluang~1/6. Ketiga tank menembak secara
bersamaan, dan setiap tank menembak lawan terkuat yang belum dihancurkan.
Bentuklah rantai Markov dengan mengambil himpunan bagian dari himpunan tank
sebagai keadaan. Carilah $\mat{N},~\mat{N}\mat{c}$, dan~$\mat{N}\mat{R}$, lalu
tafsirkan hasil Anda. \emx {Petunjuk}: Ambillah ABC, AC, BC, A, B, C, dan tidak
ada sebagai keadaan, yang menunjukkan tank-tank yang dapat bertahan jika
dimulai dalam keadaan ABC. Anda dapat mengabaikan AB karena keadaan ini tidak
dapat dicapai dari ABC.

\i\label{exer 11.2.12} Smith berada di penjara dan memiliki 3~dolar; ia dapat
keluar dengan jaminan jika memiliki 8~dolar. Seorang penjaga bersedia membuat
serangkaian taruhan dengannya. Jika Smith mempertaruhkan $A$~dolar, ia
memenangkan $A$~dolar dengan peluang~.4 dan kehilangan $A$~dolar dengan
peluang~.6. Carilah peluang bahwa ia memperoleh 8~dolar sebelum kehilangan
seluruh uangnya jika
\begin{enumerate}
\item ia mempertaruhkan 1~dolar setiap kali (strategi hati-hati).

\item setiap kali ia mempertaruhkan sebanyak mungkin, tetapi tidak lebih dari
yang diperlukan agar kekayaannya mencapai 8~dolar (strategi berani).

\item Strategi manakah yang memberi Smith peluang lebih baik untuk keluar dari
penjara?
\end{enumerate}

\i\label{exer 11.2.13} Dalam situasi Latihan~\ref{exer 11.2.12}, tinjaulah
strategi yang, untuk $i < 4$, membuat Smith mempertaruhkan $\min(i,4 - i)$, dan
untuk $i \geq 4$ membuatnya bertaruh menurut strategi berani, dengan $i$ sebagai
kekayaannya saat itu. Carilah peluang bahwa ia keluar dari penjara dengan
strategi ini. Bagaimanakah peluang ini dibandingkan dengan peluang pada strategi
berani?

\i\label{exer 11.2.14} Tinjaulah permainan tenis\index{tenis} ketika keadaan
\emx {deuce} tercapai. Jika seorang pemain memenangkan poin berikutnya, ia
memiliki \emx {keuntungan.} Pada poin setelahnya, ia memenangkan gim atau gim
kembali ke keadaan \emx {deuce.} Anggaplah bahwa pada setiap poin pemain A
memiliki peluang .6 untuk memenangkan poin tersebut dan pemain B memiliki
peluang .4.
\begin{enumerate}
\item Bentuklah rantai Markov dengan keadaan~1: A menang; 2: B menang; 3:
keuntungan A; 4: deuce; 5: keuntungan B.

\item Carilah peluang-peluang penyerapannya.

\item Pada keadaan deuce, carilah nilai harapan durasi gim dan peluang bahwa B
akan menang.
\end{enumerate}

\medbreak
Latihan \ref{exer 11.2.15}~dan~\ref{exer 11.2.16} membahas pewarisan buta
warna,\index{buta warna} yaitu sifat yang terpaut kelamin. Terdapat sepasang gen,
g~dan~G; gen pertama cenderung menghasilkan buta warna, sedangkan gen kedua
menghasilkan penglihatan normal. Gen G bersifat dominan. Namun, seorang pria
hanya memiliki satu gen, dan jika gen tersebut adalah~g, ia buta warna. Seorang
pria mewarisi salah satu dari dua gen ibunya, sedangkan seorang wanita mewarisi
satu gen dari setiap orang tua. Jadi, seorang pria dapat bertipe G~atau~g,
sedangkan seorang wanita dapat bertipe GG, Gg, atau gg. Kita akan mempelajari
proses perkawinan sekerabat yang serupa dengan proses pada
Contoh~\ref{exam 11.1.9} dengan membentuk rantai Markov.

\i\label{exer 11.2.15} Daftarkan keadaan-keadaan rantai tersebut. \emx
{Petunjuk}: Terdapat enam keadaan. Hitunglah peluang-peluang transisinya. Carilah
matriks fundamental $\mat{N}$, $\mat{N}\mat{c}$, dan $\mat{N}\mat{R}$.

\i\label{exer 11.2.16} Tunjukkan bahwa baik dalam Contoh~\ref{exam 11.1.9}
maupun dalam contoh yang baru diberikan, peluang terserap dalam keadaan yang
memiliki gen bertipe tertentu sama dengan proporsi gen bertipe tersebut dalam
keadaan awal. Tunjukkan bahwa hal ini dapat dijelaskan oleh fakta bahwa permainan
yang kekayaan Anda di dalamnya sama dengan jumlah gen bertipe tertentu dalam
keadaan rantai Markov merupakan permainan adil.\footnote{H. Gonshor, ``An Application of
Random Walk to a Problem in Population Genetics," \emx {American Math Monthly,}
vol.~94 (1987), pp.~668--671}\index{GONSHOR, H.}

\i\label{exer 11.2.17} Anggaplah seorang mahasiswa yang berkuliah di suatu
sekolah kedokteran empat tahun di New England bagian utara setiap tahun memiliki
peluang~$q$ untuk gagal dan dikeluarkan, peluang~$r$ untuk harus mengulang tahun
tersebut, serta peluang~$p$ untuk melanjutkan ke tahun berikutnya (pada tahun
keempat, melanjutkan berarti lulus).
\begin{enumerate}
\item Bentuklah matriks transisi bagi proses ini dengan mengambil F,~1, 2, 3,
4, dan~G sebagai keadaan, dengan F menyatakan gagal dan dikeluarkan, G menyatakan
lulus, serta keadaan lainnya menyatakan tahun studi.

\item Untuk kasus $q = .1$, $r = .2$, dan $p = .7$, carilah nilai harapan waktu
seorang mahasiswa baru berada pada tahun kedua. Berapa lama nilai harapan
mahasiswa tersebut berada di sekolah kedokteran?

\item Carilah peluang bahwa mahasiswa baru tersebut akan lulus.
\end{enumerate}

\i\label{exer 11.2.18} (E. Brown\footnote{Komunikasi pribadi.})\index{BROWN,
E.} 
Mary dan John memainkan permainan berikut: mereka memiliki satu set berisi tiga
kartu yang ditandai dengan angka 1,~2, dan~3, serta sebuah pemutar dengan angka
1,~2, dan~3. Permainan dimulai dengan membagikan kartu sehingga pembagi mendapat
satu kartu dan pemain lain mendapat dua kartu. Satu langkah permainan terdiri
atas satu putaran pemutar. Orang yang memegang kartu dengan angka yang muncul
pada pemutar menyerahkan kartu itu kepada orang lain. Permainan berakhir ketika
seseorang memiliki semua kartu.
\begin{enumerate}

\item Susunlah matriks transisi bagi rantai Markov penyerap ini, dengan keadaan
yang bersesuaian dengan jumlah kartu yang dimiliki Mary.

\item Carilah matriks fundamentalnya.

\item Secara rata-rata, berapa langkah permainan berlangsung?

\item Jika Mary membagikan kartu, berapakah peluang John memenangkan permainan?
\end{enumerate}

\i\label{exer 11.2.19} Anggaplah suatu percobaan memiliki $m$ hasil yang sama
mungkin. Tunjukkan bahwa nilai harapan jumlah percobaan saling bebas sebelum
kemunculan pertama $k$ hasil sejenis berturut-turut adalah
$(m^k - 1)/(m - 1)$. \emx {Petunjuk}: Bentuklah rantai Markov penyerap dengan
keadaan 1,~2, \ldots,~$k$, dengan keadaan~$i$ menyatakan panjang rentetan saat
ini. Nilai harapan waktu hingga rentetan sepanjang~$k$ adalah 1~lebih besar
daripada nilai harapan waktu hingga penyerapan bagi rantai yang dimulai dalam
keadaan~1. Telah ditemukan bahwa, dalam perluasan desimal pi, mulai dari digit
ke-24{,}658{,}601 terdapat rentetan sembilan angka 7. Menurut hasil Anda,
berapakah nilai harapan jumlah digit yang diperlukan untuk menemukan rentetan
seperti itu jika digit dihasilkan secara acak?

\i\label{exer 11.2.20} (Roberts\footnote{F. Roberts, \emx {Discrete
Mathematical 
Models} (Englewood Cliffs, NJ: Prentice Hall, 1976).})\index{ROBERTS, F.} A
kota dibagi menjadi 3 wilayah, yaitu 1,~2, dan~3. Diperkirakan bahwa setiap hari
ketiga wilayah ini masing-masing mengeluarkan polusi sebanyak $u_1$,~$u_2$,
dan~$u_3$. Sebagian $q_{ij}$ dari polusi wilayah~$i$ berakhir di wilayah~$j$
pada hari berikutnya. Sebagian $q_i = 1 - \sum_j q_{ij} > 0$ masuk ke atmosfer
dan lolos. Misalkan $w_i^{(n)}$ adalah jumlah polusi di wilayah~$i$ setelah
$n$~hari.
\begin{enumerate}

\item Tunjukkan bahwa $\mat{w}^{(n)} = \mat{u} + \mat{u} \mat{Q} +\cdots +
\mat{u}\mat{Q}^{n - 1}$.

\item Tunjukkan bahwa $\mat{w}^{(n)} \to \mat{w}$, dan tunjukkan cara menghitung
\mat{w}~dari~\mat{u}.
 
\item Pemerintah ingin membatasi tingkat polusi pada suatu tingkat yang
ditetapkan dengan menentukan~$\mat{w}.$ Tunjukkan cara menentukan tingkat polusi
$\mat{u}$ yang akan menghasilkan nilai limit~$\mat{w}$ yang ditetapkan.
\end{enumerate}

\i\label{exer 11.2.21} Dalam model ekonomi Leontief,\index{LEONTIEF, W.
W.}\footnote{W.~W. Leontief, 
\emx {Input-Output Economics} (Oxford: Oxford University Press, 1966).} terdapat
$n$ industri, yaitu 1,~2, \ldots,~$n$. Industri ke-$i$ memerlukan barang senilai
$0 \leq q_{ij} \leq 1$ dolar dari perusahaan~$j$ untuk menghasilkan barang
senilai 1~dolar. Permintaan eksternal terhadap industri-industri tersebut,
dalam nilai dolar, diberikan oleh vektor $\mat{d} = (d_1,d_2,\ldots,d_n)$.
Misalkan $\mat{Q}$ adalah matriks dengan entri~$q_{ij}$.
\begin{enumerate}

\item Tunjukkan bahwa jika industri-industri tersebut menghasilkan jumlah total
yang diberikan oleh vektor $\mat{x} = (x_1,x_2,\ldots,x_n)$, jumlah barang dari
setiap jenis yang diperlukan industri-industri tersebut sekadar untuk memenuhi
permintaan internal diberikan oleh vektor~$\mat{x} \mat{Q}$.

\item Tunjukkan bahwa untuk memenuhi permintaan eksternal~$\mat{d}$ maupun
permintaan internal, industri-industri tersebut harus menghasilkan jumlah total
yang diberikan oleh suatu vektor $\mat{x} = (x_1,x_2,\ldots,x_n)$ yang memenuhi
persamaan $\mat{x} = \mat{x} \mat{Q} + \mat{d}$.

\item Tunjukkan bahwa jika $\mat{Q}$ adalah matriks $\mat{Q}$ bagi suatu rantai
Markov penyerap, setiap permintaan eksternal~$\mat{d}$ dapat dipenuhi.

\item Anggaplah jumlah setiap baris~$\mat{Q}$ kurang dari atau sama dengan~1.
Berikan tafsiran ekonomi bagi kondisi ini. Bentuklah rantai Markov dengan
mengambil industri-industri sebagai keadaan dan~$q_{ij}$ sebagai peluang
transisi. Tambahkan satu keadaan penyerap~0. Definisikan
$$
q_{i0} = 1 - \sum_j q_{ij}\ .
$$
Tunjukkan bahwa rantai ini bersifat menyerap jika setiap perusahaan menghasilkan
laba atau pada akhirnya bergantung pada perusahaan yang menghasilkan laba.

\item Definisikan $\mat{x} \mat{c}$ sebagai produk nasional bruto. Carilah
ungkapan produk nasional bruto dalam vektor permintaan $\mat{d}$ dan vektor
$\mat{t}$ yang memberikan nilai harapan waktu hingga penyerapan.
\end{enumerate}

\i\label{exer 11.2.22} Seorang penjudi memainkan permainan yang pada setiap
putaran memberinya kemenangan satu dolar dengan peluang~$p$ dan kerugian satu
dolar dengan peluang $q = 1 - p$. \index{kebangkrutan penjudi}\emx {Masalah
Kebangkrutan Penjudi} adalah masalah mencari peluang~$w_x$ untuk memperoleh
jumlah~$T$ sebelum kehilangan semuanya jika dimulai dengan keadaan~$x$.
Tunjukkan bahwa masalah ini dapat dipandang sebagai rantai Markov penyerap
dengan keadaan 0,~1, 2, \ldots,~$T$, dengan 0~dan~$T$ sebagai keadaan penyerap.
Misalkan seorang penjudi memiliki peluang $p = .48$ untuk menang pada setiap
putaran. Misalkan pula penjudi tersebut mulai dengan 50~dolar dan $T = 100$
dolar. Simulasikan permainan ini 100 kali dan perhatikan seberapa sering penjudi
tersebut bangkrut. Hasil ini menaksir $w_{50}$.

\i\label{exer 11.2.23} Tunjukkan bahwa $w_x$ pada
Latihan~\ref{exer 11.2.22} memenuhi kondisi-kondisi berikut:
\begin{enumerate}

\item $w_x = pw_{x + 1} + qw_{x - 1}$ untuk $x = 1$,~2, \ldots,\ $T - 1$.

\item $w_0 = 0$.

\item $w_T = 1$.
\end{enumerate}

\noindent Tunjukkan bahwa kondisi-kondisi ini menentukan $w_x$. Tunjukkan bahwa
jika $p = q = 1/2$, maka
$$
w_x = \frac xT
$$
memenuhi (a), (b), dan (c), sehingga merupakan solusinya. Jika $p \ne q$,
tunjukkan bahwa
$$
w_x = \frac{(q/p)^x - 1}{(q/p)^T - 1}
$$
memenuhi kondisi-kondisi tersebut, sehingga memberikan peluang penjudi menang.

\i\label{exer 11.2.24} Tulislah program untuk menghitung peluang~$w_x$ pada
Latihan~\ref{exer 11.2.23} bagi nilai $x$,~$p$, dan~$T$ yang diberikan. Pelajari
peluang bahwa penjudi akan membangkrutkan bank dalam permainan yang hanya
sedikit tidak menguntungkan baginya, misalnya $p = .49$, jika bank memiliki uang
jauh lebih banyak daripada penjudi.

\istar\label{exer 11.2.25} Kita meninjau dua contoh perjalanan acak seorang
pemabuk\index{contoh perjalanan acak pemabuk} yang bersesuaian dengan kasus $n =
4$ dan $n = 5$ blok (lihat Contoh~\ref{exam 11.2.1}). Periksalah bahwa dalam
kedua contoh tersebut, nilai harapan waktu hingga penyerapan jika dimulai
di~$x$ sama dengan $x(n - x)$. Cobalah membuktikan bahwa hal ini berlaku secara
umum. \emx {Petunjuk}: Tunjukkan bahwa jika $f(x)$ adalah nilai harapan waktu
hingga penyerapan, maka $f(0) = f(n) = 0$ dan
$$f(x) = (1/2)f(x - 1) + (1/2)f(x + 1) + 1$$
untuk $0 < x < n$. Tunjukkan bahwa jika $f_1(x)$ dan $f_2(x)$ adalah dua
solusi, selisihnya $g(x)$ merupakan solusi persamaan
$$g(x) = (1/2)g(x - 1) + (1/2)g(x + 1)\ .$$
Selain itu, $g(0) = g(n) = 0$. Tunjukkan bahwa $g(x)$ tidak mungkin memiliki
maksimum ketat ataupun minimum ketat di titik $i$, dengan $1 \le i \le n-1$.
Gunakan hal ini untuk menunjukkan bahwa $g(i) = 0$ bagi semua i. Ini menunjukkan
bahwa paling banyak ada satu solusi. Kemudian, periksalah bahwa fungsi
$f(x) = x(n-x)$ merupakan suatu solusi.

\i\label{exer 11.2.29} Tinjaulah suatu rantai Markov penyerap dengan ruang
keadaan~$S$. Misalkan $f$ adalah fungsi yang didefinisikan pada~$S$ dengan sifat
$$
f(i) = \sum_{j \in S} p_{ij} f(j)\ ,
$$
atau, dalam bentuk vektor,
$$
\mat{f} = \mat{Pf}\ .
$$
Maka $f$ disebut \emx {fungsi harmonik} bagi $\mat{P}$.\index{fungsi harmonik}
Jika Anda membayangkan suatu permainan yang kekayaan Anda di dalamnya adalah
$f(i)$ ketika berada dalam keadaan~$i$, kondisi harmonik berarti bahwa permainan
tersebut \emx {adil}, dalam arti bahwa nilai harapan kekayaan Anda setelah satu
langkah sama dengan sebelum langkah tersebut.
\begin{enumerate}

\item Tunjukkan bahwa untuk $f$ yang harmonik,
$$
\mat{f} = \mat{P}^n{\mat{f}}
$$
untuk semua $n$.

\item Dengan menggunakan (a), tunjukkan bahwa untuk $f$ yang harmonik,
$$
\mat{f} = \mat{P}^\infty \mat{f}\ ,
$$
di mana
$$
 \mat{P}^\infty = \lim_{n \to \infty} \mat{P}^n =
\pmatrix{$$\begin{tabular}{l|r} \mat{0} & \mat{B} \\ \hline 
                                \mat{0} & \mat{I} \\ \end{tabular}$$ \cr}\ .
$$

\item Dengan menggunakan (b), buktikan bahwa jika Anda mulai dalam keadaan
transien~$i$, nilai harapan kekayaan akhir Anda
$$
\sum_k b_{ik} f(k)
$$
sama dengan kekayaan awal Anda~$f(i)$. Dengan kata lain, permainan adil pada
ruang keadaan berhingga tetap adil hingga akhir. (Permainan adil secara umum
disebut \emx {martingal.}\index{martingal} Permainan adil pada ruang keadaan tak
berhingga tidak harus tetap adil jika jumlah putaran tak terbatas diizinkan.
Sebagai contoh, tinjaulah permainan Kepala atau Ekor (lihat
Contoh~\ref{exam 1.3}). Misalkan Peter mulai dengan 1~sen dan bermain hingga
memiliki~2. Maka Peter pasti berakhir dengan keuntungan 1~sen.)
\end{enumerate} 

\i\label{exer 11.2.26} Sebuah koin dilempar berulang kali. Kita ingin mencari
nilai harapan jumlah lemparan hingga suatu pola tertentu, misalnya B = HTH,
muncul untuk pertama kalinya. Sebagai contoh, jika hasil lemparannya adalah
HHTTHTH, kita mengatakan bahwa pola B muncul untuk pertama kalinya setelah
7~lemparan. Misalkan $T^B$ adalah waktu untuk memperoleh pola B untuk pertama
kalinya. Li\footnote{S-Y.~R. Li, ``A Martingale Approach to the Study of
Occurrence of Sequence Patterns in Repeated Experiments,'' \emx {Annals of
Probability,} vol.~8 (1980), pp.~1171--1176.} memberikan metode berikut untuk
menentukan $E(T^B)$.

Kita berada di sebuah kasino dan, sebelum setiap lemparan koin, seorang penjudi
masuk, membayar 1~dolar untuk bermain, lalu bertaruh bahwa pola B = HTH akan
muncul pada tiga lemparan berikutnya. Jika H muncul, ia memenangkan 2~dolar dan
mempertaruhkan jumlah tersebut bahwa hasil berikutnya adalah~T. Jika menang, ia
memenangkan 4~dolar dan mempertaruhkan jumlah tersebut bahwa H akan muncul
berikutnya. Jika menang, ia memenangkan 8~dolar dan pola tersebut telah muncul.
Jika kalah pada tahap mana pun, ia pergi tanpa membawa kemenangan.

Misalkan A dan B adalah dua pola. Misalkan AB adalah jumlah yang dimenangkan
para penjudi yang datang ketika pola A muncul dan bertaruh bahwa B akan muncul.
Sebagai contoh, jika A = HT dan B = HTH, maka AB = 2 + 4 = 6 karena penjudi
pertama bertaruh pada H dan memenangkan 2 dolar, kemudian bertaruh pada T dan
memenangkan 4 dolar lagi. Penjudi kedua bertaruh pada H dan kalah. Jika A = HH
dan B = HTH, maka AB = 2 karena penjudi pertama bertaruh pada H dan menang,
tetapi kemudian bertaruh pada T dan kalah, sedangkan penjudi kedua bertaruh pada
H dan menang. Jika A = B = HTH, maka AB = BB = 8 + 2 = 10.
\par
Untuk setiap penjudi yang masuk, kasino menerima 1~dolar. Jadi, kasino menerima
$T^B$~dolar. Berapa banyak yang dibayarkannya? Para penjudi yang pulang membawa
uang hanyalah mereka yang datang selama pola B muncul, dan mereka memenangkan
jumlah BB. Namun, karena semua taruhan yang dibuat benar-benar adil, secara
intuitif nilai harapan jumlah yang diterima kasino semestinya sama dengan nilai
harapan jumlah yang dibayarkannya. Artinya, $E(T^B)$ = BB.
\par
Karena kita telah melihat bahwa untuk B = HTH berlaku BB = 10, nilai harapan
waktu untuk mencapai pola HTH pertama kali adalah 10. Jika kita mencoba
memperoleh pola B = HHH, maka BB $= 8 + 4 + 2 = 14$ karena dalam kasus ini tiga
penjudi terakhir semuanya dibayar. Jadi, nilai harapan waktu untuk memperoleh
pola HHH adalah 14. Untuk membenarkan argumen ini, Li menggunakan suatu teorema
dari teori martingal (permainan adil).
\par
Kita dapat memperoleh nilai-nilai harapan ini dengan meninjau rantai Markov yang
keadaannya merupakan ruas-ruas awal yang mungkin dari urutan HTH; keadaan
tersebut adalah HTH, HT, H, dan $\emptyset$, dengan $\emptyset$ sebagai himpunan
kosong. Untuk contoh ini, matriks transisinya adalah
$$
\bordermatrix{
&\mbox{HTH} & \mbox{HT} & \mbox{H} & \emptyset \cr
\mbox{HTH}  & 1 & 0 & 0 & 0 \cr
\mbox{HT}   & .5 & 0 & 0 & .5 \cr
\mbox{H}    & 0 & .5 & .5 & 0 \cr
\emptyset   & 0 & 0 & .5 & .5 }\ ,
$$
dan jika B = HTH, $E(T^B)$ adalah nilai harapan waktu hingga penyerapan bagi
rantai ini jika dimulai dalam keadaan~$\emptyset$.
\par
Dengan menggunakan rantai Markov yang terkait, tunjukkan bahwa nilai
$E(T^B)$ = 10 dan $E(T^B)$ = 14 benar bagi nilai harapan waktu untuk mencapai
pola HTH dan HHH, berturut-turut.

\i\label{exer 11.2.27} Kita dapat menggunakan tafsiran perjudian yang diberikan
dalam Latihan~\ref{exer 11.2.26} untuk mencari nilai harapan jumlah lemparan yang
diperlukan untuk mencapai pola B jika kita mulai dengan pola A. Agar masalahnya
bermakna, kita menganggap pola A tidak memuat pola B sebagai subpola. Misalkan
$E_A(T^B)$ adalah nilai harapan waktu untuk mencapai pola B jika dimulai dengan
pola A. Kita menggunakan skema perjudian kita dan menganggap bahwa k lemparan
koin pertama menghasilkan pola A. Selama waktu ini, para penjudi memenangkan
jumlah AB. Jumlah total yang akan dimenangkan para penjudi ketika pola B muncul
adalah BB. Jadi, jumlah yang dimenangkan para penjudi setelah pola A muncul
adalah BB - AB. Sekali lagi, menurut argumen permainan adil,
$E_A(T^B)$ = BB-AB.
\par
Sebagai contoh, misalkan kita mulai dengan pola A = HT dan mencoba memperoleh
pola B = HTH. Dalam Latihan~\ref{exer 11.2.26}, kita melihat bahwa AB = 4 dan
BB = 10, sehingga $E_A(T
^B)$ = BB-AB= 6.
\par
Periksalah bahwa tafsiran perjudian ini menghasilkan jawaban yang benar untuk
semua keadaan awal dalam contoh-contoh yang Anda kerjakan pada
Latihan~\ref{exer 11.2.26}.

\i\label{exer 11.2.28} Berikut adalah metode elegan dari Guibas dan
Odlyzko\footnote{L.~J.
Guibas and A.~M. Odlyzko, ``String Overlaps, Pattern Matching, and
Non-transitive Games,"
\emx {Journal of Combinatorial Theory,} Series~A, vol.~30 (1981),
pp.~183--208.} untuk memperoleh nilai harapan waktu hingga suatu pola, misalnya
HTH, tercapai untuk pertama kalinya. Misalkan $f(n)$ adalah jumlah urutan
sepanjang~$n$ yang tidak memuat pola HTH. Misalkan $f_p(n)$ adalah jumlah urutan
yang memuat pola tersebut untuk pertama kalinya setelah $n$~lemparan. Pada
setiap elemen dari~$f(n)$, tambahkan pola HTH. Kemudian, bagi urutan yang
dihasilkan menjadi tiga himpunan bagian: himpunan tempat HTH muncul untuk pertama
kalinya pada waktu $n + 1$ (untuk ini, urutan semula harus berakhir dengan HT);
himpunan tempat HTH muncul untuk pertama kalinya pada waktu $n + 2$ (tidak dapat
terjadi untuk pola ini); dan himpunan tempat urutan HTH muncul untuk pertama
kalinya pada waktu $n + 3$ (urutan semula berakhir dengan apa pun selain HT).
Dengan melakukan hal ini, kita memperoleh
$$
f(n) = f_p(n + 1) + f_p(n + 3)\ .
$$
Dengan demikian,
$$
\frac{f(n)}{2^n} = \frac{2f_p(n + 1)}{2^{n + 1}} + \frac{2^3f_p(n + 3)}{2^{n +
3}}\ .
$$
Jika $T$ adalah waktu pola tersebut muncul untuk pertama kalinya, kesamaan ini
menyatakan bahwa
$$
P(T > n) = 2P(T = n + 1) + 8P(T = n + 3)\ .
$$
Tunjukkan bahwa dengan menjumlahkan kesamaan ini atas semua~$n$, Anda memperoleh
$$
\sum_{n = 0}^\infty P(T > n) = 2 + 8 = 10\ .
$$
Tunjukkan bahwa untuk sebarang peubah acak bernilai bilangan bulat,
$$
E(T) = \sum_{n = 0}^\infty P(T > n)\ ,
$$
dan simpulkan bahwa $E(T) = 10$. Perhatikan bahwa metode bukti ini memperjelas
bahwa secara umum $E(T)$ sama dengan nilai harapan jumlah yang dibayarkan kasino,
dan metode ini menghindari teorema sistem martingal yang digunakan Li.

\i\label{exer 11.2.30} Dalam Contoh~\ref{exam 11.1.9}, definisikan $f(i)$ sebagai
proporsi gen G dalam keadaan~$i$. Tunjukkan bahwa $f$ adalah fungsi harmonik
(lihat Latihan~\ref{exer 11.2.29}). Mengapa hal ini menunjukkan bahwa peluang
terserap dalam keadaan $(\mbox{GG},\mbox{GG})$ sama dengan proporsi gen G dalam
keadaan awal? (Lihat Latihan~\ref{exer 11.2.16}.)

\i\label{exer 11.2.31} Tunjukkan bahwa model batu loncatan
(Contoh~\ref{exam 11.1.10}) merupakan rantai Markov penyerap. Anggaplah Anda
memainkan permainan dengan persegi merah dan hijau, dan kekayaan Anda setiap
saat sama dengan proporsi persegi merah pada saat itu. Berikan argumen yang
menunjukkan bahwa permainan ini adil, dalam arti bahwa nilai harapan kemenangan
Anda setelah setiap langkah sama persis dengan sebelum langkah tersebut.\emx
{Petunjuk}: Tunjukkan bahwa untuk setiap hasil yang mungkin yang membuat
kekayaan Anda berkurang satu, terdapat hasil lain dengan peluang tepat sama yang
membuatnya bertambah satu.
\par
Gunakan fakta ini dan hasil Latihan~\ref{exer 11.2.29} untuk menunjukkan bahwa
peluang suatu warna tertentu akhirnya mendominasi sama dengan proporsi persegi
yang pada awalnya berwarna tersebut.

\i\label{exer 11.2.32} Tinjaulah pejalan acak yang bergerak pada bilangan bulat
0,~1, \ldots,~$N$, dengan bergerak satu langkah ke kanan dengan peluang~$p$ dan
satu langkah ke kiri dengan peluang $q = 1 - p$. Jika pejalan mencapai 0~atau~$N$,
ia tetap di sana. (Ini adalah masalah Kebangkrutan Penjudi pada
Latihan~\ref{exer 11.2.22}.) Jika $p = q$, tunjukkan bahwa fungsi
$$
f(i) = i
$$
merupakan fungsi harmonik (lihat Latihan~\ref{exer 11.2.29}), dan jika $p \ne q$,
maka
$$
f(i) = \biggl(\frac  {q}{p}\biggr)^i
$$
merupakan fungsi harmonik. Gunakan hal ini dan hasil
Latihan~\ref{exer 11.2.29} untuk menunjukkan bahwa peluang $b_{iN}$ untuk
terserap dalam keadaan~$N$ jika dimulai dalam keadaan~$i$ adalah
$$
 b_{iN} = \left \{ \matrix{
                                   \frac iN, &\mbox{jika}\,\, p = q, \cr
          \frac{({q \over p})^i - 1}{({q \over p})^{N} - 1}, & 
\mbox{jika}\,\,p \ne q.\cr}\right.
$$
Untuk penurunan alternatif hasil-hasil ini, lihat Latihan~\ref{exer 11.2.23}.

\i\label{exer 11.2.33} Lengkapilah bukti alternatif berikut untuk
Teorema~\ref{thm 11.2.3}. Misalkan $s_i$ adalah keadaan transien dan $s_j$
adalah keadaan penyerap. Jika kita menghitung $b_{ij}$ berdasarkan kemungkinan
hasil langkah pertama, kita memperoleh persamaan
$$
b_{ij} = p_{ij} + \sum_k p_{ik} b_{kj}\ ,
$$
di mana penjumlahan dilakukan atas semua keadaan transien~$s_k$. Tuliskan
persamaan ini dalam bentuk matriks dan turunkan pernyataan berikut darinya:
$$\mat{B} = \mat{N}\mat{R}\ .$$

\i\label{exer 11.2.34} Dalam rolet Monte Carlo\index{rolet} (lihat
Contoh~\ref{exam 6.1.5}), pada pilihan (c), terdapat enam keadaan ($S$, $W$,
$L$, $E$, $P_1$, dan $P_2$). Lihat Gambar~\ref{fig 6.1.5}, yang memuat pohon
untuk pilihan ini. Bentuklah rantai Markov bagi pilihan ini dan gunakan program
{\bf AbsorbingChain} untuk mencari peluang bahwa Anda menang, kalah, atau impas
pada taruhan 1 franc pada merah. Dengan menggunakan peluang-peluang ini, carilah
nilai harapan kemenangan untuk taruhan tersebut. Untuk pembahasan yang lebih
umum tentang penerapan rantai Markov pada rolet, lihat artikel H. Sagan yang
dirujuk dalam Contoh~\ref{exam 6.7}.

\i\label{exer 11.5.26} Selanjutnya, kita meninjau permainan yang disebut
\emx {Penney-ante} oleh penciptanya, W. Penney.\footnote{W. Penney, ``Problem: Penney-Ante," \emx {Journal
of
Recreational Math,} vol.~2 (1969), p.~241.}\index{PENNEY, W.} Terdapat dua
pemain; pemain pertama memilih pola A yang terdiri atas H dan T, lalu pemain
kedua, setelah mengetahui pilihan pemain pertama, memilih pola B yang berbeda.
Kita menganggap tidak satu pun pola merupakan subpola dari pola lainnya. Sebuah
koin dilempar berulang kali, dan pemain yang polanya muncul lebih dahulu menjadi
pemenang. Untuk menganalisis permainan ini, kita perlu mencari peluang $p_A$
bahwa pola A muncul sebelum pola B dan peluang $p_B = 1- p_A$ bahwa pola B
muncul sebelum pola A. Untuk menentukan peluang-peluang ini, kita menggunakan
hasil Latihan~\ref{exer 11.2.26} dan \ref{exer 11.2.27}. Dalam latihan tersebut,
Anda diminta menunjukkan bahwa nilai harapan waktu untuk pertama kali mencapai
pola B adalah
$$
E(T^B) = BB\ ,
$$
dan, jika dimulai dengan pola A, nilai harapan waktu untuk mencapai pola B adalah
$$
E_A(T^B) = BB - AB\ .
$$

\begin{enumerate}
\item Tunjukkan bahwa rasio peluang pemain pertama menang diberikan oleh rumus
John Conway\footnote{M. Gardner, ``Mathematical Games," \emx {Scientific American,}
 vol.~10 (1974), pp. 120--125.}\index{CONWAY, J.}:
$$
\frac{p_A}{1-p_A} = \frac{p_A}{p_B} = \frac{BB - BA}{AA - AB}\ .
$$
\emx {Petunjuk}: Jelaskan mengapa
$$
E(T^B) = E(T^{A\ {\rm atau}\ B}) + p_A E_A(T^B)
$$

dan dengan demikian
$$
BB = E(T^{A\ {\rm atau}\ B}) + p_A (BB - AB)\ .
$$
Pertukarkan A dan B untuk memperoleh persamaan serupa yang melibatkan $p_B$.
Terakhir, perhatikan bahwa
$$
p_A + p_B = 1\ .$$
Gunakan persamaan-persamaan ini untuk mencari $p_A$ dan $p_B$.

\item Anggaplah kedua pemain memilih pola dengan panjang k yang sama. Tunjukkan
bahwa jika $k = 2$, permainan ini adil, tetapi jika $k = 3$, pemain kedua
memiliki keuntungan apa pun pilihan pemain pertama. (Telah ditunjukkan bahwa,
untuk $k \geq 3$, jika pemain pertama memilih $a_1$,~$a_2$, \ldots,~$a_k$,
strategi optimal pemain kedua berbentuk $b$,~$a_1$, \ldots,~$a_{k - 1}$, dengan
$b$ sebagai pilihan yang lebih baik di antara H~atau~T.\footnote{Guibas and
Odlyzko, op. cit.})
\end{enumerate}

\end{LJSItem}

\section{Rantai Markov Ergodik}\label{sec 11.3}
 
Jenis penting kedua dari rantai Markov yang akan kita pelajari secara terperinci
adalah rantai Markov \emx{ergodik}, yang didefinisikan sebagai berikut.

\begin{definition}\label{defn 11.3.6.5}
Suatu rantai Markov disebut rantai \emx {ergodik} jika setiap keadaan dapat
dicapai dari setiap keadaan lain (tidak harus dalam satu langkah).
\index{rantai Markov ergodik}
\index{rantai Markov!ergodik}
\end{definition}
\par
Dalam banyak buku, rantai Markov ergodik disebut \emx {tak
tereduksi}.\index{rantai Markov tak tereduksi}\index{rantai Markov!tak tereduksi}

\begin{definition}\label{defn 11.3.7}
Suatu rantai Markov disebut rantai \emx {reguler} jika semua elemen suatu
pangkat matriks transisinya positif.\index{rantai Markov reguler}
\index{rantai Markov!reguler}
\end{definition}
\par
Dengan kata lain, terdapat suatu $n$ sedemikian rupa sehingga setiap keadaan
dapat dicapai dari setiap keadaan lain dalam tepat $n$ langkah. Dari definisi ini
jelas bahwa setiap rantai reguler bersifat ergodik. Sebaliknya, rantai ergodik
tidak harus reguler, sebagaimana ditunjukkan oleh contoh-contoh berikut.

\begin{example}
Misalkan matriks transisi suatu rantai Markov didefinisikan oleh
$$\mat{P} = \bordermatrix{
   &  1  &  2  \cr
1  &  0  &  1  \cr
2  &  1  &  0}\ .
$$
Jelas bahwa setiap keadaan dapat dicapai dari setiap keadaan lain, sehingga
rantai ini ergodik. Namun, jika $n$ ganjil, keadaan 0 tidak dapat dicapai dari
keadaan 0 dalam $n$ langkah, dan jika $n$ genap, keadaan 1 tidak dapat dicapai
dari keadaan 0 dalam $n$ langkah; jadi, rantai ini tidak reguler.
\end{example}
Contoh yang lebih menarik dari rantai Markov ergodik tak reguler diberikan oleh
model guci Ehrenfest.

\begin{example}
Ingat kembali model guci Ehrenfest (Contoh~\ref{exam 11.1.6}).\index{model
Ehrenfest}
\index{difusi gas!model Ehrenfest} Matriks transisi untuk contoh ini adalah
$$
\mat{P} = \bordermatrix{
  &  0  &  1  &  2  &  3  &  4  \cr
0 &  0  &  1  &  0  &  0  &  0  \cr
1 & 1/4 &  0  & 3/4 &  0  &  0  \cr
2 &  0  & 1/2 &  0  & 1/2 &  0  \cr
3 &  0  &  0  & 3/4 &  0  & 1/4 \cr
4 &  0  &  0  &  0  &  1  &  0}\ .
$$
Dalam contoh ini, jika kita mulai dalam keadaan~0, setelah sebarang jumlah
langkah genap kita akan berada dalam keadaan 0, 2, atau 4, dan setelah sebarang
jumlah langkah ganjil kita akan berada dalam keadaan 1~atau~3. Jadi, rantai ini
ergodik tetapi tidak reguler.
\end{example}

\subsection*{Rantai Markov Reguler}

Setiap matriks transisi yang tidak memiliki entri nol menentukan rantai Markov
reguler. Namun, suatu rantai Markov reguler mungkin memiliki matriks transisi
yang memuat entri nol. Matriks transisi pada contoh Negeri Oz dalam
Bagian~\ref{sec 11.1} memiliki $p_{NN} = 0$, tetapi pangkat keduanya
$\mat{P}^2$ tidak memiliki entri nol, sehingga rantai ini merupakan rantai
Markov reguler.
\par
Salah satu contoh rantai Markov tak reguler adalah rantai penyerap. Sebagai
contoh, misalkan
$$
\mat{P} = \pmatrix{
1 & 0 \cr
1/2 & 1/2 \cr}
$$
merupakan matriks transisi suatu rantai Markov. Maka semua pangkat $\mat{P}$
memiliki 0 di sudut kanan atas.
\par
Sekarang kita akan membahas dua teorema penting mengenai rantai reguler.

\begin{theorem}\label{thm 11.3.6}
Misalkan $\mat{P}$ adalah matriks transisi bagi suatu rantai reguler. Ketika
$n \to \infty$, pangkat-pangkat $\mat{P}^n$ mendekati suatu matriks limit
$\mat{W}$ yang semua barisnya merupakan vektor $\mat{w}$ yang sama. Vektor
$\mat{w}$ adalah vektor peluang positif ketat (yakni, semua komponennya positif
dan jumlahnya satu).
\end{theorem}
\par
Dalam bagian berikutnya, kita memberikan dua bukti bagi teorema fundamental ini.
Di sini kita memberikan gagasan dasar bukti pertama.
\par
Kita ingin menunjukkan bahwa pangkat-pangkat $\mat{P}^n$ dari suatu matriks
transisi reguler menuju matriks yang semua barisnya sama. Hal ini setara dengan
menunjukkan bahwa $\mat{P}^n$ konvergen ke matriks dengan kolom-kolom konstan.
Kolom ke-$j$ dari $\mat{P}^n$ adalah $\mat{P}^{n} \mat{y}$, dengan $\mat{y}$
sebagai vektor kolom yang memiliki $1$ pada entri ke-$j$ dan 0 pada entri
lainnya. Jadi, kita hanya perlu membuktikan bahwa untuk sebarang vektor kolom
$\mat{y},\mat{P}^{n} \mat{y}$ mendekati vektor konstan ketika $n$ menuju tak
hingga.
\par
Karena setiap baris $\mat{P}$ merupakan vektor peluang, $\mat{Py}$ mengganti
$\mat{y}$ dengan rata-rata komponennya. Berikut adalah contohnya:
$$\pmatrix{1/2 & 1/4 & 1/4 \cr 1/3 & 1/3 & 1/3 \cr 1/2 & 1/2 & 0\cr} \pmatrix
{1 \cr 2 \cr
3 \cr} = 
\pmatrix{1/2 \cdot 1+ 1/4 \cdot 2+ 1/4 \cdot 3\cr 
         1/3 \cdot 1+ 1/3 \cdot 2+ 1/3 \cdot 3\cr 
         1/2 \cdot 1+ 1/2 \cdot 2+ 0   \cdot 3\cr}
=\pmatrix {7/4 \cr 2 \cr 3/2 \cr}\ .   
$$
Hasil proses perataan ini membuat komponen-komponen $\mat{Py}$ lebih serupa
daripada komponen-komponen $\mat{y}$. Secara khusus, komponen maksimum menurun
(dari 3 menjadi 2) dan komponen minimum meningkat (dari 1 menjadi 3/2). Bukti
kita akan menunjukkan bahwa semakin sering perataan ini dilakukan untuk
memperoleh $\mat{P}^{n} \mat{y}$, selisih antara komponen maksimum dan minimum
akan menuju 0 ketika $n \rightarrow \infty$. Artinya, $\mat{P}^{n} \mat{y}$
menuju vektor konstan. Entri ke-$ij$ dari $\mat{P}^n$, yaitu $p_{ij}^{(n)}$,
adalah peluang bahwa proses berada dalam keadaan~$s_j$ setelah $n$ langkah jika
dimulai dalam keadaan~$s_i$. Jika baris yang sama pada $\mat{W}$ dinyatakan dengan
$\mat{w}$, Teorema~\ref{thm 11.3.6} menyatakan bahwa peluang jangka panjang
berada dalam~$s_j$ kira-kira~$w_j$, yaitu entri ke-$j$ dari $\mat{w}$, dan tidak
bergantung pada keadaan awal.

\begin{example}\label{exam 11.3.1}
Ingatlah bahwa untuk contoh Negeri Oz pada Bagian~\ref{sec 11.1}, pangkat keenam
matriks transisi~$\mat{P}$, hingga tiga tempat desimal, adalah
$$
\mat{P}^6 = \bordermatrix{
         & \mbox{R} & \mbox{N} & \mbox{S} \cr
\mbox{R} & .4 & .2 & .4 \cr
\mbox{N} & .4 & .2 & .4 \cr
\mbox{S}  & .4 & .2 & .4}\ .
$$
Jadi, pada tingkat ketelitian ini, peluang hujan enam hari setelah hari hujan
sama dengan peluang hujan enam hari setelah hari cerah ataupun enam hari setelah
hari bersalju. Teorema~\ref{thm 11.3.6} memprediksi bahwa untuk $n$ besar,
baris-baris~$\mat{P}$ mendekati suatu vektor bersama. Menarik bahwa hal ini
terjadi begitu cepat dalam contoh kita.
\end{example}

\begin{theorem}\label{thm 11.3.8}
Misalkan $\mat {P}$ adalah matriks transisi reguler, misalkan
$$\mat{W} = \lim_{n \rightarrow \infty} \mat{P}^n\ ,$$
misalkan $\mat{w}$ adalah baris yang sama pada $\mat{W}$, dan misalkan $\mat{c}$ adalah
vektor kolom yang semua komponennya bernilai 1. Maka
\begin{description}
\item[(a)]
$\mat{w}\mat{P} = \mat{w}$, dan setiap vektor baris~$\mat{v}$ sedemikian rupa
sehingga $\mat{v}\mat{P} = \mat{v}$ merupakan kelipatan konstan dari~$\mat{w}$.
\item[(b)]
$\mat{P}\mat{c} = \mat{c}$, dan setiap vektor kolom~\mat{x} sedemikian rupa
sehingga $\mat{P}\mat{x} = \mat{x}$ merupakan kelipatan dari~$\mat{c}$.
\end{description}
\proof
Untuk membuktikan bagian (a), dari Teorema~\ref{thm 11.3.6} kita perhatikan bahwa
$$
\mat{P}^n \to \mat{W}\ .
$$
Jadi,
$$
\mat{P}^{n + 1} = \mat{P}^n \cdot \mat{P} \to \mat{W}\mat{P}\ .
$$
Namun $\mat{P}^{n + 1} \to \mat{W}$, sehingga $\mat{W} = \mat{W}\mat{P}$, dan
$\mat{w} =
\mat{w}\mat{P}$. 
\par
Misalkan $\mat{v}$ adalah sebarang vektor 
dengan $\mat{v} \mat{P} = \mat{v}$. Maka $\mat{v} = \mat{v} \mat{P}^n$, dan
dengan mengambil limit,
$\mat{v} = \mat{v} \mat{W}$. Misalkan $r$ adalah jumlah komponen-komponen
$\mat{v}$.
Mudah diperiksa bahwa $\mat{v}\mat{W} = r\mat{w}$. Jadi, $\mat{v} =
r\mat{w}$.
\par
Untuk membuktikan bagian (b), andaikan $\mat{x} = \mat{P} \mat{x}$. Maka $\mat{x} =
\mat{P}^n
\mat{x}$, dan sekali lagi dengan mengambil limit, $\mat{x} =  \mat{W}\mat{x}$.
Karena semua baris
$\mat{W}$ sama, komponen-komponen $\mat{W}\mat{x}$ semuanya sama, sehingga
$\mat{x}$ merupakan
kelipatan $\mat{c}$.
\end{theorem}

Perhatikan bahwa akibat langsung Teorema~\ref{thm 11.3.8} adalah hanya ada
satu vektor peluang $\mat{v}$ sedemikian rupa sehingga
$\mat{v}\mat{P} = \mat{v}$.

\subsection*{Vektor Tetap}
\begin{definition}\label{def 11.3.1}
Vektor baris~$\mat{w}$ dengan sifat $\mat{w}\mat{P} = \mat{w}$ disebut
\emx {vektor baris tetap} untuk~$\mat{P}$.\index{vektor baris tetap} Demikian pula,
vektor kolom~$\mat{x}$ sedemikian rupa sehingga $\mat{P}\mat{x} = \mat{x}$
disebut \emx {vektor kolom tetap} untuk~$\mat{P}$.\index{vektor kolom tetap}  
\end{definition}

Jadi, baris yang sama pada~$\mat{W}$ adalah satu-satunya vektor~$\mat{w}$ yang
sekaligus merupakan vektor baris tetap untuk~$\mat{P}$ dan vektor peluang.
Teorema~\ref{thm 11.3.8} menunjukkan bahwa setiap vektor baris tetap untuk
$\mat{P}$ merupakan kelipatan~$\mat{w}$ dan setiap vektor kolom tetap untuk
$\mat{P}$ merupakan vektor konstan.
\par
Definisi~\ref{def 11.3.1} juga dapat dinyatakan dalam bentuk nilai eigen dan
vektor eigen. Vektor baris tetap adalah vektor eigen kiri dari matriks
$\mat{P}$ yang bersesuaian dengan nilai eigen 1. Pernyataan serupa dapat dibuat
untuk vektor kolom tetap.  
\par
Sekarang kita akan memberikan beberapa metode berbeda untuk menghitung vektor
baris tetap \mat{w} bagi suatu rantai Markov reguler.

\begin{example}\label{exam 11.3.2}
Menurut Teorema~\ref{thm 11.3.6}, kita dapat menemukan vektor limit~$\mat{w}$
untuk Negeri Oz dari kenyataan bahwa
$$
w_1 + w_2 + w_3 = 1
$$
dan
$$
 \pmatrix{ w_1 & w_2 & w_3 } \pmatrix{
 1/2 & 1/4 & 1/4 \cr
 1/2 & 0   & 1/2 \cr
 1/4 & 1/4 & 1/2 \cr} = \pmatrix{ w_1 & w_2 & w_3 }\ .
$$

Relasi-relasi ini menghasilkan empat persamaan berikut dalam tiga peubah tak dikenal:
\begin{eqnarray*}
                            w_1   + w_2 + w_3 &=& 1\ ,   \\
             (1/2)w_1 + (1/2)w_2  + (1/4)w_3  &=& w_1\ , \\
                        (1/4)w_1  + (1/4)w_3  &=& w_2\ , \\
             (1/4)w_1 + (1/2)w_2  + (1/2)w_3  &=& w_3\ .
\end{eqnarray*}

Teorema kita menjamin bahwa persamaan-persamaan ini mempunyai solusi tunggal.
Dengan menyelesaikannya, kita memperoleh solusi
$$
\mat{w} = \pmatrix{ .4 & .2 & .4 }\ ,
$$
yang sesuai dengan prediksi dari~$\mat{P}^6$ dalam Contoh~\ref{exam 11.1.1.5}.
\end{example}

Untuk menghitung vektor tetap, kita dapat mengandaikan bahwa nilai pada suatu
keadaan tertentu, misalnya keadaan satu, adalah~1, lalu menggunakan semua
persamaan linear dari $\mat{w}\mat{P} = \mat{w}$ kecuali satu. Sistem persamaan
ini akan mempunyai solusi tunggal, dan kita dapat memperoleh $\mat{w}$ dari
solusi itu dengan membagi setiap entrinya dengan jumlah semua entri sehingga
menghasilkan vektor peluang~$\mat{w}$. Sekarang kita ilustrasikan gagasan ini
untuk contoh di atas. 
\begin{example}\label{exam 11.3.2.5}(Lanjutan Contoh~\ref{exam 11.3.2})
Kita tetapkan $w_1 = 1$, lalu menyelesaikan persamaan linear pertama dan kedua dari
$\mat{w}\mat{P} =
\mat{w}$. Kita memperoleh
\begin{eqnarray*}
             (1/2) + (1/2)w_2  + (1/4)w_3  &=& 1\ , \\
                        (1/4)  + (1/4)w_3  &=& w_2\ . \\
\end{eqnarray*}
Dengan menyelesaikannya, kita memperoleh 
$$\pmatrix{w_1&w_2&w_3} = \pmatrix{1&1/2&1}\ .$$
Sekarang kita bagi vektor ini dengan jumlah komponennya untuk memperoleh
jawaban akhir:
$$\mat{w} = \pmatrix{.4&.2&.4}\ .$$
Metode ini dapat dengan mudah diprogram untuk dijalankan pada komputer.
\end{example}
\par
Seperti disebutkan di atas, vektor baris tetap $\mat{w}$ juga dapat dipandang
sebagai vektor eigen kiri dari matriks transisi $\mat{P}$. Jadi, jika
$\mat{I}$ menyatakan matriks identitas, maka $\mat{w}$ memenuhi persamaan matriks
$$\mat{w}\mat{P} = \mat{w}\mat{I}\ ,$$
atau, secara ekuivalen,
$$\mat{w}(\mat{P} - \mat{I}) = \mat{0}\ .$$
Jadi, $\mat{w}$ berada dalam ruang nol kiri matriks $\mat{P} - \mat{I}$.
Selain itu, Teorema~\ref{thm 11.3.8} menyatakan bahwa ruang nol kiri ini
berdimensi 1. Bahasa pemrograman komputer tertentu dapat menemukan ruang nol
suatu matriks. Dalam bahasa tersebut, vektor peluang baris tetap untuk matriks
$\mat{P}$ dapat ditemukan dengan menghitung ruang nol kiri, lalu menormalkan
sebuah vektor di dalam ruang nol itu agar jumlah komponennya sama dengan 1.
\par
Program {\bf FixedVector}\index{FixedVector (program)} menggunakan salah satu
metode di atas (bergantung pada bahasa penulisannya) untuk menghitung vektor
peluang baris tetap bagi rantai Markov reguler.
\par
Sejauh ini kita selalu mengandaikan bahwa proses dimulai dalam keadaan tertentu.
Teorema berikut memperumum Teorema~\ref{thm 11.3.6} ke kasus ketika keadaan awal
itu sendiri ditentukan oleh suatu vektor peluang.

\begin{theorem}\label{thm 11.3.9}
Misalkan $\mat{P}$ adalah matriks transisi untuk suatu rantai reguler dan
$\mat{v}$ suatu vektor peluang sebarang. Maka
$$
\lim_{n \to \infty} \mat{v} \mat{P}^n = \mat{w}\ ,
$$
di mana $\mat{w}$ adalah vektor peluang tetap tunggal untuk $\mat{P}$.
\proof
Menurut Teorema~\ref{thm 11.3.6},
$$
\lim_{n \to \infty} \mat{P}^n = \mat {W}\ .
$$
Oleh karena itu,
$$
\lim_{n \to \infty} \mat {v} \mat {P}^n = \mat {v} \mat {W}\ .
$$
Namun jumlah entri-entri $\mat {v}$ sama dengan 1, dan setiap baris $\mat {W}$
sama dengan $\mat {w}$. Dari pernyataan-pernyataan ini, mudah diperiksa bahwa
$$
\mat {v} \mat {W} = \mat {w}\ .
$$
\end{theorem}

Jika kita memulai rantai Markov dengan peluang awal yang diberikan oleh
~$\mat{v}$, maka vektor peluang $\mat{v} \mat{P}^n$ memberikan peluang berada
dalam berbagai keadaan setelah $n$~langkah. Teorema~\ref{thm 11.3.9} dengan
demikian menetapkan bahwa, bahkan dalam kelas proses yang lebih umum ini,
peluang berada dalam~$s_j$ mendekati $w_j$.

\subsection*{Kesetimbangan}
Kita juga memperoleh penafsiran baru untuk~$\mat{w}$. Andaikan vektor awal kita
memilih keadaan~$s_i$ sebagai keadaan awal dengan peluang~$w_i$, untuk setiap
~$i$. Maka peluang berada dalam berbagai keadaan setelah $n$ langkah diberikan
oleh $\mat {w} \mat {P}^n = \mat {w}$ dan sama pada setiap langkah. Cara memulai
ini menghasilkan proses yang disebut ``stasioner.'' Kenyataan bahwa $\mat{w}$
adalah satu-satunya vektor peluang yang memenuhi $\mat {w} \mat {P} =
\mat {w}$ menunjukkan bahwa untuk memperoleh proses stasioner, kita harus
mempunyai vektor peluang awal tepat seperti yang dijelaskan.
\par
Banyak hasil menarik mengenai rantai Markov reguler hanya bergantung pada
kenyataan bahwa rantai itu mempunyai vektor peluang tetap positif yang tunggal.
Sifat ini berlaku untuk semua rantai Markov ergodik.

\begin{theorem}\label{thm 11.3.10}
Untuk suatu rantai Markov ergodik, terdapat vektor peluang tunggal~$\mat{w}$
sedemikian rupa sehingga $\mat {w} \mat {P} = \mat {w}$ dan $\mat{w}$ positif
ketat. Setiap vektor baris sedemikian rupa sehingga
$\mat {v} \mat {P} = \mat {v}$ merupakan kelipatan~$\mat{w}$. Setiap vektor
kolom~$\mat{x}$ sedemikian rupa sehingga $\mat {P} \mat {x} = \mat {x}$
merupakan vektor konstan.
\proof
Teorema ini menyatakan bahwa Teorema~\ref{thm 11.3.8} berlaku untuk rantai
ergodik. Hasilnya mudah diperoleh dari kenyataan bahwa, jika $\mat{P}$ adalah
matriks transisi ergodik, maka
$\bar{\mat{P}} = (1/2)\mat {I} + (1/2)\mat {P}$ adalah matriks transisi reguler
dengan vektor-vektor tetap yang sama (lihat Latihan~\ref{exer
11.3.24}--\ref{exer 11.3.27}).
\end{theorem}
\par
Untuk rantai ergodik, vektor peluang tetap mempunyai penafsiran yang agak
berbeda. Dua teorema berikut, yang tidak akan kita buktikan di sini, memberikan
penafsiran bagi vektor tetap ini.

\begin{theorem}\label{thm 11.3.11}
Misalkan $\mat{P}$ adalah matriks transisi untuk suatu rantai ergodik. Misalkan
$\mat {A}_n$ adalah matriks yang didefinisikan oleh
$$
\mat {A}_n = \frac{\mat {I} + \mat {P} + \mat {P}^2 +\cdots + \mat {P}^n}{n +
1}\ .
$$
Maka $\mat {A}_n \to \mat {W}$, di mana $\mat{W}$ adalah matriks yang semua
barisnya sama dengan vektor peluang tetap tunggal $\mat{w}$ untuk $\mat{P}$.  
\end{theorem}

\par
Jika $\mat{P}$ adalah matriks transisi suatu rantai ergodik, maka
Teorema~\ref{thm 11.3.8} menyatakan bahwa hanya ada satu vektor peluang baris
tetap untuk $\mat{P}$. Jadi, kita dapat menggunakan teknik yang sama seperti
untuk rantai reguler guna mencari vektor tetap ini. Secara khusus, program
{\bf FixedVector} dapat digunakan untuk rantai ergodik.
\par
Untuk menafsirkan Teorema~\ref{thm 11.3.11}, andaikan kita mempunyai rantai
ergodik yang dimulai dalam keadaan~$s_i$. Misalkan $X^{(m)} = 1$ jika langkah
ke-$m$ menuju keadaan~$s_j$ dan 0 jika tidak. Maka rata-rata banyaknya kali
berada dalam keadaan~$s_j$ selama $n$ langkah pertama diberikan oleh
$$
H^{(n)} = \frac{X^{(0)} + X^{(1)} + X^{(2)} +\cdots+ X^{(n)}}{n + 1}\ .
$$
Namun $X^{(m)}$ bernilai~1 dengan peluang~$p_{ij}^{(m)}$ dan 0 jika tidak.
Jadi $E(X^{(m)}) = p_{ij}^{(m)}$, dan entri ke-$ij$ dari~$\mat {A}_n$
memberikan nilai harapan~$H^{(n)}$, yaitu proporsi waktu yang diharapkan dalam
keadaan~$s_j$ selama $n$ langkah pertama jika rantai dimulai dalam keadaan~$s_i$.

Jika berada dalam keadaan~$s_j$ kita sebut \emx{keberhasilan} dan berada dalam
keadaan lain \emx{kegagalan,} kita dapat menanyakan apakah berlaku teorema yang
analog dengan hukum bilangan besar untuk percobaan independen. Jawabannya ya,
dan diberikan oleh teorema berikut.

\begin{theorem}{\bf (Hukum Bilangan Besar untuk Rantai Markov Ergodik)}\label{thm
11.3.12}
\index{Hukum Bilangan Besar!untuk Rantai Markov Ergodik}
Misalkan $H_j^{(n)}$ adalah proporsi waktu selama $n$ langkah ketika suatu
rantai ergodik berada dalam keadaan~$s_j$. Maka, untuk setiap $\epsilon > 0$,
$$
P\Bigl(|H_j^{(n)} - w_j| > \epsilon\Bigr) \to 0\ ,
$$
terlepas dari keadaan awal~$s_i$.
\end{theorem}

Kita telah melihat bahwa setiap rantai Markov reguler juga merupakan rantai
ergodik. Oleh karena itu, Teorema~\ref{thm 11.3.11}~dan~\ref{thm 11.3.12} juga
berlaku untuk rantai reguler. Sebagai contoh, hal ini memberi kita penafsiran
baru untuk vektor tetap $\mat {w} = (.4,.2,.4)$ dalam contoh Negeri Oz.
Teorema~\ref{thm 11.3.11} memprediksi bahwa dalam jangka panjang, di Negeri Oz
hujan akan turun 40~persen dari waktu, cuaca akan cerah 20~persen dari waktu,
dan salju akan turun 40~persen dari waktu.

\subsection*{Simulasi}
Kita mengilustrasikan Teorema~\ref{thm 11.3.12} dengan menulis program untuk
menyimulasikan perilaku rantai Markov. {\bf SimulateChain}\index{SimulateChain (program)}
adalah program semacam itu.  

\begin{example}\label{exam 11.3.2.6}
Di Negeri Oz terdapat 525 hari dalam setahun.
\index{Oz, Negeri} Dengan menggunakan program {\bf SimulateChain}, kita telah
menyimulasikan cuaca selama satu tahun di Negeri Oz. Hasilnya ditunjukkan dalam
Tabel~\ref{table 11.2}. 

\begin{table}[h]
\centering
\begin{verbatim}
     SSRNRNSSSSSSNRSNSSRNSRNSSSNSRRRNSSSNRRSSSSNRSSNSRRRRRRNSSS 
     SSRRRSNSNRRRRSRSRNSNSRRNRRNRSSNSRNRNSSRRSRNSSSNRSRRSSNRSNR
     RNSSSSNSSNSRSRRNSSNSSRNSSRRNRRRSRNRRRNSSSNRNSRNSNRNRSSSRSS
     NRSSSNSSSSSSNSSSNSNSRRNRNRRRRSRRRSSSSNRRSSSSRSRRRNRRRSSSSR
     RNRRRSRSSRRRRSSRNRRRRRRNSSRNRSSSNRNSNRRRRNRRRNRSNRRNSRRSNR
     RRRSSSRNRRRNSNSSSSSRRRRSRNRSSRRRRSSSRRRNRNRRRSRSRNSNSSRRRR
     RNSNRNSNRRNRRRRRRSSSNRSSRSNRSSSNSNRNSNSSSNRRSRRRNRRRRNRNRS
     SSNSRSNRNRRSNRRNSRSSSRNSRRSSNSRRRNRRSNRRNSSSSSNRNSSSSSSSNR
     NSRRRNSSRRRNSSSNRRSRNSSRRNRRNRSNRRRRRRRRRNSNRRRRRNSRRSSSSN
     SNS
\end{verbatim}
\begin{tabular}{lcc}
Keadaan  &   Frekuensi    &      Proporsi\\
         &                &           \\
R        &        217     &      .413 \\
N        &        109     &      .208 \\
S        &        199     &      .379 \\
\end{tabular}
\caption{Cuaca di Negeri Oz.}
\label{table 11.2}
\end{table}
 
Kita perhatikan bahwa simulasi ini memberikan proporsi waktu dalam setiap
keadaan yang tidak terlalu berbeda dari prediksi jangka panjang .4,~.2,
dan~.4 yang dijamin oleh Teorema~\ref{thm 11.3.6}. Untuk memperoleh hasil yang
lebih baik, kita harus menyimulasikan rantai selama waktu yang lebih panjang.
Kita melakukannya selama 10{,}000 hari tanpa mencetak cuaca setiap hari.
Hasilnya ditunjukkan dalam Tabel~\ref{table 11.3}. Kita melihat bahwa hasilnya
sekarang cukup dekat dengan nilai teoretis .4,~.2, dan~.4.

\begin{table}[h]
\centering
\begin{tabular}{lcc}
Keadaan &   Frekuensi   &     Proporsi\\
      &                &           \\
R     &      4010      &      .401 \\
N     &      1902      &      .19  \\
S     &      4088      &      .409 \\
\end{tabular}
\caption{Perbandingan frekuensi teramati dan frekuensi prediksi untuk Negeri Oz.}
\label{table 11.3}
\end{table}
\end{example}

\subsection*{Contoh Rantai Ergodik}

Perhitungan vektor tetap~$\mat{w}$ mungkin sulit jika matriks transisinya sangat
besar. Kadang-kadang berguna untuk menduga vektor tetap semata-mata berdasarkan
intuisi. Berikut contoh sederhana untuk mengilustrasikan situasi semacam ini.

\begin{example}\label{exam 11.3.3}
Seekor tikus putih ditempatkan dalam labirin pada Gambar~\ref{fig
11.4}.\index{tikus}\index{labirin} Terdapat sembilan kompartemen dengan hubungan
antar-kompartemen seperti yang ditunjukkan. Tikus bergerak secara acak melalui
kompartemen-kompartemen itu. Artinya, jika terdapat $k$~jalan keluar dari suatu
kompartemen, tikus memilih masing-masing dengan peluang yang sama. Perjalanan
tikus dapat kita representasikan sebagai proses rantai Markov dengan matriks
transisi
$$
\mat {P} = \bordermatrix{
& 1 & 2 & 3 & 4 & 5 & 6 & 7 & 8 & 9 \cr
1 & 0 & 1/2 & 0 & 0 & 0 & 1/2 & 0 & 0 & 0 \cr
2 & 1/3 & 0 & 1/3 & 0 & 1/3 & 0 & 0 & 0 & 0 \cr
3 & 0 & 1/2 & 0 & 1/2 & 0 & 0 & 0 & 0 & 0 \cr
4 & 0 & 0 & 1/3 & 0 & 1/3 & 0 & 0 & 0 & 1/3 \cr
5 & 0 & 1/4 & 0 & 1/4 & 0 & 1/4 & 0 & 1/4 & 0 \cr
6 & 1/3 & 0 & 0 & 0 & 1/3 & 0 & 1/3 & 0 & 0 \cr
7 & 0 & 0 & 0 & 0 & 0 & 1/2 & 0 & 1/2 & 0 \cr
8 & 0 & 0 & 0 & 0 & 1/3 & 0 & 1/3 & 0 & 1/3 \cr
9 & 0 & 0 & 0 & 1/2 & 0 & 0 & 0 & 1/2 & 0}\ .
$$

\putfig{2truein}{PSfig11-4}{Masalah labirin.}{fig 11.4}
 
%\begin{figure}
%\centerline{\epsfxsize=2truein 
%\epsffile{PSfig11-4}}
%\caption{The maze problem.}\label{fig 11.4}
%\end{figure}
\par
Bahwa rantai ini tidak reguler dapat dilihat sebagai berikut: dari keadaan
bernomor ganjil, proses hanya dapat menuju keadaan bernomor genap, dan dari
keadaan bernomor genap proses hanya dapat menuju keadaan bernomor ganjil.
Karena itu, jika dimulai dalam keadaan $i$, proses akan bergantian berada dalam
keadaan bernomor genap dan bernomor ganjil. Dengan demikian, pangkat ganjil
dari~$\mat{P}$ akan mempunyai 0 pada entri bernomor ganjil di baris~1. Di sisi
lain, sekilas melihat labirin menunjukkan bahwa setiap keadaan dapat dicapai
dari setiap keadaan lain, sehingga rantai ini ergodik.
\par
Untuk menemukan vektor peluang tetap bagi matriks ini, kita harus menyelesaikan
sepuluh persamaan dalam sembilan peubah tak dikenal. Namun, tampaknya masuk akal
bahwa dalam jangka panjang, waktu yang dihabiskan dalam setiap kompartemen
sebanding dengan banyaknya jalan masuk ke kompartemen tersebut. Karena itu,
kita mencoba vektor yang komponen ke-$j$-nya adalah banyaknya jalan masuk ke
kompartemen ke-$j$:
$$
\mat {x} = \pmatrix{ 2 & 3 & 2 & 3 & 4 & 3 & 2 & 3 & 2}\ .
$$
Mudah diperiksa bahwa vektor ini memang vektor tetap, sehingga vektor peluang
tunggalnya adalah vektor ini yang dinormalkan agar berjumlah~1:
$$
\mat {w} = \pmatrix{ \frac1{12} & \frac18 & \frac1{12} & \frac18 & \frac16 &
\frac18 & \frac 1{12} & \frac18 & \frac1{12} }\ .
$$
\end{example}

\begin{example}\label{exam 11.3.4}(Lanjutan Contoh~\ref{exam 11.1.6})
Ingat kembali model guci Ehrenfest pada Contoh~\ref{exam 11.1.6}.
\index{model Ehrenfest}\index{difusi gas!model Ehrenfest untuk}
Matriks transisi untuk rantai ini adalah sebagai berikut:

$$
\mat {P} = \bordermatrix{
  & 0    & 1     &  2   & 3    & 4   \cr
0 & .000 & 1.000 & .000 & .000 & .000\cr
1 & .250 &  .000 & .750 & .000 & .000\cr
2 & .000 &  .500 & .000 & .500 & .000\cr
3 & .000 &  .000 & .750 & .000 & .250\cr
4 & .000 &  .000 & .000 &1.000 & .000}\ .
$$
Jika kita menjalankan program {\bf FixedVector} untuk rantai ini, kita
memperoleh vektor
$$
\mat {w} = \bordermatrix{ 
  &     0  &    1  &    2  &    3  &    4\cr
  & .0625  &.2500  &.3750  &.2500  &.0625}\ .
$$

Menurut Teorema~\ref{thm 11.3.12}, nilai-nilai~$w_i$ ini dapat kita tafsirkan
sebagai proporsi waktu ketika proses berada dalam setiap keadaan dalam jangka
panjang. Sebagai contoh, proporsi waktu dalam keadaan~0 adalah .0625 dan
proporsi waktu dalam keadaan~1 adalah .375. Pembaca yang cermat akan melihat
bahwa angka-angka ini membentuk distribusi binomial
1/16,~4/16, 6/16, 4/16,~1/16. Jawaban ini dapat kita duga sebagai berikut: jika
kita memperhatikan satu bola tertentu, bola itu hanya bergerak bolak-balik
secara acak di antara kedua guci. Hal ini menunjukkan bahwa keadaan
kesetimbangan seharusnya sama seperti jika keempat bola didistribusikan secara
acak ke dalam kedua guci. Jika kita melakukannya, peluang terdapat tepat
$j$~bola dalam salah satu guci diberikan oleh distribusi binomial $b(n,p,j)$
dengan $n = 4$ dan $p = 1/2$.
\end{example}

\exercises
\begin{LJSItem}

\i\label{exer 11.3.1} Manakah di antara matriks-matriks berikut yang merupakan
matriks transisi rantai Markov reguler?
\begin{enumerate}
\item $\mat {P} = \pmatrix{ .5 & .5 \cr .5 & .5 }$.
\smallskip
\item $\mat {P} = \pmatrix{ .5 & .5 \cr 1 & 0 }$.
\smallskip
\item $\mat {P} = \pmatrix{ 1/3 & 0 & 2/3 \cr 0 & 1 & 0 \cr 0 & 1/5 & 4/5}$.
\smallskip
\item $\mat {P} = \pmatrix{ 0 & 1 \cr 1 & 0}$.
\smallskip
\item $\mat {P} = \pmatrix{ 1/2 & 1/2 & 0 \cr 0 & 1/2 & 1/2 \cr 1/3 & 1/3 &
1/3}$.
\smallskip
\end{enumerate}

\i\label{exer 11.3.2} Perhatikan rantai Markov dengan matriks transisi
$$
\mat {P} = \pmatrix{ 1/2 & 1/3 & 1/6 \cr3/4 & 0 & 1/4 \cr 0 & 1 & 0}\ .
$$
\begin{enumerate}

\item Tunjukkan bahwa ini adalah rantai Markov reguler.

\item Proses dimulai dalam keadaan~1; carilah peluang bahwa proses berada dalam
keadaan~3 setelah dua langkah.

\item Carilah vektor peluang limit~$\mat{w}$.
\end{enumerate}

\i\label{exer 11.3.3} Perhatikan rantai Markov dengan matriks transisi umum
$2 \times 2$
$$
\mat {P} = \pmatrix{ 1 - a & a \cr b & 1 - b}\ .
$$
\begin{enumerate}

\item Dalam kondisi apa $\mat{P}$ bersifat menyerap?

\item Dalam kondisi apa $\mat{P}$ bersifat ergodik tetapi tidak reguler?

\item Dalam kondisi apa $\mat{P}$ bersifat reguler?
\end{enumerate}

\i\label{exer 11.3.4} Carilah vektor peluang tetap~$\mat{w}$ untuk
matriks-matriks ergodik dalam Latihan~\ref{exer 11.3.3}.

\i\label{exer 11.3.5} Carilah vektor peluang tetap~$\mat{w}$ untuk setiap
matriks reguler berikut.
\begin{enumerate}
\item $\mat {P} = \pmatrix{ .75 & .25 \cr .5 & .5}$.
\smallskip
\item $\mat {P} = \pmatrix{ .9 & .1 \cr .1 & .9}$.
\smallskip
\item $\mat {P} = \pmatrix{ 3/4 & 1/4 & 0 \cr 0 & 2/3 & 1/3 \cr 1/4 & 1/4 &
1/2}$.
\smallskip
\end{enumerate}

\i\label{exer 11.3.6} Perhatikan rantai Markov dengan matriks transisi dalam
Latihan~\ref{exer 11.3.3}, dengan $a = b = 1$. Tunjukkan bahwa rantai ini
ergodik tetapi tidak reguler. Carilah vektor peluang tetapnya dan tafsirkan.
Tunjukkan bahwa $\mat {P}^n$ tidak menuju suatu limit, tetapi
$$
\mat {A}_n = \frac{\mat {I} + \mat {P} + \mat {P}^2 +\cdots + \mat {P}^n}{n +
1}
$$
menuju suatu limit.

\i\label{exer 11.3.7} Perhatikan rantai Markov dengan matriks transisi dalam
Latihan~\ref{exer 11.3.3}, dengan $a = 0$ dan $b = 1/2$. Hitunglah secara
langsung vektor peluang tetap tunggalnya, lalu gunakan hasil Anda untuk
membuktikan bahwa rantai itu tidak ergodik.

\i\label{exer 11.3.8} Tunjukkan bahwa matriks
$$
\mat {P} = \pmatrix{ 1 & 0 & 0 \cr 1/4 & 1/2 & 1/4 \cr 0 & 0 & 1}
$$
mempunyai lebih dari satu vektor peluang tetap. Carilah matriks yang didekati
oleh $\mat {P}^n$ ketika $n \to \infty$, dan verifikasikan bahwa matriks itu
tidak mempunyai baris-baris yang semuanya sama.

\i\label{exer 11.3.9} Buktikan bahwa jika suatu matriks transisi 3-kali-3
mempunyai sifat bahwa jumlah setiap \emx {kolomnya} adalah~1, maka
$(1/3, 1/3, 1/3)$ merupakan vektor peluang tetap. Nyatakan hasil serupa untuk
matriks transisi $n$-kali-$n$. Tafsirkan hasil-hasil ini untuk rantai ergodik.

\i\label{exer 11.3.10} Apakah rantai Markov dalam Contoh~\ref{exam 11.1.8}
bersifat ergodik?

\i\label{exer 11.3.10.5} Apakah rantai Markov dalam Contoh~\ref{exam 11.1.9}
bersifat ergodik?

\i\label{exer 11.3.11} Perhatikan Contoh~\ref{exam 11.2.1} (Perjalanan Si
Pemabuk).\index{contoh Perjalanan Si Pemabuk} Andaikan bahwa jika pejalan
mencapai keadaan~0, ia berbalik dan kembali ke keadaan~1 pada langkah
berikutnya; demikian pula, jika ia mencapai 4, ia kembali ke keadaan~3 pada
langkah berikutnya. Apakah rantai baru ini ergodik? Apakah rantai ini reguler?

\i\label{exer 11.3.12} Untuk Contoh~\ref{exam 11.1.2}, ketika $\mat{P}$
ergodik, berapakah proporsi orang yang diberi tahu bahwa Presiden akan mencalonkan
diri? Tafsirkan kenyataan bahwa proporsi ini tidak bergantung pada keadaan awal.

\i\label{exer 11.3.13} Pandanglah suatu proses percobaan independen sebagai
rantai Markov yang keadaannya adalah hasil-hasil yang mungkin dari setiap
percobaan. Apakah vektor peluang tetapnya? Apakah rantai itu selalu reguler?
Ilustrasikan untuk Contoh~\ref{exam 11.1.3}.

\i\label{exer 11.3.14} Tunjukkan bahwa Contoh~\ref{exam 11.1.6} adalah rantai
ergodik, tetapi bukan rantai reguler. Tunjukkan bahwa vektor peluang tetapnya
$\mat{w}$ merupakan distribusi binomial.

\i\label{exer 11.3.15} Tunjukkan bahwa Contoh~\ref{exam 11.1.7} bersifat
reguler dan carilah vektor limitnya.

\i\label{exer 11.3.16} Lemparkan sebuah dadu adil berulang kali. Misalkan
$S_n$ menyatakan jumlah hasil sampai lemparan ke-$n$. Tunjukkan bahwa terdapat
nilai limit bagi proporsi di antara $n$ nilai pertama~$S_n$ yang habis dibagi
~7, dan hitunglah nilai limit ini. \emx {Petunjuk}: Limit yang dicari adalah
vektor peluang kesetimbangan untuk rantai Markov tujuh-keadaan yang sesuai.

\i\label{exer 11.3.17} Misalkan $\mat{P}$ adalah matriks transisi suatu rantai
Markov reguler. Andaikan terdapat $r$~keadaan dan misalkan $N(r)$ adalah
bilangan bulat terkecil~$n$ sedemikian rupa sehingga $\mat{P}$ reguler jika dan
hanya jika $\mat {P}^{N(r)}$ tidak mempunyai entri nol. Carilah batas atas
hingga untuk~$N(r)$. Selidiki apakah Anda dapat menentukan $N(3)$ secara tepat.

\istar\label{exer 11.3.18} Definisikan $f(r)$ sebagai bilangan bulat terkecil
$n$ sedemikian rupa sehingga, untuk semua rantai Markov reguler dengan $r$
keadaan, pangkat ke-$n$ dari matriks transisinya mempunyai semua entri positif.
Telah dibuktikan,\footnote{E. Seneta, \emx
{Non-Negative Matrices:  An Introduction to Theory and Applications,}
Wiley, New York, 1973, pp.~52-54.\index{SENETA, E.}}
bahwa $f(r) = r^2 - 2r + 2$.  
\begin{enumerate}
\item
Definisikan matriks transisi suatu rantai Markov dengan $r$ keadaan sebagai
berikut: untuk keadaan $s_i$, dengan $i = 1$,~2, \ldots,~$r - 2$, 
$\mat {P}(i,i + 1) = 1$, $\mat {P}(r - 1,r) = \mat {P}(r - 1, 1) = 1/2$, dan 
$\mat {P}(r,1) = 1$. Tunjukkan bahwa ini adalah rantai Markov reguler.

\item Untuk $r = 3$, verifikasikan bahwa pangkat kelima adalah pangkat pertama
yang tidak mempunyai nol.

\item Tunjukkan bahwa, untuk $r$ umum, $n$ terkecil sedemikian rupa sehingga
$\mat {P}^n$ mempunyai semua entri positif adalah $n = f(r)$.
\end{enumerate}

\i\label{exer 11.3.19} Suatu sistem antrean waktu diskret berkapasitas~$n$
terdiri atas orang yang sedang dilayani dan mereka yang menunggu dilayani.
Panjang antrean~$x$ diamati setiap detik. Jika $0 < x < n$, maka dengan
peluang~$p$, ukuran antrean bertambah satu karena kedatangan dan, secara
independen, dengan peluang~$r$, berkurang satu karena orang yang sedang
dilayani selesai dilayani. Jika $x = 0$, hanya kedatangan (dengan peluang~$p$)
yang mungkin. Jika $x = n$, orang yang datang akan pergi tanpa menunggu
pelayanan, sehingga hanya keberangkatan (dengan peluang~$r$) orang yang sedang
dilayani yang mungkin. Bentuklah rantai Markov dengan keadaan yang diberikan
oleh banyaknya pelanggan dalam antrean. Ubahlah program {\bf FixedVector} agar
Anda dapat memasukkan $n$,~$p$, dan~$r$, dan program akan menyusun matriks
transisi serta menghitung vektor tetap. Besaran $s = p/r$ disebut
\emx {intensitas lalu lintas.} Jelaskan perbedaan vektor-vektor tetap ketika
$s < 1$, $s = 1$, atau $s > 1$.

\i\label{exer 11.3.20} Tulislah program komputer untuk menyimulasikan antrean
dalam Latihan~\ref{exer 11.3.19}. Minta program Anda melacak proporsi waktu
ketika panjang antrean adalah $j$, untuk $j = 0$,~1, \ldots,~$n$, serta panjang
antrean rata-rata. Tunjukkan bahwa perilaku panjang antrean sangat berbeda,
bergantung pada apakah intensitas lalu lintas~$s$ memenuhi $s < 1$, $s = 1$,
atau $s > 1$.

\i\label{exer 11.3.21} Dalam masalah antrean pada Latihan~\ref{exer 11.3.19},
misalkan $S$ adalah waktu pelayanan total yang diperlukan seorang pelanggan dan
$T$ adalah waktu antara kedatangan pelanggan.
\begin{enumerate}
\item Tunjukkan bahwa $P(S = j) = (1 - r)^{j - 1}r$ dan $P(T = j) = (1 - p)^{j -
1}p$, untuk $j > 0$.

\item Tunjukkan bahwa $E(S) = 1/r$ dan $E(T) = 1/p$.

\item Tafsirkan kondisi $s < 1$, $s = 1$, dan $s > 1$ dalam bentuk nilai-nilai
harapan ini.
\end{enumerate}

\i\label{exer 11.3.22} Dalam Latihan~\ref{exer 11.3.19}, waktu pelayanan~$S$
mempunyai distribusi geometrik dengan $E(S) = 1/r$. Sebagai gantinya, andaikan
waktu pelayanan adalah waktu konstan $t$~detik. Ubahlah program komputer Anda
dari Latihan~\ref{exer 11.3.20} agar menyimulasikan distribusi waktu pelayanan
konstan. Bandingkan panjang antrean rata-rata untuk kedua jenis distribusi
ketika keduanya mempunyai waktu pelayanan harapan yang sama (yakni, ambil
$t = 1/r$). Distribusi mana yang rata-rata menghasilkan antrean lebih panjang?

\i\label{exer 11.3.23} Suatu percobaan diyakini dapat dijelaskan oleh rantai
Markov dua-keadaan dengan matriks transisi $\mat{P}$, di mana
$$
\mat {P} = \pmatrix{ .5 & .5 \cr p & 1 - p}
$$
dan parameter~$p$ tidak diketahui. Ketika percobaan dilakukan berkali-kali,
rantai berakhir dalam keadaan satu sekitar 20~persen dari waktu dan dalam
keadaan dua sekitar 80~persen dari waktu. Hitunglah taksiran yang masuk akal
untuk parameter~$p$ yang tidak diketahui dan jelaskan cara Anda menemukannya.

\i\label{exer 11.3.24} Buktikan bahwa dalam suatu rantai ergodik dengan $r$
keadaan, setiap keadaan dapat dicapai dari setiap keadaan lain dalam paling
banyak $r - 1$ langkah.

\i\label{exer 11.3.25} Misalkan $\mat{P}$ adalah matriks transisi suatu rantai
ergodik dengan $r$ keadaan. Buktikan bahwa jika entri diagonal $p_{ii}$ positif,
maka rantai tersebut reguler.

\i\label{exer 11.3.26} Buktikan bahwa jika $\mat{P}$ adalah matriks transisi
suatu rantai ergodik, maka $(1/2)(\mat {I} + \mat {P})$ adalah matriks transisi
suatu rantai reguler. \emx {Petunjuk}: Gunakan Latihan~\ref{exer 11.3.25}.

\i\label{exer 11.3.27} Buktikan bahwa $\mat{P}$ dan
$(1/2)(\mat {I} + \mat {P})$ mempunyai vektor-vektor tetap yang sama.

\i\label{exer 11.3.28} Dalam bukunya, \emx {Wahrscheinlichkeitsrechnung und
Statistik,}\footnote{A. Engle, \emx {Wahrscheinlichkeitsrechnung und
Statistik,} vol.~2 
(Stuttgart: Klett Verlag, 1976).}\index{ENGLE, A.} A.~Engle mengusulkan
algoritma untuk menemukan vektor tetap suatu rantai Markov ergodik ketika
peluang transisinya merupakan bilangan rasional. Berikut algoritmanya: untuk
setiap keadaan~$i$, misalkan $a_i$ adalah kelipatan persekutuan terkecil dari
penyebut entri-entri tak nol pada baris ke-$i$. Engle menjelaskan algoritmanya
dalam bentuk pemindahan keping di antara keadaan-keadaan---bahkan, untuk contoh
kecil, ia menyarankan agar algoritma diterapkan dengan cara ini. Mulailah dengan
menempatkan $a_i$ keping pada keadaan~$i$ untuk setiap~$i$. Kemudian, pada
setiap keadaan, distribusikan kembali $a_i$ keping itu dengan mengirimkan
$a_i p_{ij}$ ke keadaan~$j$. Banyaknya keping pada keadaan~$i$ setelah
redistribusi ini tidak harus merupakan kelipatan~$a_i$. Untuk setiap
keadaan~$i$, tambahkan keping secukupnya agar banyaknya keping pada keadaan~$i$
menjadi kelipatan~$a_i$. Kemudian distribusikan kembali keping-keping dengan
cara yang sama. Proses ini pada akhirnya akan mencapai titik ketika banyaknya
keping pada setiap keadaan setelah redistribusi sama seperti sebelum
redistribusi. Pada titik ini, kita telah menemukan vektor tetap. Berikut sebuah
contoh:
$$
\mat {P} = \bordermatrix{
& 1 & 2 & 3 \cr
1 & 1/2 & 1/4 & 1/4 \cr
2 & 1/2 & 0 & 1/2 \cr
3 & 1/2 & 1/4 & 1/4}\ .
$$
Kita mulai dengan $\mat {a} = (4,2,4)$. Keping-keping setelah redistribusi
berturut-turut ditunjukkan dalam Tabel~\ref{table 11.4}.

\begin{table}
\centering
$$\begin{array}{lll}
(4 & 2\;\; & 4)\\
(5 & 2 & 3)\\
(8 & 2 & 4)\\
(7 & 3 & 4)\\
(8 & 4 & 4)\\
(8 & 3 & 5)\\
(8 & 4 & 8)\\
(10 & 4 & 6)\\
(12 & 4 & 8)\\
(12 & 5 & 7)\\
(12 & 6 & 8)\\
(13 & 5 & 8)\\
(16 & 6 & 8)\\
(15 & 6 & 9)\\
(16 & 6 & 12)\\
(17 & 7 & 10)\\
(20 & 8 & 12)\\
(20 & 8 & 12)\ .
\end{array}
$$
\caption{Distribusi keping.}
\label{table 11.4}
\end{table}
Kita memperoleh bahwa $\mat {a} = (20,8,12)$ adalah vektor tetap.

\begin{enumerate}
\item Tulislah program komputer untuk menerapkan algoritma ini.

\item Buktikan bahwa algoritma akan berhenti. \emx {Petunjuk}: Misalkan
$\mat{b}$ adalah vektor berkomponen bilangan bulat yang merupakan vektor tetap
untuk~$\mat{P}$ dan sedemikian rupa sehingga setiap koordinat vektor awal
$\mat{a}$ kurang dari atau sama dengan komponen~$\mat{b}$ yang bersesuaian.
Tunjukkan bahwa dalam iterasi, komponen-komponen vektor selalu bertambah dan
selalu kurang dari atau sama dengan komponen~$\mat{b}$ yang bersesuaian.
\end{enumerate}

\i\label{exer 11.3.29} (Coffman, Kaduta, dan Shepp\footnote{E.~G. Coffman, 
J.~T. Kaduta, dan L.~A. Shepp, ``On the Asymptotic Optimality of
First-Storage
Allocation," \emx {IEEE Trans.\ Software Engineering,} vol.~II (1985),
pp.~235-239.}) Sebuah pusat komputasi menyimpan informasi pada pita dalam
posisi-posisi sepanjang satu satuan. Selama setiap satuan waktu, terdapat satu
permintaan untuk menempati satu satuan pita. Ketika permintaan ini tiba, satuan
bebas pertama digunakan. Selain itu, selama setiap detik, setiap satuan yang
terisi menjadi kosong dengan peluang~$p$. Simulasikan proses ini dengan memulai
dari pita kosong. Taksirlah banyaknya posisi terisi yang diharapkan untuk suatu
nilai~$p$. Jika $p$ kecil, dapatkah Anda memilih pita yang cukup panjang
sehingga peluang pekerjaan baru harus ditolak (yakni, semua posisi terisi)
kecil? Bentuklah rantai Markov dengan keadaan berupa banyaknya posisi terisi.
Ubahlah program {\bf FixedVector} untuk menghitung vektor tetap. Gunakan ini
untuk memeriksa dugaan Anda melalui simulasi.

\istar\label{exer 11.3.30} (Bukti alternatif Teorema~\ref{thm 11.3.8})
Misalkan $\mat{P}$ adalah matriks transisi suatu rantai Markov ergodik. Misalkan
$\mat{x}$ adalah sebarang vektor kolom sedemikian rupa sehingga
$\mat{P}
\mat{x} =
\mat{ x}$. Misalkan $M$ adalah nilai maksimum komponen-komponen $\mat{x}$.
Andaikan $x_i = M$. Tunjukkan bahwa jika $p_{ij} > 0$, maka $x_j = M$.
Gunakan ini untuk membuktikan bahwa $\mat{x}$ harus merupakan vektor konstan.

\i\label{exer 11.3.31} Misalkan $\mat{P}$ adalah matriks transisi suatu rantai
Markov ergodik. Misalkan $\mat{w}$ adalah vektor peluang tetap (yakni,
$\mat{w}$ adalah vektor baris dengan $\mat {w}\mat {P} = \mat {w}$).
Tunjukkan bahwa jika $w_i = 0$ dan $p_{ji} > 0$, maka $w_j = 0$. Gunakan ini
untuk menunjukkan bahwa vektor peluang tetap suatu rantai ergodik tidak dapat
mempunyai entri 0.

\i\label{exer 11.3.32} Carilah rantai Markov yang tidak menyerap dan juga tidak
ergodik.

\end{LJSItem}

\section[Teorema Limit Fundamental]{Teorema Limit Fundamental untuk Rantai
\newline Reguler}\label{sec 11.4}

Teorema limit fundamental untuk rantai Markov reguler menyatakan bahwa jika
$\mat{P}$ adalah
matriks transisi reguler, maka
$$
\lim_{n \to \infty} \mat {P}^n = \mat {W}\ ,
$$
di mana $\mat{W}$ adalah matriks yang setiap barisnya sama dengan vektor baris
peluang tetap tunggal
$\mat{w}$ bagi $\mat{P}$. Dalam bagian ini, kita akan memberikan dua
pembuktian yang sangat berbeda bagi
teorema tersebut.
\par
Pembuktian pertama dilakukan dengan menunjukkan bahwa, untuk setiap vektor kolom
$\mat{y}$,
$\mat{P}^n \mat {y}$ mendekati suatu vektor konstan. Sebagaimana ditunjukkan dalam
Bagian~\ref{sec 11.3}, hal ini
akan memperlihatkan bahwa $\mat{P}^n$ konvergen menuju matriks dengan
kolom-kolom konstan atau, secara ekuivalen, menuju
matriks yang semua barisnya sama.
\par
Lemma berikut menyatakan bahwa jika suatu matriks transisi berukuran $r$-kali-$r$
tidak mempunyai entri nol,
dan $\mat {y}$ adalah sembarang vektor kolom dengan $r$ entri, maka vektor $\mat
{P}\mat{y}$ mempunyai
entri-entri yang ``lebih dekat satu sama lain" daripada entri-entri dalam $\mat {y}$.

\begin{lemma}
Misalkan $\mat{P}$ adalah matriks transisi berukuran $r$-kali-$r$ tanpa entri nol.
Misalkan $d$ adalah
entri terkecil matriks tersebut. Misalkan $\mat{y}$ adalah vektor kolom dengan
$r$ komponen, dengan komponen terbesar $M_0$ dan komponen terkecil $m_0$.
Misalkan $M_1$~dan~$m_1$ masing-masing adalah komponen terbesar dan terkecil dari
vektor $\mat {P} \mat {y}$. Maka
$$
M_1 - m_1 \leq (1 - 2d)(M_0 - m_0)\ .
$$

\proof
Dalam pembahasan setelah Teorema~\ref{thm 11.3.6}, telah dicatat bahwa setiap
entri dalam vektor
$\mat {P}\mat{y}$ merupakan rata-rata berbobot dari entri-entri dalam $\mat {y}$.
Rata-rata berbobot terbesar
yang dapat diperoleh dalam kasus ini akan terjadi jika semua entri
$\mat {y}$ kecuali satu
bernilai $M_0$ dan satu entri bernilai $m_0$, serta entri kecil tersebut
diberi bobot sekecil mungkin, yaitu $d$. Dalam kasus ini, rata-rata berbobot
tersebut sama dengan
$$dm_0 + (1-d)M_0\ .$$
Demikian pula, rata-rata berbobot terkecil yang mungkin adalah
$$dM_0 + (1-d)m_0\ .$$
Dengan demikian,
\begin{eqnarray*}
M_1 - m_1 &\le& \Bigl(dm_0 + (1-d)M_0\Bigr) - \Bigl(dM_0 + (1-d)m_0\Bigr) \\
          &=& (1 - 2d)(M_0 - m_0)\ .
\end{eqnarray*}
Hal ini menyelesaikan pembuktian lemma.
\end{lemma}

Sekarang kita beralih ke pembuktian teorema limit fundamental untuk rantai
Markov reguler.

\begin{theorem}{\bf (Teorema Limit Fundamental untuk Rantai Reguler)}
\index{Teorema Limit Fundamental untuk Rantai Markov Reguler}\index{rantai
Markov!Teorema Limit Fundamental untuk Rantai Reguler}
Jika $\mat{P}$ adalah matriks transisi bagi suatu rantai Markov reguler, maka
$$
\lim_{n \to \infty} \mat {P}^n = \mat {W}\ ,
$$
di mana $\mat {W}$ adalah matriks yang semua barisnya sama. Lebih lanjut, semua
entri dalam $\mat{W}$
bernilai positif.

\proof
Kita membuktikan teorema ini untuk kasus khusus ketika $\mat{P}$ tidak mempunyai
entri~0. Perluasannya
ke kasus umum ditunjukkan dalam Latihan~\ref{exer 11.4.6}. Misalkan \mat {y}
adalah sembarang
vektor kolom dengan $r$ komponen, di mana $r$ adalah banyaknya keadaan dalam
rantai tersebut. Kita mengasumsikan
$r > 1$, sebab jika tidak, teorema ini sepele. Misalkan
$M_n$ dan
$m_n$ masing-masing adalah komponen maksimum dan minimum dari
vektor~$\mat {P}^n \mat { y}$. Vektor~$\mat {P}^n \mat {y}$ diperoleh dari
vektor~$\mat {P}^{n - 1}
\mat {y}$ dengan mengalikannya dari kiri dengan matriks~$\mat{P}$. Oleh karena
itu, setiap komponen
$\mat {P}^n \mat {y}$ merupakan rata-rata dari komponen-komponen $\mat {P}^{n - 1}
\mat {y}$. Dengan demikian,
$$
M_0 \geq M_1 \geq M_2 \geq\cdots
$$
dan
$$
m_0 \leq m_1 \leq m_2 \leq\cdots\ .
$$
Setiap barisan bersifat monoton dan terbatas:
$$
m_0 \leq m_n \leq M_n \leq M_0\ .
$$
Oleh karena itu, masing-masing barisan tersebut mempunyai limit ketika $n$
menuju tak hingga.
\par
Misalkan $M$ adalah limit dari~$M_n$ dan $m$ adalah limit dari~$m_n$. Kita tahu
bahwa $m \leq M$. Kita akan membuktikan bahwa $M - m = 0$. Hal ini akan terbukti jika
$M_n - m_n$ menuju~0. Misalkan $d$ adalah elemen terkecil dari~$\mat{P}$.
Karena semua entri $\mat{P}$ bernilai positif, kita mempunyai $d > 0$. Menurut
lemma di atas,
$$
M_n - m_n \leq (1 - 2d)(M_{n - 1} - m_{n - 1})\ .
$$
Dari sini kita melihat bahwa
$$
M_n - m_n \leq (1 - 2d)^n(M_0 - m_0)\ .
$$
Karena $r \ge 2$, haruslah $d \leq 1/2$, sehingga $0 \leq 1 - 2d < 1$; dengan
demikian, selisih $M_n - m_n$ menuju~0 ketika $n$ menuju tak hingga. Karena setiap
komponen dari~$\mat {P}^n \mat {y}$ terletak di antara $m_n$~dan~$M_n$, setiap
komponen harus
mendekati bilangan yang sama, yaitu $u = M = m$.
Hal ini menunjukkan bahwa
$$
\lim_{n \to \infty} \mat {P}^n \mat {y} = \mat{u}\ ,  \label{eq 11.4.4}
$$
di mana $\mat{u}$ adalah vektor kolom yang semua komponennya sama dengan $u$.
\par
Sekarang, misalkan $\mat{y}$ adalah vektor yang komponen ke-$j$-nya sama dengan~1
dan semua komponen lainnya sama dengan~0. Maka $\mat {P}^n \mat {y}$ adalah
kolom ke-$j$ dari~$\mat {P}^n$. Dengan melakukan ini untuk setiap~$j$, terbukti
bahwa kolom-kolom dari~$\mat {P}^n$ mendekati vektor-vektor kolom konstan.
Artinya, baris-baris dari~$\mat {P}^n$ mendekati suatu vektor baris yang sama,
yaitu~$\mat{w}$, atau
$$
\lim_{n \to \infty} \mat {P}^n = \mat {W}\ .  \label{eq 11.4.5}
$$
\par
Tinggal ditunjukkan bahwa semua entri dalam $\mat{W}$ bernilai positif. Seperti
sebelumnya, misalkan
$\mat{y}$ adalah vektor yang komponen ke-$j$-nya sama dengan~1 dan semua komponen
lainnya sama dengan~0.
Maka $\mat{P}\mat{y}$ adalah kolom ke-$j$ dari $\mat{P}$, dan semua entri kolom
ini bernilai positif. Komponen minimum dari vektor $\mat{P}\mat{y}$ didefinisikan
sebagai $m_1$, sehingga
$m_1 > 0$. Karena $m_1 \le m$, kita mempunyai $m > 0$. Terakhir, perhatikan
bahwa nilai $m$ ini tepat merupakan
komponen ke-$j$ dari $\mat{w}$, sehingga semua komponen $\mat{w}$ bernilai
positif.
\end{theorem}

\subsection*{Bukti Doeblin}
Sekarang kita memberikan pembuktian yang sangat berbeda bagi bagian utama
teorema limit fundamental untuk
rantai Markov reguler. Pembuktian ini pertama kali diberikan oleh
Doeblin,\index{DOEBLIN,
W.}\footnote{W. Doeblin, ``Expos\'e de la Th\'eorie des Chaines Simple Constantes
de Markov \`a un
Nombre Fini d'Etats," \emx {Rev.\ Mach.\ de l'Union Interbalkanique,} vol.~2
(1937), pp.~77--105.}, seorang matematikawan muda cemerlang yang tewas
ketika masih berusia dua puluhan dalam Perang Dunia Kedua.

\begin{theorem}\label{thm 11.4.1}
Misalkan $\mat {P}$ adalah matriks transisi bagi suatu rantai Markov reguler
dengan vektor tetap
$\mat {w}$. Maka, untuk setiap vektor peluang awal $\mat {u}$,
$\mat {uP}^n \rightarrow \mat {w}$ ketika $n \rightarrow \infty.$

\proof Misalkan $ X_0,\ X_1,\ \ldots$ adalah rantai Markov
dengan matriks transisi $\mat {P}$ yang dimulai pada keadaan $s_i$. Misalkan
$Y_0,\ Y_1,\ \ldots$ adalah rantai Markov
dengan matriks transisi $\mat {P}$ yang dimulai dengan peluang-peluang awal
yang diberikan oleh
$\mat {w}$. Proses $X$ dan $Y$ dijalankan secara saling bebas.
\par
Kita juga mempertimbangkan rantai Markov ketiga, $\mat{P}^*$, yang mengamati
kedua proses
$X$ dan $Y$. Keadaan-keadaan bagi $\mat{P}^*$
adalah pasangan $(s_i, s_j)$. Peluang transisinya diberikan oleh
$$\mat{P}^{*}[(i,j),(k,l)] = \mat{P}(i,k) \cdot  \mat{P}(j,l)\ .$$   
Karena $\mat{P}$ reguler, terdapat $N$ sedemikian rupa sehingga
$\mat{P}^{N}(i,j) > 0$ untuk semua $i$
dan
$j$. Dengan demikian, untuk rantai
$\mat{P}^*$ juga dimungkinkan berpindah dari sembarang keadaan $(s_i, s_j)$ ke
sembarang keadaan lain
$(s_k,s_l)$ dalam paling banyak $N$ langkah. Artinya, $\mat{P}^*$ juga merupakan
rantai Markov reguler.
\par
Kita tahu bahwa rantai Markov reguler akan mencapai sembarang keadaan dalam
waktu hingga. Misalkan $T$ adalah
saat pertama rantai $\mat{P}^*$ berada dalam keadaan berbentuk $(s_k,s_k)$.
Dengan kata lain,
$T$ adalah saat pertama proses $X$ dan $Y$ berada dalam keadaan yang sama.
Dengan demikian, kita telah menunjukkan bahwa
$$
P[T > n] \rightarrow 0 \;\;\mbox{ketika}\;\; n \rightarrow \infty\ .
$$
Jika kita mengamati proses $X$ dan $Y$ setelah keduanya untuk pertama kali
berada dalam keadaan yang sama, kita tidak akan memperkirakan adanya perbedaan
dalam perilaku jangka panjang mereka. Karena hal ini akan
terjadi bagaimanapun kedua proses tersebut dimulai, tampak jelas bahwa perilaku
jangka panjang
seharusnya tidak bergantung pada keadaan awal. Sekarang kita menunjukkan bahwa
hal ini benar.
\par
Pertama, kita mencatat bahwa jika $n \ge T$, maka karena $X$ dan $Y$ keduanya
berada dalam keadaan yang sama pada saat $T$,
$$
P(X_n = j\ |\ n \ge T) = P(Y_n = j\ |\ n \ge T)\ .
$$
Jika kedua ruas persamaan ini dikalikan dengan $P(n \ge T)$, kita memperoleh
\begin{equation}
P(X_n = j,\ n \ge T) = P(Y_n = j,\ n \ge T)\ .
\label{eq 11.4.1}
\end{equation}
Kita tahu bahwa untuk semua $n$,
$$P(Y_n = j) = w_j\ .$$
Namun,
$$P(Y_n = j) = P(Y_n = j,\ n \ge T) + P(Y_n = j,\ n < T)\ ,$$
dan suku kedua pada ruas kanan persamaan ini menuju~0 ketika $n$ menuju
$\infty$, sebab $P(n < T)$ menuju~0 ketika $n$ menuju $\infty$. Jadi,
$$P(Y_n = j,\ n \ge T) \rightarrow w_j\ ,$$
ketika $n$ menuju $\infty$. Dari Persamaan~\ref{eq 11.4.1}, kita melihat bahwa
$$P(X_n = j,\ n \ge T) \rightarrow w_j\ ,$$
ketika $n$ menuju $\infty$. Namun, dengan penalaran yang serupa dengan yang
digunakan di atas, selisih
antara pernyataan terakhir ini dan $P(X_n = j)$ menuju~0 ketika $n$ menuju
$\infty$. Oleh karena itu,
$$P(X_n = j) \rightarrow w_j\ ,$$
ketika $n$ menuju $\infty$. Hal ini menyelesaikan pembuktian.
\end{theorem}
Dalam pembuktian di atas, kita belum mengatakan apa pun mengenai laju konvergensi
distribusi $X_n$
menuju distribusi tetap $\mat {w}$. Sebenarnya, dapat ditunjukkan
bahwa\footnote{T. Lindvall, \emx {Lectures on the Coupling Method}   (New York:
Wiley 1992).}
$$
\sum ^{r}_{j = 1} \mid P(X_{n} = j) - w_j \mid \leq 2 
P(T > n)\ .
$$
Ruas kiri pertidaksamaan ini dapat dipandang sebagai jarak antara distribusi
rantai Markov setelah $n$ langkah, yang dimulai pada keadaan $s_i$, dan distribusi
limit $\mat {w}$.

\exercises
\begin{LJSItem}

\i\label{exer 11.4.1} Definisikan $\mat{P}$ dan $\mat{y}$ dengan
$$
\mat {P} = \pmatrix{ .5 & .5 \cr.25 & .75 }, \qquad
\mat {y} = \pmatrix{ 1 \cr 0 }\ .
$$
Hitunglah $\mat {P}\mat{y}$, $\mat {P}^2 \mat {y}$, dan $\mat {P}^4 \mat {y}$,
lalu tunjukkan bahwa hasil-hasilnya mendekati suatu vektor konstan. Vektor apakah
itu?

\i\label{exer 11.4.2} Misalkan $\mat{P}$ adalah matriks transisi reguler berukuran
$r \times r$ dan $\mat{y}$ adalah sebarang vektor kolom dengan $r$ komponen.
Tunjukkan bahwa nilai vektor konstan yang menjadi limit bagi~$\mat {P}^n \mat {y}$
adalah~$\mat{w}\mat{y}$.

\i\label{exer 11.4.3} Misalkan
$$
\mat {P} = \pmatrix{ 1 & 0 & 0 \cr .25 & 0 & .75 \cr 0 & 0 & 1 }
$$
adalah matriks transisi suatu rantai Markov. Carilah dua vektor tetap dari
$\mat {P}$ yang bebas linear. Apakah hal ini menunjukkan bahwa rantai Markov
tersebut tidak reguler?

\i\label{exer 11.4.4} Uraikan himpunan semua vektor kolom tetap bagi rantai yang
diberikan dalam Latihan~\ref{exer 11.4.3}.

\i\label{exer 11.4.6} Teorema bahwa $\mat {P}^n \to \mat {W}$ hanya dibuktikan
untuk kasus ketika $\mat{P}$ tidak mempunyai entri nol. Lengkapilah perincian
perluasan berikut untuk kasus ketika $\mat{P}$ reguler. Karena $\mat{P}$ reguler,
untuk suatu $N, \mat {P}^N$ tidak mempunyai entri nol. Jadi, pembuktian yang
diberikan menunjukkan bahwa $M_{nN} - m_{nN}$ mendekati 0 ketika $n$ menuju tak
hingga. Akan tetapi, selisih $M_n - m_n$ tidak pernah dapat bertambah. (Mengapa?)
Oleh karena itu, jika kita tahu bahwa selisih-selisih yang diperoleh dengan
mengamati setiap kelipatan~$N$ menuju~0, maka seluruh barisan juga harus menuju~0.

\i\label{exer 11.4.7} Misalkan $\mat{P}$ adalah matriks transisi reguler dan
misalkan $\mat{w}$ adalah vektor tetap taknol tunggal dari $\mat{P}$. Tunjukkan
bahwa tidak ada entri $\mat{w}$ yang bernilai 0.

\i\label{exer 11.4.8} Berikut ini sebuah trik untuk dicoba pada teman-teman Anda.
Kocok satu pak kartu dan keluarkan kartu-kartu itu satu per satu. Berilah setiap
kartu bergambar nilai sepuluh. Mintalah teman Anda melihat salah satu dari
sepuluh kartu pertama; jika kartu ini bernilai enam, ia harus melihat kartu yang
muncul enam posisi kemudian; jika kartu ini bernilai tiga, ia harus melihat kartu
yang muncul tiga posisi kemudian, dan seterusnya. Pada akhirnya ia akan mencapai
suatu titik ketika ia harus melihat kartu yang muncul $x$~kartu kemudian, tetapi
kartu yang tersisa kurang dari $x$. Anda lalu menyebutkan kartu terakhir yang ia
lihat, meskipun Anda tidak mengetahui titik awalnya. Katakan kepadanya bahwa Anda
melakukan hal ini dengan mengawasinya, dan ia tidak dapat menyembunyikan kapan ia
melihat kartu-kartu itu. Sebenarnya, Anda hanya melakukan prosedur yang sama dan,
meskipun Anda tidak memulai dari titik yang sama dengannya, kemungkinan besar
Anda akan berakhir pada titik yang sama. Mengapa?

\i\label{exer 11.4.9} Tulislah program untuk memainkan permainan dalam Latihan~\ref{exer
11.4.8}.

\i\label{exer 11.4.10} (Disarankan oleh Peter Doyle)\index{DOYLE, P. G.} Dalam
pembuktian Teorema~\ref{thm 11.4.1}, kita mengasumsikan keberadaan suatu vektor
tetap $\mat{w}$. Untuk menghindari asumsi ini, perkuatlah argumen penggandengan
untuk menunjukkan (tanpa mengasumsikan keberadaan distribusi stasioner
$\mat{w}$) bahwa, untuk konstanta-konstanta $C$ dan $r<1$ yang sesuai, jarak antara
$\alpha P^n$ dan $\beta P^n$ tidak lebih dari $C r^n$ bagi sebarang distribusi awal
$\alpha$ dan $\beta$. Terapkan hal ini pada kasus $\beta = \alpha P$ untuk
menyimpulkan bahwa barisan $\alpha P^n$ adalah barisan Cauchy dan bahwa limitnya
adalah matriks $W$ yang semua barisnya sama dengan suatu vektor peluang $w$
dengan $wP=w$. Perhatikan bahwa jarak antara $\alpha P^n$ dan $w$ tidak lebih dari
$C r^n$. Dengan demikian, dengan membebaskan diri dari asumsi mengenai keberadaan
vektor tetap, kita telah membuktikan bahwa konvergensi menuju kesetimbangan
berlangsung dengan laju eksponensial.
\end{LJSItem}

\section[Waktu Lintasan Pertama Rata-Rata]{Waktu Lintasan Pertama Rata-Rata
untuk Rantai Ergodik}\label{sec 11.5}
Dalam bagian ini kita meninjau dua besaran deskriptif yang berkaitan erat dan
penting bagi rantai ergodik: waktu rata-rata untuk kembali ke suatu keadaan dan
waktu rata-rata untuk berpindah dari satu keadaan ke keadaan lain.
\par
Misalkan $\mat{P}$ adalah matriks transisi suatu rantai ergodik dengan keadaan
$s_1$,~$s_2$, \ldots,~$s_r$. Misalkan $\mat {w} = (w_1,w_2,\ldots,w_r)$ adalah
vektor peluang tetap tunggal sedemikian rupa sehingga
$\mat {w} \mat {P} = \mat {w}$.
Menurut Hukum Bilangan Besar untuk rantai Markov, dalam jangka panjang proses
tersebut akan menghabiskan fraksi waktu sebesar~$w_j$ dalam keadaan~$s_j$. Jadi,
jika kita memulai dari keadaan mana pun, rantai itu pada akhirnya akan mencapai
keadaan~$s_j$; bahkan, rantai akan berada dalam keadaan~$s_j$ tak berhingga kali.
\par
Cara lain untuk melihat hal ini adalah sebagai berikut. Bentuklah rantai Markov
baru dengan menjadikan $s_j$ keadaan menyerap, yaitu tetapkan $p_{jj} = 1$. Jika
kita memulai dari keadaan mana pun selain $s_j$, proses baru ini akan berperilaku
persis seperti rantai semula sampai keadaan $s_j$ dicapai untuk pertama kalinya.
Karena rantai semula ergodik, $s_j$ dapat dicapai dari setiap keadaan lain. Jadi,
rantai baru tersebut merupakan rantai menyerap dengan satu keadaan menyerap
$s_j$, yang pada akhirnya akan dicapai. Dengan demikian, jika kita memulai rantai
semula dari keadaan $s_i$ dengan $i \ne j$, pada akhirnya kita akan mencapai
keadaan $s_j$.
\par
Misalkan $\mat{N}$ adalah matriks fundamental bagi rantai baru. Entri-entri
$\mat{N}$ memberikan banyaknya kunjungan yang diharapkan ke setiap keadaan
sebelum penyerapan. Dalam rantai semula, besaran-besaran ini memberikan banyaknya
kunjungan yang diharapkan ke setiap keadaan sebelum mencapai keadaan~$s_j$ untuk
pertama kalinya. Komponen ke-$i$ dari vektor~$\mat {N} \mat {c}$ memberikan
banyaknya langkah yang diharapkan sebelum penyerapan dalam rantai baru jika
dimulai dari keadaan~$s_i$. Dalam rantai lama, ini adalah banyaknya langkah yang
diharapkan untuk mencapai keadaan~$s_j$ untuk pertama kalinya jika dimulai dari
keadaan~$s_i$.

\subsection*{Waktu Lintasan Pertama Rata-Rata}
\begin{definition}
Jika suatu rantai Markov ergodik dimulai dalam keadaan~$s_i$, banyaknya langkah
yang diharapkan untuk mencapai keadaan~$s_j$ untuk pertama kalinya disebut
\emx {waktu lintasan pertama rata-rata}\index{waktu lintasan pertama rata-rata}
dari~$s_i$ ke~$s_j$. Besaran ini dinotasikan dengan $m_{ij}$. Berdasarkan
konvensi, $m_{ii} = 0$.
\end{definition}
\begin{example}\label{exam 11.5.1}
Mari kembali ke contoh labirin (Contoh~\ref{exam
11.3.3}).\index{tikus}\index{labirin} Rantai ergodik ini kita ubah menjadi rantai
menyerap dengan menjadikan keadaan~5 sebagai keadaan menyerap. Sebagai contoh,
kita dapat menganggap bahwa makanan diletakkan di tengah labirin dan, setelah
tikus menemukan makanan itu, tikus tetap di sana untuk menikmatinya (lihat
Gambar~\ref{fig 11.5}).
\putfig{2truein}{PSfig11-5}{Masalah labirin.}{fig 11.5} 
%\begin{figure}
%\centerline{\epsfxsize=2truein 
%\epsffile{PSfig11-5}}
%\caption{Masalah labirin.}\label{fig 11.5}
%\end{figure}
\par
Matriks transisi baru dalam bentuk kanonik adalah
$$
\mat {P}\;= \bordermatrix{&                 
\hbox{1}&\hbox{2}&\hbox{3}&\hbox{4}&\hbox{6}&\hbox{7}&\hbox{8}
&\hbox{9}&\omit\hfil&\hbox{5}\cr
\hbox{1}\strut   &  0     &1/2     &  0     &   0    & 1/2    &  0     & 0      
&    0   &\srule    & 0 \cr
\hbox{2}\strut   &1/3     &  0     &1/3     &   0    &  0     &  0     & 0     
&    0   &\srule    &1/3\cr
\hbox{3}\strut   &  0     &1/2     &  0     & 1/2    &  0     &  0     & 0     
&    0   &\srule    & 0 \cr
\hbox{4}\strut   &  0     &  0     &1/3     &   0    &  0     &1/3     & 0     
& 1/3    &\srule    &1/3\cr
\hbox{6}\strut   &1/3     &  0     &  0     &   0    &  0     &  0     & 0     
&   0    &\srule    &1/3\cr
\hbox{7}\strut   &  0     &  0     &  0     &   0    &1/2     &  0     & 1/2   
&   0    &\srule    &0  \cr
\hbox{8}\strut   &  0     &  0     &  0     &   0    &  0     &1/3     & 0     
& 1/3    &\srule    &1/3\cr
\hbox{9}\strut   &  0     &  0     &  0     &1/2     &  0     &  0     & 1/2&  
0        &\srule    & 0 \cr
\middlehrule{8}{1}
\hbox{5}\strut   &  0     &  0     &  0     &   0    & 0      & 0      & 0     
& 0      &\srule   
& 1}\ .
$$
Jika kita menghitung matriks fundamental~$\mat{N}$, kita memperoleh
$$
\mat {N} = \frac18 \pmatrix{
14 & 9 & 4 & 3 & 9 & 4 & 3 & 2 \cr
6 & 14 & 6 & 4 & 4 & 2 & 2 & 2 \cr
4 & 9 & 14 & 9 & 3 & 2 & 3 & 4 \cr
2 & 4 & 6 & 14 & 2 & 2 & 4 & 6 \cr
6 & 4 & 2 & 2 & 14 & 6 & 4 & 2 \cr
4 & 3 & 2 & 3 & 9 & 14 & 9 & 4 \cr
2 & 2 & 2 & 4 & 4 & 6 & 14 & 6 \cr
2 & 3 & 4 & 9 & 3 & 4 & 9 & 14 \cr}\ .
$$
Waktu penyerapan yang diharapkan untuk berbagai keadaan awal diberikan oleh
vektor~$\mat {N} \mat {c}$, dengan
$$
\mat {N} \mat {c} = \pmatrix{ 6 \cr 5 \cr 6 \cr 5 \cr 5 \cr 6 \cr 5 \cr 6 \cr}\
.
$$
Kita melihat bahwa, jika dimulai dari kompartemen~1, rata-rata diperlukan enam
langkah untuk mencapai makanan. Berdasarkan simetri, jelas bahwa kita seharusnya
memperoleh jawaban yang sama jika dimulai dari keadaan~3, 7, atau~9. Jelas pula
bahwa, jika dimulai dari salah satu keadaan tersebut, diperlukan satu langkah
lebih banyak daripada jika dimulai dari 2,~4, 6, atau~8. Sebagian hasil yang
diperoleh dari~$\mat{N}$ tidak begitu jelas. Misalnya, perhatikan bahwa banyaknya
kunjungan yang diharapkan ke keadaan awal adalah 14/8, terlepas dari keadaan
tempat kita memulai.
\end{example}

\subsection*{Waktu Kembali Rata-Rata}
Besaran yang berkaitan erat dengan waktu lintasan pertama rata-rata adalah
\emx{waktu kembali rata-rata,} yang didefinisikan sebagai berikut. Andaikan kita
memulai dalam keadaan~$s_i$; tinjaulah lamanya waktu sebelum kita kembali
ke~$s_i$ untuk pertama kalinya. Jelas bahwa kita pasti kembali, sebab pada
langkah pertama kita tetap di~$s_i$ atau berpindah ke suatu keadaan lain~$s_j$,
dan dari setiap keadaan lain~$s_j$ pada akhirnya kita akan mencapai $s_i$ karena
rantai tersebut ergodik.

\begin{definition}
Jika suatu rantai Markov ergodik dimulai dalam keadaan~$s_i$, banyaknya langkah
yang diharapkan untuk kembali ke~$s_i$ untuk pertama kalinya disebut \emx {waktu
kembali rata-rata}\index{waktu kembali rata-rata} bagi~$s_i$. Besaran ini
dinotasikan dengan~$r_i$.
\end{definition}

Kita perlu mengembangkan beberapa sifat dasar waktu lintasan pertama rata-rata.
Tinjaulah waktu lintasan pertama rata-rata dari $s_i$~ke~$s_j$ dan andaikan
$i \ne j$. Besaran ini dapat dihitung sebagai berikut: ambil banyaknya langkah
yang diharapkan dengan syarat hasil langkah pertama, kalikan dengan peluang
terjadinya hasil tersebut, lalu jumlahkan. Jika langkah pertama menuju~$s_j$,
banyaknya langkah yang diharapkan adalah~1; jika langkah itu menuju suatu
keadaan lain~$s_k$, banyaknya langkah yang diharapkan adalah $m_{kj}$~ditambah~1
untuk langkah yang telah diambil. Jadi,
$$
m_{ij} = p_{ij} + \sum_{k \ne j} p_{ik}(m_{kj} + 1)\ ,
$$
atau, karena $\sum_k p_{ik} = 1$,
\begin{equation}
m_{ij} = 1 + \sum_{k \ne j}p_{ik} m_{kj}\ .\label{eq 11.5.1}
\end{equation}

Demikian pula, jika dimulai dari~$s_i$, diperlukan setidaknya satu langkah untuk
kembali. Dengan meninjau semua langkah pertama yang mungkin, kita memperoleh
\begin{eqnarray}
r_i &=& \sum_k p_{ik}(m_{ki} + 1) \\
    &=& 1 + \sum_k p_{ik} m_{ki}\ .\label{eq 11.5.2}
\end{eqnarray}

\subsection*{Matriks Waktu Lintasan Pertama Rata-Rata dan Matriks Waktu Kembali
Rata-Rata}
Sekarang kita definisikan dua matriks $\mat{M}$~dan~$\mat{D}$. Entri
ke-$ij$, yaitu~$m_{ij}$, dari~$\mat{M}$ adalah waktu lintasan pertama rata-rata
untuk berpindah dari~$s_i$ ke~$s_j$ jika $i \ne j$; entri-entri diagonalnya
adalah 0. Matriks $\mat{M}$ disebut \emx {matriks waktu lintasan pertama
rata-rata.}\index{matriks waktu lintasan pertama rata-rata} Matriks $\mat{D}$
adalah matriks yang semua entrinya 0 kecuali entri-entri diagonal $d_{ii} = r_i$.
Matriks $\mat{D}$ disebut \emx {matriks waktu kembali rata-rata.}
\index{matriks waktu kembali rata-rata}
Misalkan $\mat {C}$ adalah matriks berukuran $r \times r$ yang semua entrinya~1.
Dengan menggunakan Persamaan~\ref{eq 11.5.1} untuk kasus $i \ne j$ dan
Persamaan~\ref{eq 11.5.2} untuk kasus $i = j$, kita memperoleh persamaan matriks
\begin{equation}
\mat{M} = \mat{P} \mat{M} + \mat{C} - \mat{D}\ ,  
\label{eq 11.5.3}
\end{equation}
atau
\begin{equation}
(\mat{I} - \mat{P}) \mat{M} = \mat{C} - \mat{D}\ . 
\label{eq 11.5.4}
\end{equation}
Persamaan~\ref{eq 11.5.4}, bersama $m_{ii} = 0$, mengimplikasikan
Persamaan~\ref{eq 11.5.1} dan~\ref{eq 11.5.2}. Sekarang kita siap membuktikan
teorema dasar pertama.

\begin{theorem}
Untuk suatu rantai Markov ergodik, waktu kembali rata-rata bagi keadaan $s_i$
adalah $r_i = 1/w_i$, dengan $w_i$ merupakan komponen ke-$i$ dari vektor peluang
tetap bagi matriks transisi.
\proof
Dengan mengalikan kedua ruas Persamaan~\ref{eq 11.5.4} dengan~$\mat{w}$ dan
menggunakan kenyataan bahwa
$$
\mat {w}(\mat {I} - \mat {P}) = \mat {0}
$$ 
menghasilkan
$$
\mat {w} \mat {C} - \mat {w} \mat {D} = \mat {0}\ .
$$
Di sini $\mat {w} \mat {C}$ adalah vektor baris yang semua entrinya 1 dan
$\mat {w} \mat {D}$ adalah vektor baris dengan entri ke-$i$ sebesar $w_i r_i$.
Jadi,
$$
(1,1,\ldots,1) = (w_1r_1,w_2r_2,\ldots,w_nr_n)
$$
dan
$$
r_i = 1/w_i\ ,
$$
sebagaimana hendak dibuktikan.
\end{theorem}

\begin{corollary}\label{cor 11.5.17}
Untuk suatu rantai Markov ergodik, komponen-komponen vektor peluang tetap~{\bf
w} bernilai positif.
\proof
Kita tahu bahwa nilai-nilai~$r_i$ berhingga, sehingga $w_i = 1/r_i$ tidak mungkin
bernilai 0.
\end{corollary}

\begin{example}
Dalam Contoh~\ref{exam 11.3.3}, kita memperoleh vektor peluang tetap untuk contoh
labirin sebagai
$$
\mat {w} = \pmatrix{ \frac1{12} & \frac18 & \frac1{12} & \frac18 & \frac16 &
\frac18 & \frac1{12} & \frac18 & \frac1{12}}\ .
$$
Dengan demikian, waktu-waktu kembali rata-rata diberikan oleh kebalikan dari
peluang-peluang ini. Artinya,
$$
\mat {r} = \pmatrix{ 12 & 8 & 12 & 8 & 6 & 8 & 12 & 8 & 12 }\ .
$$
\end{example}

Kembali ke Negeri Oz, kita menemukan bahwa cuaca di Negeri Oz dapat
direpresentasikan oleh rantai Markov dengan keadaan hujan, cerah, dan salju.
Dalam Bagian~\ref{sec 11.3}, kita memperoleh vektor limit $\mat {w} = 
(2/5,1/5,2/5)$. Dari sini kita melihat bahwa banyaknya hari rata-rata antara dua
hari hujan adalah 5/2, antara dua hari cerah adalah 5, dan antara dua hari
bersalju adalah 5/2.

\subsection*{Matriks Fundamental}
Sekarang kita akan mengembangkan matriks fundamental bagi rantai ergodik yang
memainkan peran serupa dengan matriks fundamental $\mat {N} = (\mat {I} - \mat
{Q})^{-1}$ bagi rantai menyerap. Seperti pada rantai menyerap, matriks fundamental
dapat digunakan untuk menemukan sejumlah besaran menarik yang melibatkan rantai
ergodik. Dengan matriks ini, kita akan memberikan metode untuk menghitung
waktu-waktu lintasan pertama rata-rata bagi rantai ergodik yang lebih mudah
digunakan daripada metode di atas. Selain itu, kita akan menyatakan (tetapi tidak
membuktikan) Teorema Limit Pusat untuk Rantai Markov, yang pernyataannya
menggunakan matriks fundamental.
\par
Kita mulai dengan meninjau kasus ketika \mat{P} adalah matriks transisi suatu
rantai Markov reguler. Karena tidak ada keadaan menyerap, kita mungkin tergoda
untuk mencoba $\mat {Z} = (\mat {I} - \mat {P})^{-1}$ sebagai matriks fundamental.
Namun, $\mat {I} - \mat {P}$ tidak mempunyai invers. Untuk melihatnya, ingatlah
bahwa suatu matriks~$\mat{R}$ mempunyai invers jika dan hanya jika
$\mat {R} \mat {x} = \mat {0}$ mengimplikasikan $\mat {x} = \mat {0}$. Akan
tetapi, karena $\mat {P} \mat {c} = \mat {c}$, kita mempunyai
$(\mat {I} - \mat {P})\mat {c} = \mat {0}$, sehingga $\mat {I} - \mat {P}$ tidak
mempunyai invers.
\par
Ingatlah bahwa jika kita mempunyai rantai Markov menyerap dan \mat{Q} adalah
pembatasan matriks transisi pada himpunan keadaan sementara, maka matriks
fundamental \mat {N} dapat ditulis sebagai
$$\mat {N} = \mat {I} + \mat {Q} + \mat {Q}^2 + \cdots\ .$$
Deret pangkat ini konvergen karena $\mat {Q}^n \rightarrow 0$, sehingga deret
tersebut berperilaku seperti deret geometri konvergen.
\par
Gagasan ini mungkin mendorong kita mencari deret serupa bagi rantai reguler.
Karena kita tahu bahwa $\mat {P}^n \rightarrow \mat {W}$, kita dapat meninjau
deret
\begin{equation}
\mat {I} + (\mat {P} -\mat {W}) + (\mat {P}^2 - \mat{W}) + \cdots\ .\label{eq
11.5.8}
\end{equation}
Sekarang kita menggunakan sifat-sifat khusus \mat {P} dan \mat {W} untuk menulis
ulang deret ini. Sifat-sifat khusus tersebut adalah: 1)
$\mat {P}\mat {W} = \mat {W}$ dan 2) $\mat {W}^k = \mat {W}$ untuk semua bilangan
bulat positif $k$. Fakta-fakta ini mudah diperiksa dan diserahkan sebagai latihan
(lihat Latihan~\ref{exer 11.5.28}). Dengan menggunakan fakta-fakta tersebut, kita
melihat bahwa
\begin{eqnarray*}
(\mat {P} - \mat {W})^n &=& \sum_{i = 0}^n (-1)^i{n \choose i}\mat
{P}^{n-i}\mat {W}^i \\
&=& \mat {P}^n + \sum_{i = 1}^n (-1)^i{n \choose i} \mat {W}^i \\
&=& \mat {P}^n + \sum_{i = 1}^n (-1)^i{n \choose i} \mat {W} \\
&=& \mat {P}^n + \Biggl(\sum_{i = 1}^n (-1)^i{n \choose i}\Biggr) \mat {W}\ .
\end{eqnarray*}
Jika kita mengembangkan ekspresi $(1-1)^n$ dengan menggunakan Teorema Binomial,
kita memperoleh ekspresi dalam tanda kurung di atas, kecuali bahwa kita memiliki
satu suku tambahan (yang sama dengan 1). Karena $(1-1)^n = 0$, kita melihat bahwa
ekspresi di atas sama dengan -1. Jadi, kita mempunyai
$$(\mat {P} - \mat {W})^n = \mat {P}^n - \mat {W}\ ,$$
bagi semua $n \ge 1$.
\par
Sekarang kita dapat menulis ulang deret dalam \ref{eq 11.5.8} sebagai
$$\mat {I} + (\mat {P} - \mat {W}) + (\mat {P} - \mat {W})^2 + \cdots\ .$$
Karena suku ke-$n$ dalam deret ini sama dengan $\mat {P}^n - \mat {W}$, suku
ke-$n$ menuju 0 ketika $n$ menuju tak hingga. Hal ini cukup untuk menunjukkan
bahwa deret tersebut konvergen dan jumlahnya merupakan invers dari matriks
$\mat {I} - \mat {P} + \mat {W}$. Invers ini kita sebut \emx {matriks
fundamental} yang berkaitan dengan rantai tersebut, dan kita notasikan dengan
\mat {Z}.\index{matriks fundamental!untuk rantai Markov reguler}
\par
Dalam kasus ketika rantai ergodik tetapi tidak reguler, tidak benar bahwa
$\mat {P}^n \rightarrow \mat {W}$ ketika $n \rightarrow \infty$. Meskipun
demikian, matriks $\mat {I} - \mat {P} + \mat {W}$ tetap mempunyai invers,
sebagaimana akan kita tunjukkan sekarang.

\begin{proposition}
Misalkan \mat {P} adalah matriks transisi suatu rantai ergodik dan \mat {W}
adalah matriks yang semua barisnya merupakan vektor baris peluang tetap bagi
\mat {P}. Maka matriks
$$\mat {I} - \mat {P} + \mat {W}$$
mempunyai invers.

\proof
Misalkan $\mat{x}$ adalah vektor kolom sedemikian rupa sehingga
$$
(\mat {I} - \mat {P} + \mat {W})\mat {x} = \mat {0}\ .
$$
Untuk membuktikan proposisi tersebut, cukup ditunjukkan bahwa \mat {x} harus
merupakan vektor nol. Dengan mengalikan persamaan ini dengan~$\mat{w}$ dan
menggunakan kenyataan bahwa $\mat{w}(\mat{I}
- \mat{
P}) = \mat{0}$ dan
$\mat{w} \mat{W} = \mat {w}$, kita memperoleh
$$
\mat {w}(\mat {I} - \mat {P} + \mat {W})\mat {x} = \mat {w} \mat {x} = \mat
{0}\ .
$$
Oleh karena itu,
$$
(\mat {I} - \mat {P})\mat {x} = \mat {0}\ .
$$

Namun, ini berarti bahwa $\mat {x} = \mat {P} \mat {x}$ adalah vektor kolom tetap
bagi~$\mat{P}$. Menurut Teorema~\ref{thm 11.3.10}, hal ini hanya dapat terjadi
jika $\mat{x}$ adalah vektor konstan. Karena $\mat {w}\mat {x} = 0$ dan
\mat {w} mempunyai entri-entri yang positif secara ketat, kita melihat bahwa
$\mat {x} = \mat {0}$. Hal ini menyelesaikan pembuktian.
\end{proposition}
\par
Seperti dalam kasus reguler, kita menyebut invers dari matriks $\mat {I} - \mat
{P} + \mat {W}$ sebagai \emx {matriks fundamental}\index{matriks
fundamental!untuk rantai Markov ergodik} bagi rantai ergodik dengan matriks
transisi \mat {P}, dan kita menggunakan \mat {Z} untuk menotasikan matriks
fundamental ini.

\begin{example}\label{exam 11.5.2} 
Misalkan $\mat{P}$ adalah matriks transisi bagi cuaca di Negeri Oz. Maka
\begin{eqnarray*}
\mat{I} - \mat{P} + \mat{W} &=& \pmatrix{
1 & 0 & 0\cr
0 & 1 & 0\cr
0 & 0 & 1\cr} - \pmatrix{
1/2 & 1/4 & 1/4\cr
1/2 & 0   & 1/2\cr
1/4 & 1/4 & 1/2\cr} + \pmatrix{
2/5 & 1/5 & 2/5\cr
2/5 & 1/5 & 2/5\cr
2/5 & 1/5 & 2/5\cr} \cr
&=& \pmatrix{
9/10 & -1/20 & 3/20\cr
-1/10 & 6/5   & -1/10\cr
3/20 & -1/20 & 9/10\cr}\ ,\cr
\end{eqnarray*}
sehingga
$$
\mat{Z} = (\mat{I} - \mat{P} + \mat{W})^{-1} =
\pmatrix{                   
86/75 & 1/25 & -14/75\cr
2/25 & 21/25 & 2/25\cr
-14/75 & 1/25 & 86/75\cr}\ .         
$$
\end{example}

\subsection*{Menggunakan Matriks Fundamental untuk Menghitung Matriks Waktu
Lintasan Pertama Rata-Rata}

Kita akan menunjukkan cara memperoleh matriks waktu lintasan pertama rata-rata
\mat{M} dari matriks fundamental \mat{Z} bagi suatu rantai Markov ergodik.
Sebelum menyatakan teorema yang memberikan waktu-waktu lintasan pertama, kita
memerlukan beberapa fakta mengenai \mat{Z}.

\begin{lemma}\label{thm 11.5.18}
Misalkan $\mat{Z} = (\mat{I} - \mat{P} + \mat{W})^{-1}$ dan misalkan $\mat{c}$
adalah vektor kolom yang semua entrinya 1. Maka
$$\mat{Z}\mat{c} = \mat{c}\ ,$$
$$\mat{w}\mat{Z} = \mat{w}\ ,$$
dan
$$\mat{Z}(\mat{I} - \mat{P}) = \mat{I} - \mat{W}\ .$$
\proof
Karena $\mat{P}\mat{c} = \mat{c}$ dan $\mat{W}\mat{c} = \mat{c}$,
$$\mat{c} = (\mat{I} - \mat{P} + \mat{W}) \mat{c}\ .$$
Jika kita mengalikan kedua ruas persamaan ini dari kiri dengan $\mat{Z}$, kita
memperoleh
$$\mat{Z}\mat{c} = \mat{c}\ .$$
\par
Demikian pula, karena $\mat{w}\mat{P} = \mat{w}$ dan
$\mat{w}\mat{W} = \mat{w}$,
$$\mat{w} = \mat{w}(\mat{I} - \mat{P} + \mat{W})\ .$$
Jika kita mengalikan kedua ruas persamaan ini dari kanan dengan $\mat{Z}$, kita
memperoleh
$$\mat{w}\mat{Z} = \mat{w}\ .$$
\par
Terakhir, kita mempunyai
\begin{eqnarray*}
(\mat{I} - \mat{P} + \mat{W})(\mat{I} - \mat{W}) &=& 
\mat{I} - \mat{W} - \mat{P} + \mat{W} + \mat{W} - \mat{W}\\
    &=& \mat{I} - \mat{P}\ .
\end{eqnarray*}
Dengan mengalikan dari kiri dengan \mat{Z}, kita memperoleh
$$\mat{I} - \mat{W} = \mat{Z}(\mat{I} - \mat{P})\ .$$
Hal ini menyelesaikan pembuktian.
\end{lemma}

Teorema berikut menunjukkan cara memperoleh waktu-waktu lintasan pertama
rata-rata dari matriks fundamental.

\begin{theorem}\label{thm 11.5.19}
Matriks waktu lintasan pertama rata-rata~$\mat{M}$ bagi suatu rantai ergodik
ditentukan dari matriks fundamental $\mat{Z}$ dan vektor baris peluang tetap
$\mat{w}$ melalui
$$
m_{ij} = \frac{z_{jj} - z_{ij}}{w_j}\ .
$$
\proof
Dalam Persamaan~\ref{eq 11.5.4}, kita telah menunjukkan bahwa
$$
(\mat {I} - \mat {P})\mat {M} = \mat {C} - \mat {D}\ .
$$
Jadi,
$$
\mat {Z}(\mat {I} - \mat {P})\mat {M} = \mat {Z} \mat {C} - \mat {Z} \mat {D}\
,
$$
dan, menurut Lemma~\ref{thm 11.5.18},
$$
\mat {Z}(\mat {I} - \mat {P})\mat {M} = \mat {C} - \mat {Z} \mat {D}\ .
$$
Dengan kembali menggunakan Lemma~\ref{thm 11.5.18}, kita mempunyai
$$
\mat {M} - \mat{W}\mat {M} = \mat {C} - \mat {Z} \mat {D}
$$
atau
$$
\mat {M} = \mat {C} - \mat {Z} \mat {D} + \mat{W}\mat {M}\ .
$$
Dari persamaan ini, kita melihat bahwa
\begin{equation}
m_{ij} = 1 - z_{ij}r_j + (\mat{w} \mat{M})_j\ .  \label{eq 11.5.6}
\end{equation}
Namun, $m_{jj} = 0$, sehingga
$$
0 = 1 - z_{jj}r_j + (\mat {w} \mat {M})_j\ ,
$$
atau
\begin{equation}
(\mat{w} \mat{M})_j = z_{jj}r_j - 1\ .  \label{eq 11.5.7}
\end{equation}
Dari Persamaan~\ref{eq 11.5.6} dan \ref{eq 11.5.7}, kita memperoleh
$$
m_{ij} = (z_{jj} - z_{ij}) \cdot r_j\ .
$$
Karena $r_j = 1/w_j$,
$$
m_{ij} = \frac{z_{jj} - z_{ij}}{w_j}\ .
$$
\end{theorem}

\begin{example}\label{exam 11.5.3} (Lanjutan Contoh~\ref{exam 11.5.2})
Dalam contoh Negeri Oz, kita memperoleh
$$
\mat{Z} = (\mat{I} - \mat{P} + \mat{W})^{-1} =
\pmatrix{                   
86/75 & 1/25 & -14/75\cr
2/25 & 21/25 & 2/25\cr
-14/75 & 1/25 & 86/75\cr}\ .         
$$
Kita juga telah melihat bahwa $\mat {w} = (2/5,1/5, 2/5)$. Jadi, sebagai contoh,
\begin{eqnarray*}
m_{12} &=& \frac{z_{22} - z_{12}}{w_2} \\
       &=& \frac{21/25 - 1/25}{1/5} \\
       &=& 4\ ,\\
\end{eqnarray*}
menurut Teorema~\ref{thm 11.5.19}. Dengan melakukan perhitungan bagi entri-entri
lain dari $\mat{M}$, kita memperoleh
$$
\mat {M} = \pmatrix{
   0 & 4 & 10/3\cr
 8/3 & 0 & 8/3\cr
10/3 & 4 & 0\cr}\ .
$$
\end{example}

\subsection*{Perhitungan}
Program {\bf ErgodicChain} menghitung matriks fundamental, vektor tetap, matriks
waktu kembali rata-rata~$\mat{D}$, dan matriks waktu lintasan pertama
rata-rata~$\mat{M}$. Kita telah menjalankan program tersebut untuk model guci
Ehrenfest (Contoh~\ref{exam 11.1.6}).\index{model Ehrenfest} \index{difusi
gas!model Ehrenfest}
Kita memperoleh:
$$
\mat {P} = \bordermatrix{
    &  0  & 1 & 2 & 3 & 4 \cr
  0 &   .0000 & 1.0000 &  .0000 &  .0000 &  .0000\cr
  1 &   .2500 & .0000  &  .7500 &  .0000 &  .0000\cr
  2 &   .0000 & .5000  &  .0000 &  .5000 &  .0000\cr
  3 &   .0000 & .0000  &  .7500 &  .0000 &  .2500\cr
  4 &   .0000 & .0000  &  .0000 &  1.0000&  .0000}\ ;
$$
 
$$
\mat {w} = \bordermatrix{
  &     0  &   1   &   2   &   3   &    4\cr
  & .0625  &.2500  &.3750  &.2500  &.0625}\ ;
$$

$$
\mat {r} = \bordermatrix{
  &       0  &      1  &      2 &       3 &       4\cr
  & 16.0000  & 4.0000  & 2.6667 &  4.0000 & 16.0000}\ ;
$$

$$
\mat {M} = \bordermatrix{
  &        0 &      1 &      2 &       3 &       4\cr
0 &    .0000 & 1.0000 & 2.6667 &  6.3333 & 21.3333\cr
1 &  15.0000 &  .0000 & 1.6667 &  5.3333 & 20.3333\cr
2 &  18.6667 & 3.6667 &  .0000 &  3.6667 & 18.6667\cr
3 &  20.3333 & 5.3333 & 1.6667 &   .0000 & 15.0000\cr
4 &  21.3333 & 6.3333 & 2.6667 &  1.0000 &   .0000}\ .
$$ 
\par
Dari matriks waktu lintasan pertama rata-rata, kita melihat bahwa waktu rata-rata
untuk beralih dari 0~bola dalam guci~1 ke 2~bola dalam guci~1 adalah 2.6667
langkah, sedangkan waktu rata-rata untuk beralih dari 2~bola dalam guci~1 ke
0~bola dalam guci~1 adalah 18.6667. Hal ini mencerminkan bahwa model tersebut
menunjukkan kecenderungan ke pusat. Tentu saja, fisikawan tertarik pada kasus
dengan sejumlah besar molekul, atau bola, sehingga kita seharusnya meninjau
contoh ini untuk~$n$ yang sedemikian besar sehingga kita bahkan tidak dapat
menghitungnya dengan komputer.

\subsection*{Model Ehrenfest}
\begin{example}(Lanjutan Contoh~\ref{exam 11.3.4})
Mari meninjau model Ehrenfest (lihat Contoh \ref{exam
11.1.6})\index{model Ehrenfest}
\index{difusi gas!model Ehrenfest} untuk difusi gas dalam kasus umum dengan
$2n$ bola. Setiap detik, salah satu dari $2n$ bola dipilih secara acak dan
dipindahkan dari guci tempat bola itu berada ke guci lainnya. Jika terdapat
$i$~bola dalam guci pertama, maka dengan peluang $i/2n$ kita mengambil salah
satunya dan memasukkannya ke guci kedua, sedangkan dengan peluang $(2n - i)/2n$
kita mengambil sebuah bola dari guci kedua dan memasukkannya ke guci pertama.
Pada setiap detik, banyaknya bola~$i$ dalam guci pertama kita jadikan keadaan
sistem. Maka, dari keadaan~$i$ kita hanya dapat berpindah ke keadaan $i - 1$ dan
$i + 1$, dan peluang transisinya diberikan oleh
$$
p_{ij} = \left \{ \matrix{
               \frac i{2n}\ , & \mbox{jika} \,\,j = i-1, \cr
               1 - \frac i{2n}\ , & \mbox{jika} \,\,j = i+1, \cr
               0\ , & \mbox{selain itu.}\cr}\right.
$$
\par
Hal ini mendefinisikan matriks transisi suatu rantai Markov ergodik yang tidak
reguler (lihat Latihan~\ref{exer 11.3.14}). Di sini fisikawan tertarik pada
prediksi jangka panjang mengenai keadaan yang ditempati. Dalam
Contoh~\ref{exam 11.3.4}, kita memberikan alasan intuitif untuk menduga bahwa
vektor tetap $\mat{w}$ adalah distribusi binomial dengan parameter $2n$ dan
$1/2$. Mudah diperiksa bahwa dugaan ini benar. Jadi,
$$
w_i = \frac{{2n \choose i}}{2^{2n}}\ .
$$
Dengan demikian, waktu kembali rata-rata bagi keadaan~$i$ adalah
$$
r_i = \frac{2^{2n}}{{2n \choose i}}\ .
$$
\par
Secara khusus, tinjaulah suku tengah $i = n$. Kita telah melihat bahwa suku ini
kira-kira sama dengan $1/\sqrt{\pi n}$. Jadi, kita dapat menghampiri $r_n$ dengan
$\sqrt{\pi n}$.
\par
Model ini digunakan untuk menjelaskan konsep reversibilitas dalam sistem fisik.
Andaikan kita membiarkan sistem berjalan sampai mencapai kesetimbangan. Pada
saat itu dibuat sebuah film yang memperlihatkan perkembangan sistem. Film
tersebut kemudian ditayangkan kepada Anda, dan Anda diminta menentukan apakah
film diputar maju atau mundur. Tampaknya selalu ada kecenderungan untuk bergerak
menuju proporsi bola yang sama, sehingga urutan waktu yang benar seharusnya
merupakan urutan dengan transisi terbanyak dari~$i$ ke $i - 1$ jika $i > n$, dan
dari $i$ ke $i + 1$ jika $i < n$.

\putfig{4truein}{PSfig11-6}{Simulasi Ehrenfest.}{fig 11.6}

%\begin{figure}
%\centerline{\epsfxsize=4truein 
%\epsffile{PSfig11-6}}
%\caption{Simulasi Ehrenfest.}\label{fig 11.6}
%\end{figure}
\par
Pada Gambar~\ref{fig 11.6}, kita memperlihatkan hasil simulasi model guci
Ehrenfest untuk kasus $n = 50$ dan 1000 satuan waktu dengan menggunakan program
{\bf EhrenfestUrn}.\index{EhrenfestUrn (program)} Grafik atas memperlihatkan
hasil-hasil tersebut dalam urutan terjadinya, sedangkan grafik bawah
memperlihatkan hasil yang sama dengan urutan waktu dibalik. Tidak tampak adanya
perbedaan.
\par
Perhatikan bahwa jika kita tidak memulai dalam keadaan setimbang, kedua grafik
itu biasanya akan tampak cukup berbeda.
\end{example}

\subsection*{Reversibilitas}
Jika model Ehrenfest dimulai dalam keadaan setimbang, proses tersebut tidak
memperlihatkan arah waktu. Hal ini terjadi karena proses tersebut mempunyai
sifat yang disebut \emx {reversibilitas.}\index{reversibilitas} Definisikan
$X_n$ sebagai banyaknya bola dalam guci kiri pada langkah $n$. Untuk suatu rantai
ergodik umum, kita dapat menghitung peluang transisi balik:
\begin{eqnarray*}
P(X_{n - 1} = j | X_n = i) &=& \frac{P(X_{n - 1} = j,X_n = i)}{P(X_n = i)} \\
     &=& \frac{P(X_{n - 1} = j) P(X_n = i | X_{n - 1} = j)}{P(X_n = i)} \\
     &=& \frac{P(X_{n - 1} = j) p_{ji}}{P(X_n = i)}\ .\\
\end{eqnarray*}
\par
Secara umum, nilai ini bergantung pada $n$, sebab $P(X_n = j)$ dan juga
$P(X_{n - 1} = j)$ berubah bersama~$n$. Akan tetapi, jika kita memulai dengan
vektor~$\mat{w}$ atau menunggu sampai kesetimbangan tercapai, hal ini tidak lagi
berlaku. Kemudian kita dapat mendefinisikan
$$
p_{ij}^* = \frac{w_j p_{ji}}{w_i}
$$
sebagai matriks transisi bagi proses yang diamati dengan arah waktu dibalik.
\par
Mari menghitung suatu peluang transisi khas bagi rantai balik $\mat{
P}^* = \{p_{ij}^*\}$ dalam model Ehrenfest. Sebagai contoh,
\begin{eqnarray*}
 p_{i,i - 1}^* &=& \frac{w_{i - 1} p_{i - 1,i}}{w_i} = 
\frac{{2n \choose {i - 1}}}{2^{2n}}
                  \times \frac{2n - i + 1}{2n} \times 
\frac{2^{2n}}{{2n \choose i}}\\
               &=& \frac{(2n)!}{(i - 1)!\,(2n - i + 1)!} \times 
\frac{(2n - i + 1) i!\,
                  (2n - i)!}{2n(2n)!} \\
               &=& \frac{i}{2n} = p_{i,i - 1}\ .\\
\end{eqnarray*}
\par
Perhitungan serupa untuk peluang-peluang transisi lainnya menunjukkan bahwa
$\mat {P}^* = \mat {P}$. Jika hal ini terjadi, proses tersebut disebut \emx
{reversibel.} Jelas bahwa suatu rantai ergodik reversibel jika dan hanya jika,
untuk setiap pasangan keadaan $s_i$~dan~$s_j$,
$w_i p_{ij} = w_j p_{ji}$. Secara khusus, bagi model Ehrenfest hal ini berarti
bahwa $w_i p_{i,i - 1} = w_{i - 1} p_{i - 1,i}$. Jadi, dalam keadaan setimbang,
pasangan $(i, i - 1)$ dan $(i - 1, i)$ seharusnya muncul dengan frekuensi yang
sama. Walaupun banyak rantai Markov yang muncul dalam penerapan bersifat
reversibel, syarat ini sangat kuat. Dalam Latihan~\ref{exer 11.5.12}, Anda
diminta mencari contoh rantai Markov yang tidak reversibel.
 
\subsection*{Teorema Limit Pusat untuk Rantai Markov}
Andaikan kita mempunyai rantai Markov ergodik dengan keadaan $s_1, s_2, \ldots,
s_k$. Wajar untuk meninjau distribusi peubah acak $S^{(n)}_j$, yang menyatakan
berapa kali rantai berada dalam keadaan $s_j$ selama $n$ langkah pertama. Komponen
ke-$j$, yaitu $w_j$, dari vektor baris peluang tetap $\mat{w}$ adalah proporsi
waktu rantai berada dalam keadaan $s_j$ dalam jangka panjang. Oleh karena itu,
masuk akal untuk menduga bahwa nilai harapan peubah acak $S^{(n)}_j$, ketika
$n \rightarrow \infty$, setara secara asimtotik dengan $nw_j$, dan mudah ditunjukkan bahwa
memang demikian (lihat Latihan~\ref{exer 11.5.29}).
\par
Wajar pula untuk menanyakan apakah peubah-peubah acak $S^{(n)}_j$ mempunyai
distribusi limit. Jawabannya ya; bahkan, distribusi limit ini adalah distribusi
normal. Seperti dalam kasus percobaan-percobaan bebas, peubah-peubah acak ini
harus dinormalkan. Jadi, dari $S^{(n)}_j$ kita harus mengurangkan nilai harapannya,
lalu membaginya dengan simpangan bakunya. Dalam kedua hal tersebut, kita akan
menggunakan nilai asimtotik besaran-besaran ini, bukan nilainya sendiri. Jadi,
dalam hal pertama kita akan menggunakan nilai $nw_j$. Tidak begitu jelas apa
yang harus kita gunakan dalam hal kedua. Ternyata, besaran
\begin{equation}
\sigma_j^2 = 2w_jz_{jj} - w_j - w_j^2
\label{eq 11.5.9}
\end{equation}
menyatakan varians asimtotik. Dengan gagasan-gagasan ini, kita dapat menyatakan
teorema berikut.
\begin{theorem}\label{thm 11.5.20}{\bf (Teorema Limit Pusat untuk Rantai
Markov)}
\index{Teorema Limit Pusat!untuk Rantai Markov}\index{rantai Markov!Teorema
Limit Pusat untuk}
Untuk suatu rantai ergodik dan sebarang bilangan real $r < s$, kita mempunyai
$$P\Biggl(r < {{S^{(n)}_j - nw_j}\over{\sqrt{n\sigma_j^2}}}< s\Biggr)
\rightarrow 
{1\over{\sqrt {2\pi}}}\int_r^s e^{-x^2/2}\,dx\
,$$ ketika $n \rightarrow \infty$, untuk sebarang pilihan keadaan awal, dengan
$\sigma_j^2$ merupakan besaran yang didefinisikan dalam
Persamaan~\ref{eq 11.5.9}.
\end{theorem}


\subsection*{Catatan Sejarah}

Rantai Markov diperkenalkan oleh Andre\u i Andreevich Markov\index{MARKOV, A.
A.} (1856--1922) dan dinamai untuk menghormatinya. Ia adalah mahasiswa berbakat
yang menerima medali emas atas skripsinya di Universitas St.~Petersburg. Selain
aktif sebagai matematikawan peneliti dan pengajar, ia juga aktif dalam politik
dan ikut serta dalam gerakan liberal di Rusia pada awal abad kedua puluh. Pada
1913, ketika pemerintah merayakan ulang tahun ke-300 Wangsa Romanov, Markov
menyelenggarakan perayaan tandingan untuk ulang tahun ke-200 penemuan Hukum
Bilangan Besar oleh Bernoulli.
\par
Markov terdorong mengembangkan rantai Markov sebagai perluasan alami dari
barisan peubah acak bebas. Dalam makalah pertamanya pada 1906, ia membuktikan
bahwa, bagi rantai Markov dengan peluang transisi positif dan keadaan numerik,
rata-rata hasil konvergen menuju nilai harapan distribusi limit (vektor tetap).
Dalam makalah berikutnya, ia membuktikan teorema limit pusat bagi rantai semacam
itu. Menurut ringkasan A.~P. Youschkevitch, Markov memperoleh gagasan rantainya dari
kebutuhan internal teori peluang dan tidak membahas penerapannya pada fisika. Contoh
konkretnya ialah teks sastra dengan dua keadaan, yakni vokal dan konsonan.\footnote{Lihat
\emx {Dictionary of Scientific Biography,} ed.~C.~C. Gillespie (New York: Scribner's
Sons, 1970), pp.~124--130. Ringkasan editorial; ungkapan sumber tidak diterjemahkan
atau direproduksi.}
\par
Dalam sebuah makalah yang ditulis pada 1913,\footnote{A.~A. Markov, ``An Example of Statistical
Analysis of the Text of Eugene Onegin Illustrating the Association of Trials
into a Chain," \emx {Bulletin de l'Acadamie Imperiale des Sciences de
St.~Petersburg,} ser.~6, vol.~7 (1913), pp.~153--162.} Markov memilih barisan
20{,}000 huruf dari \emx {Eugene Onegin} karya Pushkin untuk
melihat apakah barisan ini secara hampiran dapat dianggap sebagai rantai
sederhana. Ia memperoleh rantai Markov dengan matriks transisi
$$
\bordermatrix{
& \mbox{vokal} & \mbox{konsonan} \cr
\mbox{vokal} & .128 & .872 \cr
\mbox{konsonan} & .663 & .337}\ .
$$
\par
Vektor tetap bagi rantai ini adalah $(.432, .568)$, yang menunjukkan bahwa kita
seharusnya mengharapkan sekitar 43.2~persen vokal dan 56.8~persen konsonan dalam
novel tersebut; hal ini ditegaskan oleh penghitungan sebenarnya.
\par
Claude Shannon meninjau perluasan menarik dari gagasan ini dalam bukunya \emx {
The Mathematical Theory of Communication,}\index{SHANNON, C. E.}\index{WEAVER,
W.}\footnote{C.~E.
Shannon and W. Weaver, \emx {The Mathematical Theory of Communication} (Urbana:
Univ.\ of
Illinois Press, 1964).} Dalam buku itu ia mengembangkan konsep entropi dalam
teori informa\-si. Shannon meninjau serangkaian hampiran rantai Markov terhadap
prosa bahasa Inggris. Mula-mula ia melakukannya dengan rantai yang
keadaan-keadaannya berupa huruf, kemudian dengan rantai yang keadaan-keadaannya
berupa kata. Sebagai contoh, untuk kasus kata, pertama-tama ia menyajikan
simulasi dengan kata-kata yang dipilih secara independen, tetapi dengan frekuensi
yang sesuai. Contoh ilustratif mandiri berikut menunjukkan jenis keluaran tersebut;
deret Shannon tidak direproduksi.
\par
\begin{quote}
THE IN FOR A WITH IS OF THAT TO AND BY IT FROM BE ON THIS AS WERE AT OR.
\end{quote}
\par
Kemudian ia mencatat meningkatnya kemiripan dengan teks bahasa Inggris biasa
ketika kata-kata dipilih sebagai rantai Markov. Contoh mandiri berikut menggambarkan
keluaran yang lebih menyerupai bahasa Inggris biasa tanpa mereproduksi deret Shannon.

\begin{quote}
THE HOUSE OF THE OLD ROAD WAS IN THE PLACE WHERE THE STORY COULD BEGIN AGAIN.
\end{quote}
\par
Simulasi seperti simulasi terakhir dilakukan dengan membuka sebuah buku dan
memilih kata pertama, misalkan kata itu \emx {the.} Kemudian buku dibaca sampai
kata \emx {the} muncul lagi, dan kata sesudahnya dipilih sebagai kata kedua;
ternyata kata itu adalah \emx {head.} Buku lalu dibaca sampai kata \emx {head}
muncul lagi, dan kata berikutnya, \emx {and,} dipilih, demikian seterusnya.
\par
Contoh awal lain penggunaan rantai Markov muncul dalam kajian Galton pada 1889
mengenai masalah bertahannya nama keluarga dan dalam rantai Markov yang
diperkenalkan oleh P. dan T. Ehrenfest pada 1907 untuk difusi. Pada 1912,
Poincar\'{e} membahas pengocokan kartu dalam kerangka rantai Markov ergodik yang
didefinisikan pada suatu grup permutasi. Gerak Brown, versi waktu kontinu dari
jalan acak, diperkenalkan pada 1900--1901 oleh L. Bachelier dalam kajiannya
mengenai pasar saham, dan pada 1905--1907 dalam karya A. Einstein dan M.
Smoluchowsky mengenai proses fisik.
\par
Salah satu kajian sistematis pertama mengenai rantai Markov berhingga dilakukan
oleh M. Frechet.\index{FRECHET, M.}\footnote{M. Frechet, ``Th\'eorie des
\'ev\'enements en chaine dans
le cas d'un nombre fini d'\'etats possible," in \emx {Recherches th\'eoriques
Modernes sur le calcul des probabilit\'es,} vol.~2 (Paris, 1938).} Perlakuan
rantai Markov dalam kerangka dua matriks fundamental yang telah kita gunakan
dikembangkan oleh Kemeny dan Snell\index{KEMENY, J. G.}\index{SNELL,
J. L.}
\footnote{J. G. Kemeny and J. L. Snell, \emx {
Finite Markov Chains.}} untuk menghindari penggunaan nilai eigen yang dianggap
terlalu rumit oleh salah seorang penulis tersebut. Matriks fundamental~$\mat{N}$
juga muncul dalam karya J.~L. Doob dan yang lainnya ketika mengkaji hubungan
antara proses Markov dan teori potensial klasik. Matriks fundamental~$\mat{Z}$
bagi rantai ergodik pertama kali muncul dalam karya Frechet, yang menggunakannya
untuk menemukan varians limit bagi teorema limit pusat untuk rantai Markov.

\exercises
\begin{LJSItem}
\i\label{exer 11.5.1} Tinjaulah rantai Markov dengan matriks transisi
$$
\mat {P} = \pmatrix{ 1/2 & 1/2 \cr 1/4 & 3/4}\ .
$$
Carilah matriks fundamental~$\mat{Z}$ bagi rantai ini. Hitunglah matriks waktu
lintasan pertama rata-rata dengan menggunakan $\mat{Z}$.

\i\label{exer 11.5.2} Suatu kajian tentang kekuatan tim sepak bola Amerika Ivy
League menunjukkan bahwa, jika suatu perguruan tinggi mempunyai tim yang kuat
pada suatu tahun, pada tahun berikutnya tim itu sama mungkin menjadi kuat atau
sedang; jika timnya sedang, dalam separuh kasus tim tersebut tetap sedang pada tahun
berikutnya dan, jika berubah, tim itu sama mungkin menjadi kuat atau lemah; jika
timnya lemah, peluangnya untuk tetap lemah adalah 2/3 dan peluangnya untuk
menjadi sedang adalah 1/3.
\begin{enumerate}
\item Suatu perguruan tinggi mempunyai tim yang kuat. Secara rata-rata, berapa
lama waktu berlalu sebelum perguruan tinggi itu kembali mempunyai tim yang kuat?

\item Suatu perguruan tinggi mempunyai tim yang lemah; berapa lama (secara
rata-rata) para alumninya harus menunggu tim yang kuat?
\end{enumerate}

\i\label{exer 11.5.3} Tinjaulah Contoh~\ref{exam 11.1.2} dengan $a = .5$ dan
$b = .75$. Andaikan Presiden menyatakan bahwa ia akan maju. Carilah lamanya waktu
yang diharapkan sampai jawaban tersebut diteruskan secara keliru untuk pertama
kalinya.

\i\label{exer 11.5.4} Carilah waktu kembali rata-rata bagi setiap keadaan dalam
Contoh~\ref{exam 11.1.2} untuk $a = .5$ dan $b = .75$. Lakukan hal yang sama
untuk $a$~dan~$b$ umum.

\i\label{exer 11.5.5} Sebuah dadu dilempar berulang kali. Dengan menggunakan
hasil-hasil dalam bagian ini, tunjukkan bahwa waktu rata-rata antara dua
kemunculan suatu angka tertentu adalah~6.

\i\label{exer 11.5.6} Untuk contoh Negeri Oz (Contoh~\ref{exam 11.1.1}),
jadikan hujan sebagai keadaan menyerap dan carilah matriks
fundamental~$\mat{N}$. Tafsirkan hasil-hasil yang diperoleh dari rantai ini dalam
kaitannya dengan rantai semula.

\i\label{exer 11.5.7} Seekor tikus berlari melalui labirin yang ditunjukkan pada
Gambar~\ref{fig 11.6.5}. Pada setiap langkah, tikus meninggalkan ruangan yang
ditempatinya dengan memilih secara acak salah satu pintu keluar dari ruangan itu.
\putfig{2truein}{PSfig11-6-5}{Labirin untuk Latihan \protect\ref{exer 11.5.7}\protect.}{fig 11.6.5}
%\begin{figure}
%\centerline{\epsfxsize=2truein 
%\epsffile{PSfig11-6.5}}
%\caption{Labirin untuk Latihan \protect\ref{exer 11.5.7}\protect.}\label{fig 11.6.5}
%\end{figure}
\begin{enumerate}

\item Berikan matriks transisi $\mat{P}$ bagi rantai Markov ini.

\item Tunjukkan bahwa rantai ini ergodik, tetapi tidak reguler.

\item Carilah vektor tetapnya.

\item Carilah banyaknya langkah yang diharapkan sebelum mencapai Ruangan~5 untuk
pertama kalinya jika dimulai dari Ruangan~1.
\end{enumerate}

\i\label{exer 11.5.8} Ubahlah program {\bf  ErgodicChain} agar Anda dapat
menghitung besaran-besaran dasar bagi contoh antrean dalam Latihan~\ref{sec
11.3}.\ref{exer 11.3.19}. Tafsirkan waktu kembali rata-rata bagi keadaan~0.

\i\label{exer 11.5.9} Tinjaulah jalan acak pada lingkaran dengan keliling $n$.
Pejalan mengambil satu langkah satuan searah jarum jam dengan peluang~$p$ dan
satu langkah satuan berlawanan arah jarum jam dengan peluang $q = 1 - p$.
Ubahlah program {\bf  ErgodicChain} agar Anda dapat memasukkan $n$~dan~$p$ serta
menghitung besaran-besaran dasar bagi rantai ini.
\begin{enumerate}

\item Untuk nilai $n$ manakah rantai ini reguler? Ergodik?

\item Tentukan vektor limit $\mat{w}$.

\item Carilah matriks waktu lintasan pertama rata-rata untuk $n = 5$ dan
$p = .5$. Periksalah bahwa $m_{ij} = d(n - d)$, dengan $d$ merupakan jarak
searah jarum jam dari $i$~ke~$j$.
\end{enumerate}
 
\i\label{exer 11.5.10} Dua pemain mencocokkan sisi uang logam dan secara
keseluruhan memiliki 5~keping uang satu sen. Jika pada suatu saat salah seorang
pemain memiliki semua keping tersebut, agar permainan terus berlangsung ia
memberikan satu keping kembali kepada pemain lainnya, lalu permainan dilanjutkan.
Tunjukkan bahwa permainan ini dapat dirumuskan sebagai rantai ergodik. Kajilah
rantai ini dengan menggunakan program {\bf  ErgodicChain}.

\i\label{exer 11.5.11} Hitunglah matriks transisi balik bagi contoh Negeri Oz
(Contoh~\ref{exam 11.1.1}). Apakah rantai ini reversibel?

\i\label{exer 11.5.12} Berikan contoh rantai Markov ergodik dengan tiga keadaan
yang tidak reversibel.

\i\label{exer 11.5.13} Misalkan $\mat{P}$ adalah matriks transisi suatu rantai
Markov ergodik dan $\mat{P}^*$ adalah matriks transisi balik. Tunjukkan bahwa
keduanya mempunyai vektor peluang tetap yang sama, yaitu~$\mat{w}$.

\i\label{exer 11.5.14} Jika $\mat{P}$ adalah rantai Markov reversibel, apakah
waktu rata-rata untuk berpindah dari keadaan~$i$ ke keadaan~$j$ selalu sama
dengan waktu rata-rata untuk berpindah dari keadaan~$j$ ke keadaan~$i$?
\emx {Petunjuk}: Cobalah contoh Negeri Oz (Contoh~\ref{exam 11.1.1}).

\i\label{exer 11.5.15} Tunjukkan bahwa setiap rantai Markov ergodik dengan
matriks transisi simetris (yakni, $p_{ij} = p_{ji})$ bersifat reversibel.

\i\label{exer 11.5.18} (Crowell\footnote{Komunikasi
pribadi.})\index{CROWELL, R.} Misalkan $\mat{P}$ adalah matriks transisi suatu
rantai Markov ergodik. Tunjukkan bahwa
$$
(\mat {I} + \mat {P} +\cdots+ \mat {P}^{n - 1})(\mat {I} - \mat {P} + \mat {W})
= \mat {I} -
\mat {P}^n + n\mat {W}\ ,
$$
dan dari sini tunjukkan bahwa
$$
\frac{\mat {I} + \mat {P} +\cdots+ \mat {P}^{n - 1}}n \to \mat {W}\ ,
$$ 
ketika $n \rightarrow \infty$.

\i\label{exer 11.5.19} Suatu rantai Markov ergodik dimulai dalam keadaan
setimbang (yakni, dengan vektor peluang awal~$\mat{w}$). Waktu rata-rata sampai
kemunculan keadaan~$s_i$ berikutnya adalah
$\bar{m_i} = \sum_k w_k m_{ki} + w_i r_i$. Tunjukkan bahwa
$\bar {m_i} = z_{ii}/w_i$ dengan menggunakan kenyataan bahwa $\mat {w}\mat {Z}
= \mat {w}$ dan $m_{ki} = (z_{ii} - z_{ki})/w_i$.

\i\label{exer 11.5.20} Di tempat Charley berlangsung permainan
craps\index{craps} tanpa henti. Jones datang ke tempat Charley pada suatu malam
ketika 100 putaran telah dimainkan. Ia berencana bermain sampai mata ular
(sepasang angka satu) berikutnya dilempar. Jones bertanya-tanya berapa kali ia
akan bermain. Di satu pihak, ia menyadari bahwa waktu rata-rata antara dua mata
ular adalah 36, sehingga ia seharusnya bermain sekitar 18~kali karena ia sama
mungkin datang pada salah satu sisi titik tengah di antara dua kemunculan mata
ular. Di pihak lain, dadu tidak mempunyai ingatan, sehingga tampaknya ia harus
bermain 36 kali lagi, apa pun hasil-hasil sebelumnya. Manakah, jika ada, argumen
Jones yang Anda percayai? Dengan menggunakan hasil Latihan~\ref{exer 11.5.19},
hitunglah waktu yang diharapkan untuk mencapai mata ular dalam keadaan setimbang,
dan lihatlah apakah hal ini menyelesaikan paradoks yang tampak. Jika Anda masih
ragu, simulasikan percobaan tersebut untuk menentukan argumen mana yang benar.
Dapatkah Anda memberikan argumen intuitif yang menjelaskan hasil ini?

\i\label{exer 11.5.22} Tunjukkan bahwa, untuk suatu rantai Markov ergodik (lihat
Teorema \ref{thm 11.5.19}),
$$
\sum_j m_{ij} w_j = \sum_j z_{jj} - 1 = K\ .
$$
Ekspresi kedua di atas menunjukkan bahwa bilangan $K$ tidak bergantung pada $i$.
Bilangan $K$ disebut \emx {konstanta Kemeny.}\index{konstanta Kemeny} Sebuah
hadiah ditawarkan kepada orang pertama yang memberikan alasan intuitif yang
masuk akal mengapa jumlah di atas tidak bergantung pada~$i$. (Lihat juga
Latihan~\ref{exer 11.5.27}.)

\i\label{exer 11.5.23} Tinjaulah permainan yang dimainkan sebagai berikut. Anda
diberi suatu rantai Markov reguler dengan matriks transisi $\mat P$, vektor
peluang tetap $\mat{w}$, dan fungsi imbalan $\mat f$ yang menetapkan bagi setiap
keadaan $s_i$ suatu besaran $f_i$ yang dapat bernilai positif atau negatif.
Andaikan bahwa $\mat {w}\mat {f} = 0$. Anda mengamati rantai Markov ini
selama rantai itu berevolusi, dan setiap kali berada dalam keadaan $s_i$, Anda menerima
imbalan sebesar $f_i$. Tunjukkan bahwa nilai harapan kemenangan Anda setelah
$n$~langkah dapat direpresentasikan oleh vektor kolom~$\mat{g}^{(n)}$, dengan
$$\mat{g}^{(n)} = (\mat {I} + \mat {P} + \mat {P}^2 +\cdots+ \mat {P}^n)
\mat {f}.$$ Tunjukkan bahwa ketika $n \to
\infty$, $\mat {g}^{(n)} \to \mat {g}$ dengan $\mat {g} = \mat {Z} \mat {f}$.
 
\i\label{exer 11.5.24} Permainan ``Monopoly" yang sangat disederhanakan dimainkan
pada papan dengan empat petak seperti ditunjukkan pada
Gambar~\ref{fig 11.7}.\index{Monopoly} Anda mulai di GO. Anda melempar sebuah
dadu dan bergerak searah jarum jam mengelilingi papan sebanyak jumlah petak yang
sama dengan angka yang muncul pada dadu. Anda menerima atau membayar jumlah yang
tertera pada petak tempat Anda berhenti. Kemudian Anda melempar dadu lagi dan
bergerak mengelilingi papan dengan cara yang sama dari posisi terakhir. Dengan
menggunakan hasil Latihan~\ref{exer 11.5.23}, taksirlah jumlah kemenangan yang
seharusnya Anda harapkan dalam jangka panjang ketika memainkan versi Monopoly
ini.
\putfig{1truein}{PSfig11-7}{Monopoly yang disederhanakan.}{fig 11.7}

%\begin{figure}[here]
%\centerline{\epsfxsize=1truein 
%\epsffile{PSfig11-7}}
%\caption{Monopoly yang disederhanakan.}\label{fig 11.7}
%\end{figure}
  
\i\label{exer 11.5.28} Tunjukkan bahwa jika $\mat P$ adalah matriks transisi suatu
rantai Markov reguler dan $\mat W$ adalah matriks yang setiap barisnya merupakan
vektor peluang tetap yang bersesuaian dengan $\mat {P}$, maka
$\mat {P}\mat {W} = \mat {W}$ dan $\mat {W}^k = \mat {W}$ bagi semua bilangan
bulat positif $k$.

\i\label{exer 11.5.29} Andaikan suatu rantai Markov ergodik mempunyai keadaan
$s_1, s_2, \ldots, s_k$. Misalkan $S^{(n)}_j$ menyatakan berapa kali rantai
berada dalam keadaan $s_j$ selama $n$ langkah pertama. Misalkan $\mat{w}$
menyatakan vektor baris peluang tetap bagi rantai ini. Tunjukkan bahwa, terlepas
dari keadaan awal, nilai harapan $S^{(n)}_j$, dibagi dengan $n$, menuju $w_j$
ketika $n \rightarrow \infty$. \emx {Petunjuk}: Jika rantai dimulai dalam
keadaan $s_i$, nilai harapan $S^{(n)}_j$ diberikan oleh ekspresi
$$\sum_{h = 0}^n p^{(h)}_{ij}\ .$$

\i\label{exer 11.5.27} Ketika berjalan kaki bersama Snell di sepanjang Minnehaha Avenue di Minneapolis
pada musim gugur 1983, Peter Doyle\footnote{Komunikasi
pribadi.}\index{DOYLE, P. G.} mengusulkan penjelasan berikut bagi kekonstanan
\emx {konstanta Kemeny
}\index{konstanta Kemeny} (lihat
Latihan~\ref{exer 11.5.22}). Pilihlah keadaan sasaran menurut
vektor tetap~$\mat{w}$. Mulailah dari keadaan $i$ dan tunggulah sampai waktu $T$ ketika
keadaan sasaran muncul untuk pertama kalinya. Misalkan $K_i$ adalah nilai harapan
dari $T$. Perhatikan bahwa
$$
K_i + w_i \cdot 1/w_i= \sum_j P_{ij} K_j + 1\ ,
$$
dan oleh karena itu
$$
K_i = \sum_j P_{ij} K_j\ .
$$
Menurut prinsip maksimum, $K_i$ adalah konstanta.
Pantaskah Peter menerima hadiah tersebut?
\end{LJSItem}
